\chapter{Local fields}\label{ch:local}

The paper works at a bad prime $q\ne\ell$ with the restriction of $\rho=\rho_{E,\ell}$ to a decomposition group $G_{\Q_q}$, and every local statement it makes
rests on a small amount of local field theory: the fields $\Q_q$, $\Q_q^{\ur}$ and $\Z_\ell$; Hensel's lemma and $\ell$-th powers; the inertia group, its wild
part and its tame quotient; the tame relation $\Frob\,t\,\Frob^{-1}=t^q$; the Kummer cocycle of the Tate parameter on inertia and on Frobenius; quadratic
characters locally; Grothendieck's monodromy theorem and Weil-Deligne representations; and the inertia invariants, local $L$-factors and conductors of the
representations that occur at a bad prime, together with those of their adjoints. This chapter supplies all of it, with complete proofs except for local class
field theory, the identification of Kummer classes with Galois cohomology (proved in \Cref{ch:galcoh}), the integrality of the Swan conductor, upper
numbering, epsilon factors and a few classification statements, which are stated with references. It uses
\Cref{ch:groups} for profinite groups, \Cref{ch:rings} for $v_p$, $\abs{\cdot}_p$, finite fields and powers in $\F_q^\times$, \Cref{ch:modules} for discrete valuation rings and
the algebraic $\Z_\ell$, \Cref{ch:reps} for the linear algebra of inertia actions, \Cref{ch:galois} for infinite Galois theory and Kummer theory, and
\Cref{ch:ant} for number fields, decomposition groups and quadratic characters. Elliptic curves appear only in examples and in quoted statements from
\Cref{ch:reduction,ch:galrep}.

\section*{Core results at a glance}

Throughout, $K$ is $\Q_p$ or a finite extension, with $\mathcal O_K$, $\mathfrak m_K$, uniformiser $\varpi$, residue field $\F_q$; $\ell\ne p$; $\Frob$ is an arithmetic and
$\Phi=\Frob^{-1}$ a geometric Frobenius; $I_K\supseteq P_K$ are the inertia and wild inertia groups; $\norm w=q^{d(w)}$ on the Weil group.

\begin{enumerate}[label=\textbf{\thechapter.\arabic*}]
\item \textbf{Conventions.} $K$, $\mathcal O_K$, $\varpi_K$, $k_K=\F_q$, $v_K$, $\abs{\cdot}_K=q^{-v_K}$; $q$ is both $\#k_K$ and the paper's bad prime; $K^{\ur}\subseteq K^{\tame}$, $I_K$,
$P_K$; $\Z_\ell(1)=\varprojlim\mu_{\ell^k}$ and $t_\zeta$; $\varepsilon$ unramified of order $\le2$, $\theta_m$, $\eta_{\ur}$, $[a]$, $N$, $\chi_d$ by discriminant.
\Cref{conv:local:conventions}.
\item \textbf{Absolute values} (definition): nonarchimedean, equivalent, discrete; valuation ring, maximal ideal, residue field. \Cref{def:local:absval}.
\item \textbf{Basic properties.} Nonarchimedean iff $\abs{n}\le1$ on $\Z$; isosceles triangles; $\mathcal O$ local with $\mathcal O^\times=\{\abs x=1\}$; balls are open and
closed; $\sum a_n$ converges iff $a_n\to0$ (complete case); discrete implies $\abs x=c^{v(x)}$ with $\mathcal O$ a discrete valuation ring; equivalence iff
$\abs{\cdot}_1=\abs{\cdot}_2^s$. \emph{(proved here)} \Cref{prop:local:absval}.
\item \textbf{Unconditional summation.} In a complete nonarchimedean field a countable family with only finitely many terms above each $\delta>0$ can be summed in any
order and grouping. \emph{(proved here)} \Cref{lem:local:summation}.
\item \textbf{Completion.} Exists, unique up to unique isometry; the value group and the rings $\mathcal O/\varpi^n$ do not change; $\hat{\mathcal O}\simeq\varprojlim\mathcal O/\varpi^n$
topologically. \emph{(proved here)} \Cref{thm:local:completion}.
\item \textbf{$\Q_p$ and $\Z_p$} (definition): completion of $\Q$ for $\abs{\cdot}_p$, and its valuation ring. \Cref{def:local:qp}.
\item \textbf{Structure of $\Z_p$.} $\Z_p$ is the closure of $\Z$, a complete discrete valuation ring with residue field $\F_p$; $\Z_p\simeq\varprojlim\Z/p^n$ as
topological rings (so it is the algebraic $\Z_\ell$ of \Cref{ch:modules}); $\Q_p^\times=p^{\Z}\times\Z_p^\times$; $\Z_p$ compact, $\Q_p$ locally compact and totally
disconnected; $q^{-1}\in\Z_p$ for primes $q\ne p$. \emph{(proved here)} \Cref{thm:local:qp}.
\item \textbf{$p$-adic expansions.} $x=\sum_{i\ge v}a_ip^i$ uniquely, $a_i\in\{0,\dots,p-1\}$; partial sums are the residues; $x\in\Q$ iff the digits are eventually
periodic. \emph{(proved here)} \Cref{prop:local:expansions}.
\item \textbf{Hensel's lemma.} If $\abs{f(a_0)}<\abs{f'(a_0)}^2$, Newton's iteration converges to the unique root $a$ with $\abs{a-a_0}<\abs{f'(a_0)}$, and
$\abs{a-a_0}=\abs{f(a_0)/f'(a_0)}$; a simple root modulo $\mathfrak m$ lifts uniquely. \emph{(proved here)} \Cref{thm:local:hensel}.
\item \textbf{Hensel, factorisation form.} Over a complete discrete valuation ring, $\bar f=\bar g\bar h$ with $\bar g$ monic and $\bar g,\bar h$ coprime lifts to $f=gh$, $g$
monic of degree $\deg\bar g$. \emph{(proved here)} \Cref{thm:local:hensel-factor}.
\item \textbf{Coefficients of irreducible polynomials.} If $f=\sum a_iX^i\in K[X]$ is irreducible with $a_0a_n\ne0$, $\max\abs{a_i}=\max(\abs{a_0},\abs{a_n})$; a monic irreducible
polynomial with $f(0)\in\mathcal O$ lies in $\mathcal O[X]$. \emph{(proved here)} \Cref{cor:local:irreducible}.
\item \textbf{Teichm\"uller lifts.} $[a]=\lim\tilde a^{q^n}$ is the unique lift of $a\in\F_q$ with $[a]^q=[a]$; $k^\times\simeq\mu_{q-1}(K)$; $\mathcal O^\times=\mu_{q-1}(K)\times(1+\mathfrak m)$.
\emph{(proved here)} \Cref{thm:local:teichmuller}.
\item \textbf{Units.} For $p$ odd, $\exp$ and $\log$ give $p^n\Z_p\simeq1+p^n\Z_p$ ($n\ge1$), $1+p\Z_p\simeq\Z_p$ generated by $1+p$, and $\Q_p^\times\simeq\Z\times\Z/(p-1)\times\Z_p$;
for $p=2$, $\Z_2^\times=\{\pm1\}\times(1+4\Z_2)$, $1+4\Z_2\simeq\Z_2$ generated by $5$; in general $K^\times=\varpi^{\Z}\times\mu_{q-1}(K)\times U^{(1)}$ with $U^{(1)}$ pro-$p$ and uniquely
$m$-divisible for $p\nmid m$. \emph{(proved here)} \Cref{thm:local:units-structure}.
\item \textbf{Roots of unity.} $\mu(\Q_p)=\mu_{p-1}$ ($p$ odd), $\mu(\Q_2)=\{\pm1\}$; $\mu_m(K)=\mu_{\gcd(m,q-1)}$ for $p\nmid m$; $\mu_{\ell^k}\subset\Q_p$ iff $\ell^k\mid p-1$; $i\in\Q_\ell$ iff
$\ell\equiv1\pmod4$. \emph{(proved here)} \Cref{cor:local:roots-of-unity}.
\item \textbf{$\ell^k$-th powers ($\ell\ne p$).} A unit is an $\ell^k$-th power iff its reduction is one in $\F_q^\times$; $\varpi^au$ is one iff moreover $\ell^k\mid a$;
$K^\times/(K^\times)^{\ell^k}\simeq\Z/\ell^k\times\Z/\ell^{\min(k,v_\ell(q-1))}$, of order $\ell^k\#\mu_{\ell^k}(K)$; $\varprojlim K^\times/(K^\times)^{\ell^k}\simeq\Z_\ell\times\Z/\ell^{v_\ell(q-1)}$ with
torsion $\F_q^\times\otimes\Z_\ell$; square classes of $\Q_p$ and $\Q_2$. \emph{(proved here)} \Cref{thm:local:powers}.
\item \textbf{Equivalence of norms.} All norms on a finite-dimensional space over a complete field are equivalent, and the space is complete. \emph{(proved here)}
\Cref{lem:local:norms}.
\item \textbf{Extension of the absolute value.} $\abs x_L=\abs{N_{L/K}x}^{1/[L:K]}$ is the unique extension; $L$ is complete; $\mathcal O_L$ is the integral closure; conjugates
have equal absolute values. \emph{(proved here)} \Cref{thm:local:extension-absval}.
\item \textbf{$e$ and $f$.} $\mathcal O_L$ is free over $\mathcal O_K$ with basis $\alpha_i\varpi_L^j$, and $ef=[L:K]$; multiplicative in towers. \emph{(proved here)} \Cref{thm:local:ef}.
\item \textbf{Unramified, totally, tamely, wildly ramified} (definition): $e=1$, $f=1$, $p\nmid e$, $p\mid e$. \Cref{def:local:ram-types}.
\item \textbf{Eisenstein.} $L/K$ is totally ramified of degree $e$ iff it is generated by a root $\pi$ of an Eisenstein polynomial of degree $e$; then $\pi$ is a uniformiser
and $\mathcal O_L=\mathcal O_K[\pi]$. \emph{(proved here)} \Cref{thm:local:eisenstein}.
\item \textbf{Examples; $\Z_3[\zeta_3]$.} $\Q_p(\zeta_{p^n})$ is totally ramified with group $(\Z/p^n)^\times$; $\Q_3(\zeta_3)=\Q_3(\sqrt{-3})$ is tamely totally ramified with ring of
integers $\Z_3[\zeta_3]=\Z_3\otimes\Z[\zeta_3]$, uniformiser $1-\zeta_3$, $(1-\zeta_3)^2=-3\zeta_3$, the completion of $\Z[\zeta_3]$; $\Q_p(p^{1/m})$, $\Q_p(\sqrt{p^*})$.
\emph{(proved here)} \Cref{prop:local:z3zeta3}.
\item \textbf{Completions of number fields.} The extensions of $\abs{\cdot}_p$ to $F$ correspond to the $\mathfrak p\mid p$; $F_{\mathfrak p}$ has $e_{\mathfrak p}$, $f_{\mathfrak p}$ and
$\mathcal O_{F_{\mathfrak p}}=\varprojlim\mathcal O_F/\mathfrak p^n$; $\mathcal O_F\otimes\Z_p\simeq\prod_{\mathfrak p\mid p}\mathcal O_{F_{\mathfrak p}}$. \emph{(proved here)}
\Cref{thm:local:completion-numberfield}.
\item \textbf{Krasner's lemma.} If $\abs{\beta-\alpha}<\abs{\alpha-\alpha_i}$ for the other conjugates, $K(\alpha)\subseteq K(\beta)$; nearby polynomials define the same fields.
\emph{(proved here)} \Cref{lem:local:krasner}.
\item \textbf{Decomposition groups.} $\Gal(L_{\mathfrak P}/F_{\mathfrak p})\simeq D_{\mathfrak P}$ with inertia to inertia; an embedding $\Qbar\to\Qbar_p$ identifies $G_{\Q_p}$ with
$D_{\bar{\mathfrak P}}\subseteq G_\Q$, $I_p$ with $I_{\bar{\mathfrak P}}$, Frobenius with Frobenius. \emph{(proved here)} \Cref{prop:local:decomposition}.
\item \textbf{Tame totally ramified extensions.} $L=K(\varpi^{1/e})$; Galois iff $e\mid q-1$, and then cyclic, $\sigma\mapsto\sigma(\alpha)/\alpha$. \emph{(proved here)}
\Cref{prop:local:tame-totally-ramified}.
\item \textbf{Unramified extensions.} Exactly one of each degree, $K_f=K(\mu_{q^f-1})$; $\Gal(K_f/K)\simeq\Gal(\F_{q^f}/\F_q)$, cyclic, generated by $\Frob$; largest unramified
subextension $L_0$ with $L/L_0$ totally ramified. \emph{(proved here)} \Cref{thm:local:unramified}.
\item \textbf{$K^{\ur}$} (definition): $\bigcup K_f=K(\mu_m:p\nmid m)$. \Cref{def:local:qpur}.
\item \textbf{Properties of $K^{\ur}$.} $\Gal(K^{\ur}/K)\simeq\hat\Z$, generated by $\Frob$; discrete, not complete; Hensel holds; units are $m$-th powers for $p\nmid m$, and
$x$ is an $m$-th power iff $m\mid v_K(x)$; $\mu_{\ell^\infty}\subset\Q_q^{\ur}$. \emph{(proved here)} \Cref{prop:local:qpur}.
\item \textbf{Frobenius convention.} $\Frob$ arithmetic, $\Phi=\Frob^{-1}$ geometric; $\chi_\ell(\Frob)=q$; group theory uses $\Frob$, $L$-factors use $\Phi$. \Cref{conv:local:frobenius}.
\item \textbf{Unramified base change.} For $K/\Q_p$ unramified: $K^{\ur}=\Q_p^{\ur}$, $I_K=I_p$, $P_K=P_p$; images of inertia, invariants, monodromy and conductors are unchanged. In
general $I_K=I_p\cap G_K$ has index $e(K/\Q_p)$. \emph{(proved here)} \Cref{prop:local:unram-basechange}.
\item \textbf{The cyclotomic character locally.} $\chi_\ell$ is unramified at $p\ne\ell$ with $\chi_\ell(\Frob_p)=p$; $\chi_p(I_p)=\Z_p^\times$. \emph{(proved here)}
\Cref{prop:local:cyclotomic-local}.
\item \textbf{Ramification groups} (definition): $G_i=\{\sigma:v_L(\sigma x-x)\ge i+1\ \forall x\in\mathcal O_L\}$; $G_0$ inertia, $G_1$ wild inertia. \Cref{def:local:ramgroups}.
\item \textbf{Structure of the $G_i$.} $\#G_0=e$; $G_0/G_1\into k_L^\times$ by $\sigma\mapsto\sigma\pi/\pi$; $G_i/G_{i+1}\into(k_L,+)$ for $i\ge1$; so $G_1$ is the Sylow $p$-subgroup of $G_0$,
and $L/K$ is tame iff $G_1=1$. Upper numbering, Herbrand and Hasse-Arf are stated. \emph{(proved here, except (iv): stated; reference)}
\Cref{prop:local:ramgroups}.
\item \textbf{The filtration $G_K\supset I_K\supset P_K$.} $G_K/I_K\simeq\hat\Z$; $K^{\tame}=\bigcup_{p\nmid m}K^{\ur}(\varpi^{1/m})$; $P_K=\Gal(\Qbar_p/K^{\tame})$ is normal, pro-$p$, maps onto every
$G_1(L/K)$, and is the unique pro-$p$ Sylow subgroup of $I_K$. \emph{(proved here, except the upper filtration (iv): stated; reference)}
\Cref{thm:local:inertia-structure}.
\item \textbf{Semidirect product.} $G_K=I_K\rtimes\overline{\langle\Frob\rangle}$ with $\overline{\langle\Frob\rangle}\simeq\hat\Z$; a Frobenius-invariant character of $I_K$ extends to $G_K$.
\emph{(proved here)} \Cref{cor:local:split}.
\item \textbf{The tame character.} $\theta_m(\sigma)=\sigma(\varpi^{1/m})/\varpi^{1/m}$ is a surjection $I_K\to\mu_m$ independent of choices; $I_K/P_K\simeq\prod_{\ell\ne p}\Z_\ell(1)$, procyclic of
order prime to $p$; $t_\ell\colon I_K\to\Z_\ell(1)$ identifies the maximal pro-$\ell$ quotient of $I_K$ with $\Z_\ell(1)$; maps from $I_K$ to pro-$\ell$ groups factor through
$t_\ell$; $t_{\ell,K}|_{I_L}=e(L/K)t_{\ell,L}$. \emph{(proved here)} \Cref{thm:local:tame-character}.
\item \textbf{Open subgroups of inertia.} Every open $U\subseteq I_K$ is $I_L$ for a finite $L/K$; its maximal pro-$\ell$ quotient is $\Z_\ell(1)$. \emph{(proved here)}
\Cref{cor:local:open-inertia}.
\item \textbf{A Frobenius fixing $\varpi^{1/m}$.} Some arithmetic Frobenius lift fixes all $\varpi^{1/m}$, $p\nmid m$; two differ by $P_K$. \emph{(proved here)}
\Cref{cor:local:frob-fixing-roots}.
\item \textbf{The tame relation.} $t_\ell(\tau\sigma\tau^{-1})=\chi_\ell(\tau)t_\ell(\sigma)$. On $I_K/P_K$ one has $\Frob\,t\,\Frob^{-1}=t^q$, equivalently
$\Frob^{-1}t\,\Frob=t^{q^{-1}}$, and for the geometric Frobenius $\Phi\,t\,\Phi^{-1}=t^{q^{-1}}$. The group $\Gal(K^{\tame}/K)$ is
$\prod_{\ell\ne p}\Z_\ell(1)\rtimes\hat\Z$, with $1$ acting by $q$. \emph{(proved here; the presentation by two generators is stated; reference)}
\Cref{thm:local:tame-relation}.
\item \textbf{Finite tame images of inertia.} If $\rho(I_K)$ is finite of order prime to $p$, it is cyclic and $\rho(\Frob)h\rho(\Frob)^{-1}=h^q$; Frobenius centralises
(inverts) a generator of order $m$ iff $q\equiv1$ ($q\equiv-1$) $\pmod m$. \emph{(proved here)} \Cref{cor:local:finite-tame-image}.
\item \textbf{Wild inertia in $\ell$-adic representations.} $\rho(P_K)$ is a finite $p$-group embedding in $\GL_n(\mathcal O/\varpi_\lambda)$; on an open normal $I_1\subseteq I_K$, $\rho$
factors through $t_\ell$; $\rho$ is tame iff $\rho(I_K)$ has no element of order $p$. \emph{(proved here)} \Cref{prop:local:wild-image}.
\item \textbf{Kummer theory over $K$.} The Kummer cocycle $\kappa_x$ and its properties; $K^\times/(K^\times)^{\ell^k}\simeq H^1(K,\mu_{\ell^k})$ and $\widehat{K^\times}\simeq H^1(K,\Z_\ell(1))\simeq\Z_\ell\times
\Z/\ell^{v_\ell(q-1)}$. \emph{(proved here, except the isomorphisms with $H^1$: stated; \Cref{ch:galcoh})} \Cref{prop:local:kummer-qp}.
\item \textbf{The Kummer cocycle on inertia and Frobenius.} $\kappa_x=v_K(x)\kappa_\varpi+\kappa_u$; $\kappa_x|_{I_K}=v_K(x)\,t_\ell$; $\kappa_u$ is unramified with class the class of $\bar u$
in $\F_q^\times\otimes\Z_\ell$; $c'=\kappa_x(\Frob)$ is defined modulo $(v_K(x),q-1)$; $I_K$ fixes $x^{1/\ell^k}$ iff $\ell^k\mid v_K(x)$. \emph{(proved here)}
\Cref{thm:local:kummer-inertia}.
\item \textbf{Quadratic extensions of $\Q_p$.} Three for $p$ odd (one unramified, two tame), restricting to the unique order-two character $\theta_2$ of $I_p$ when
ramified; seven for $p=2$, with $\Q_2(\sqrt5)=\Q_2(\sqrt{-3})$ unramified and jumps $1$ for $\Q_2(\sqrt{-1})$, $\Q_2(\sqrt3)$, $2$ for $\Q_2(\sqrt{\pm2})$, $\Q_2(\sqrt{\pm10})$.
\emph{(proved here)} \Cref{prop:local:quadratic-ext}.
\item \textbf{Global quadratic characters locally.} $\chi_d$ is unramified at $p\nmid d$ with $\chi_d(\Frob_p)=(d/p)$; $\chi_{q^*}$ is unramified outside $q$ and equals $\theta_2$ on $I_q$;
twisting by $\chi_{q^*}$ changes nothing on inertia at $p\ne q$; every quadratic character of $G_{\Q_p}$ is a restriction of a product of $\chi_d$'s. \emph{(proved
here)} \Cref{prop:local:global-quadratic-local}.
\item \textbf{Local class field theory for $\Q_p$.} $\Q_p^{\mathrm{ab}}=\Q_p^{\ur}\Q_p(\mu_{p^\infty})$; the Artin map ($p\mapsto\Frob$, $u\mapsto(\zeta\mapsto\zeta^{u^{-1}})$); $a(\chi)$ is the least $n$
with $\chi(\Art(U^{(n)}))=1$. \emph{(stated; reference)} \Cref{thm:local:lcft}.
\item \textbf{Weil-Deligne relation.} $\norm w=q^{d(w)}$, $r(w)Nr(w)^{-1}=\norm wN$: $r(\Frob)Nr(\Frob)^{-1}=qN$, $r(\Phi)Nr(\Phi)^{-1}=q^{-1}N$. \Cref{conv:local:weil}.
\item \textbf{The Weil group} (definition): $W_K=\{w\in G_K:w|_{K^{\ur}}\in\Frob^{\Z}\}$, with $I_K$ open. \Cref{def:local:weil-group}.
\item \textbf{Properties of $W_K$.} Dense in $G_K$, $W_K/I_K\simeq\Z$, $\norm w=\chi_\ell(w)$; unramified characters are free on $\Phi$; $W_{\Q_p}^{\mathrm{ab}}\simeq\Q_p^\times$ (stated).
\emph{(proved here, except (iv): stated; reference)} \Cref{prop:local:weil-group}.
\item \textbf{Matrix logarithm.} $\exp$ and $\log$ are inverse bijections $\ell^2\Mat_n(\mathcal O)\leftrightarrows1+\ell^2\Mat_n(\mathcal O)$, with $g^a=\exp(a\log g)$. \emph{(proved here)}
\Cref{lem:local:matrix-log}.
\item \textbf{Grothendieck's monodromy theorem.} For $\rho\colon G_K\to\GL(V)$ $\ell$-adic, there is a unique nilpotent $N$ with $\rho(\sigma)=\exp(t_\zeta(\sigma)N)$ on an open normal $I_1\subseteq I_K$,
and $\rho(w)N\rho(w)^{-1}=\norm wN$; $N=0$ iff $\rho(I_K)$ is finite. \emph{(proved here)} \Cref{thm:local:grothendieck}.
\item \textbf{Weil-Deligne representations} (definition): pairs $(r,N)$, with $r$ a representation of $W_K$ such that $r(I_K)$ is finite and $N$ nilpotent,
satisfying $r(w)Nr(w)^{-1}=\norm wN$; Frobenius-semisimplicity. \Cref{def:local:wd}.
\item \textbf{Operations.} Sums, tensor products, duals ($-N^{\mathsf T}$), $\Hom$, $\ad$ ($[N,\cdot\,]$), twists. \emph{(proved here; Frobenius-semisimplification stated)}
\Cref{prop:local:wd-constructions}.
\item \textbf{$\WD(\rho)$.} $r(\Phi^n\sigma)=\rho(\Phi^n\sigma)\exp(-t_\zeta(\sigma)N)$ is a Weil-Deligne representation, independent of $\Phi$ and $\zeta$ up to isomorphism; equivalence of
categories and independence of $\ell$ for elliptic curves are stated. \emph{(proved here, except (iii), (iv): stated; reference)} \Cref{thm:local:wd-equivalence}.
\item \textbf{$\Sp(n)$} (definition): $r(w)e_i=\norm w^ie_i$, $Ne_i=e_{i+1}$; twists $\Sp(n)\otimes\mu$; unramified and ramified twists of Steinberg. \Cref{def:local:sp}.
\item \textbf{Two-dimensional Weil-Deligne representations.} $\Sp(n)\otimes\mu$ is indecomposable; an $(r,N)$ of dimension $2$ with $N\ne0$ is Frobenius-semisimple and
isomorphic to $\Sp(2)\otimes\mu$ for a unique $\mu$; the case $N=0$ is stated. \emph{(proved here, except (iii): stated; reference)} \Cref{prop:local:sp}.
\item \textbf{Inertia invariants.} $V^{\rho(I_K)}=(\ker N)^{r(I_K)}$, with $\rho(\Phi)=r(\Phi)$ there. \emph{(proved here)} \Cref{prop:local:invariants}.
\item \textbf{Local $L$-factor} (definition): $L(V,s)=\det(1-\Phi q^{-s}\mid V^{I_K})^{-1}$, $\Phi$ geometric. \Cref{def:local:L-factor}.
\item \textbf{Properties of $L$-factors.} Well defined, additive; $L(\Q_\ell(1),s)=(1-q^{-1-s})^{-1}$, $L(V(1),s)=L(V,s+1)$ (false with $\Frob$); $L(\Sp(2)\otimes\mu,s)$;
$L_p(V_\ell E,s)=L_p(E,s+1)$ and $L_p(V_\ell E^\vee,s)=L_p(E,s)$. \emph{(proved here)} \Cref{prop:local:L-factor}.
\item \textbf{Swan and Artin conductors} (definition): $\operatorname{sw}(V)=\sum_{i\ge1}\frac{g_i}{g_0}\dim V/V^{G_i}$, $a(V)=\dim V-\dim V^{I_K}+\operatorname{sw}(V)$. \Cref{def:local:swan}.
\item \textbf{Properties of the conductors.} Independence of $L$ and integrality (stated); $\operatorname{sw}=0$ iff tame; additivity, invariance under unramified twists and base change;
$a(\chi)=1$ for tame $\chi$, $a(\chi_{-4})=2$, $a(\chi_{\pm8})=3$ at $2$; the same for mod $\ell$ representations. \emph{(proved here, except (i), (ii): stated; reference)}
\Cref{thm:local:swan}.
\item \textbf{Conductor of a Weil-Deligne representation} (definition): the exponent is
$a(r,N)=\dim V-\dim(\ker N)^{r(I_K)}+\operatorname{sw}(r)$. \Cref{def:local:conductor}.
\item \textbf{Properties.} $a(\rho)=\dim V-\dim V^{I_K}+\operatorname{sw}(\rho)$; $a(V\otimes\chi)=\dim V\cdot a(\chi)$ for $V$ unramified; $a(\Sp(2)\otimes\mu)=1$, $a(\Sp(2)\otimes\mu\chi)=2a(\chi)$; the conductor
exponent of $E$ at $p$ is $a(V_\ell E)$ (definition, \Cref{ch:reduction}), equal to $v_p(\Delta_{\min})+1-m$ with $m$ the number of components of the
special fibre of the minimal regular model (Ogg, Saito). \emph{(proved here; independence of $\ell$ and Ogg's formula stated; reference)} \Cref{prop:local:conductor}.
\item \textbf{Epsilon factors.} Existence and uniqueness (Langlands, Deligne) and the Weil-Deligne correction factor. \emph{(stated; reference)} \Cref{thm:local:epsilon}.
\item \textbf{Unramified shape.} $A$ has eigenvalues $\alpha/\beta,1,\beta/\alpha$ for both Frobenii; $L(A,1)^{-1}=\frac{q-1}q\cdot\frac{(q+1)^2-a^2}{q^2}$. \emph{(proved here)}
\Cref{prop:local:unramified-shape}.
\item \textbf{Finite cyclic tame inertia.} Order $2$: $V\otimes\chi$ unramified; orders $3,4,6$: $V^{I_q}=0$, $a(V)=2$, $A^{I_q}=\Omega(2g-\tr g)$ with $\Frob$ and $\Phi$ acting by $\chi(q)$,
$\chi=\chi_{-3}$ or $\chi_{-4}$, $L(A,s)^{-1}=1-\chi(q)q^{-s}$, $a(A)=2$. \emph{(proved here)} \Cref{prop:local:tame-cyclic-shape}.
\item \textbf{Steinberg and its twists.} $a(\Sp(2)\otimes\mu)=1$, $a(\Sp(2)\otimes\mu\chi)=2a(\chi)$; $\ad\Sp(2)\simeq\Sp(3)\otimes\norm{\cdot}^{-1}$, $A^{I_q}=\Omega N$ with $\Phi$-eigenvalue $q^{-1}$,
$L(\ad V,s)=(1-q^{-1-s})^{-1}$ for both. \emph{(proved here)} \Cref{prop:local:steinberg-shape}.
\item \textbf{The table of shapes at a tame prime.} $\rho(I_q)$, $\dim V^{I_q}$, $a(V)$, $L(V,s)$, $\dim A^{I_q}$, $L(A,s)$ for the six shapes, for $q\ge5$ (and partly for
$q=2,3$). \emph{(proved here; the Kodaira column
stated, \Cref{ch:reduction,ch:galrep})} \Cref{thm:local:shapes-table}.
\item \textbf{The primes $2$ and $3$.} An inertia image of order divisible by $3$ is wild at $3$ ($a\ge3$); order $3$ is tame at $2$ with $L(A,s)^{-1}=1+2^{-s}$; non-abelian finite
images have $A^{I_q}=0$; $a(\Sp(2)\otimes\mu\chi)\in\{4,6\}$ at $2$. \emph{(proved here)} \Cref{prop:local:wild-shapes}.
\end{enumerate}

\section*{Where the paper uses this chapter}

\begin{itemize}
\item \textbf{Section 1.1} (the local terms): $I_p$ and $D_p\simeq G_{\Q_p}$, \Cref{prop:local:decomposition}; $\delta_p=\dim V^{I_p}$, \Cref{prop:local:invariants};
$L_p(A_f,s)=\det(1-\Frob_p^{-1}p^{-s}\mid A_f^{I_p})^{-1}$ and the naive factor, \Cref{def:local:L-factor}, \Cref{conv:local:frobenius}, \Cref{pl:local:Lfactor}.
\item \textbf{Definition def:terms and the table of Section 1}: the shapes and their factors, \Cref{thm:local:shapes-table}; $e_\ell=0$ at multiplicative primes,
\Cref{prop:local:steinberg-shape}; potentially good $\mathrm{I}_n^*$ at $2$, \Cref{rem:local:wild-paper}.
\item \textbf{Lemma lem:tamagawa-exp} (unramified at $q$ modulo $\ell^t$ iff $\ell^t\mid n_q$): unramified base change, \Cref{prop:local:unram-basechange}; $\mu_{\ell^\infty}\subset\Q_q^{\ur}$ and
$\ell^t$-th powers in $\Q_q^{\ur}$, \Cref{prop:local:qpur}(iv); Kummer theory on inertia, \Cref{thm:local:kummer-inertia}(v); \Cref{pl:local:tamagawa-exp-ur},
\Cref{pl:local:kummer}.
\item \textbf{Lemma lem:IV and Remark rem:types}: $\Z_3[\zeta_3]$ as the ring of integers of $\Q_3(\sqrt{-3})$, \Cref{prop:local:z3zeta3}; tame inertia images are
cyclic and Frobenius acts by the $q$-th power, \Cref{cor:local:finite-tame-image}; wild inertia, \Cref{prop:local:wild-image}; conductor exponents and the
Serre conductor, \Cref{thm:local:swan}(v), \Cref{prop:local:conductor}, \Cref{prop:local:tame-cyclic-shape}; the value at $q=2$, \Cref{prop:local:wild-shapes}(ii).
\item \textbf{Lemma lem:twist}: (i) unramified base change and the local triviality of $\chi_{q^*}$ at $p\ne q$, \Cref{prop:local:unram-basechange},
\Cref{prop:local:global-quadratic-local}(v); (ii) tame conductor exponents for the exchanges listed, \Cref{prop:local:quadratic-ext}, \Cref{prop:local:conductor};
(iii) the ramified twist of Steinberg has conductor exponent $2$, \Cref{prop:local:steinberg-shape}(ii).
\item \textbf{Lemma lem:euler}: Frobenius centralises or inverts the generator of inertia, \Cref{cor:local:finite-tame-image}; the invariant lines of $\ad V$ and their
eigenvalues, \Cref{prop:local:tame-cyclic-shape}, \Cref{prop:local:steinberg-shape}; the rings $\Z_\ell[\zeta_3]$, $\Z_\ell[i]$, \Cref{ex:local:hensel-examples}(g);
\Cref{pl:local:euler}.
\item \textbf{Proposition prop:shapes}: the four shapes of $P_q$, \Cref{thm:local:shapes-table}, \Cref{prop:local:wild-shapes}(iii); $G_{\Q_q}=I_q\rtimes\langle\phi\rangle$ and the
extension of $\psi$, \Cref{cor:local:split}; quadratic characters of $G_{\Q_q}$ as restrictions of Dirichlet characters, \Cref{prop:local:global-quadratic-local}(vi).
\item \textbf{Remark rem:wild and Section sec:wild}: wild inertia at $2$ and $3$, \Cref{rem:local:wild-ec}, \Cref{prop:local:wild-shapes}, \Cref{rem:local:wild-paper}.
\item \textbf{Lemma lem:tamfactor}: the tame character and the maximal pro-$\ell$ quotient of inertia, \Cref{thm:local:tame-character}; $\Frob^{-1}t\Frob=t^{q^{-1}}$,
\Cref{thm:local:tame-relation}; the Kummer cocycle of $q_E$, its restriction $n\,t_\ell$ to inertia, the class of $u$ and the indeterminacy of $c'$,
\Cref{thm:local:kummer-inertia}, \Cref{cor:local:frob-fixing-roots}; $\varprojlim\Q_q^\times/(\Q_q^\times)^{\ell^k}$ and $\Z_\ell(1)/(q-1)$, \Cref{thm:local:powers},
\Cref{prop:local:kummer-qp}; the matrices with corner $\varepsilon\kappa$, \Cref{ex:local:tate-shape-wd}; the class $c$ read through the reduction of $u$,
\Cref{pl:local:powers}; \Cref{pl:local:tamfactor-tame}.
\item \textbf{Lemma lem:tamadd}: the maximal pro-$\ell$ quotient of an open subgroup of inertia, \Cref{cor:local:open-inertia}.
\item \textbf{Lemma lem:fpr and the paragraph after Lemma lem:dfg21}: $L_p(M(1),s)=L_p(M,s+1)$ with the geometric Frobenius, \Cref{prop:local:L-factor}(iii),
\Cref{warn:local:geometric}.
\item \textbf{Lemmas lem:crit, lem:old, lem:ddt417}: the completions $\mathcal O_\lambda$ and $\mathcal O\otimes\Z_\ell=\prod\mathcal O_\lambda$, \Cref{thm:local:completion-numberfield}.
\item \textbf{Theorems thm:twist, thm:quant, thm:qual}: $\chi_S$ and $\chi_{q^*}$ locally, \Cref{prop:local:global-quadratic-local}.
\end{itemize}

\section{Absolute values, completions, \texorpdfstring{$\Q_p$}{Qp} and \texorpdfstring{$\Z_p$}{Zp}}\label{sec:local:qp}

This section builds the fields in which every local statement of the paper lives: $\Q_q$ at a bad prime $q$, and $\Q_\ell$,
$\Z_\ell$ for the coefficients. \Cref{ch:rings} defined $v_p$ and $\abs{\cdot}_p$ on $\Q$, and \Cref{ch:modules} built $\Z_\ell$
algebraically as an inverse limit; here we complete $\Q$ analytically and show that the two constructions agree
(\Cref{thm:local:qp}(ii)). The paper uses these fields without comment in Sections 2 and 4; nothing in this section is a
statement of the paper, but everything after it depends on it.

\begin{convention}[local fields]\label{conv:local:conventions}
\begin{enumerate}[label=(\alph*)]
\item $p$ is a prime and $K$ denotes $\Q_p$ or a finite extension of $\Q_p$ (\Cref{sec:local:ext}), with valuation ring $\mathcal O_K$, maximal
ideal $\mathfrak m_K$, uniformiser $\varpi_K$, residue field $k_K=\F_q$, $q=p^{f}$, normalised valuation $v_K$ ($v_K(\varpi_K)=1$) and normalised
absolute value $\abs{x}_K=q^{-v_K(x)}$. On $\Q_p$ we write $v_p$, $\abs{\cdot}_p$ as in \Cref{ch:rings}. The letter $q$ is both $\#k_K$ and,
when the paper's prime is meant, the bad prime $q$; the two agree for $K=\Q_q$, the field over which the paper works (when its unramified
quadratic extension appears, its residue field is written $\F_{q^2}$). In the general theory
the residue characteristic is $p$ and $\ell\ne p$ is another prime; at the paper's prime, $\ell\ne q$.
\item $\Qbar_p$ is a fixed algebraic closure of $\Q_p$ and $G_K=\Gal(\Qbar_p/K)$. Inside $\Qbar_p$ we have $K\subseteq K^{\ur}\subseteq K^{\tame}$,
the maximal unramified and the maximal tamely ramified extensions (\Cref{def:local:qpur}, \Cref{thm:local:inertia-structure});
$I_K=\Gal(\Qbar_p/K^{\ur})$ is the inertia group and $P_K=\Gal(\Qbar_p/K^{\tame})$ the wild inertia group; for $K=\Q_p$ we write $I_p$, $P_p$.
\item $\Z_\ell(1)=\varprojlim_k\mu_{\ell^k}$ (transition maps $\zeta\mapsto\zeta^\ell$), on which $G_K$ acts through the cyclotomic character
$\chi_\ell$; it is written additively when it receives the values of cocycles. Where a $\Z_\ell$-valued version is needed we fix a
generator $\zeta=(\zeta_{\ell^k})_k$ of $\Z_\ell(1)$ and identify $\Z_\ell(1)\simeq\Z_\ell$ by $\zeta\mapsto1$.
\item $\Frob$ is an \emph{arithmetic} Frobenius and $\Phi=\Frob^{-1}$ the \emph{geometric} one (\Cref{conv:local:frobenius}); $\Phi$ without subscript
is always the geometric Frobenius of the local field in question, and the geometric Frobenius at a named prime $p$ is written $\Frob_p^{-1}$. With a
numerical subscript, as in $\Phi_{p^n}$ and $\Phi_3$, the letter denotes a cyclotomic polynomial (\Cref{ch:rings}); $\Phig_q$ is the geometric component
group of the N\'eron model at $q$ (\Cref{ch:reduction}).
\item Other letters: $\varepsilon$ is an unramified character of order at most $2$ (the paper's Tate data), $\epsilon$ an
epsilon factor; $\theta_m$ are the tame characters of \Cref{thm:local:tame-character}, $\theta_2$ the character of order two of inertia, and
$\eta_{\ur}$ the unramified quadratic character; $[a]$ is the Teichm\"uller lift of $a$; $N$ is a monodromy operator (Sections
\ref{sec:local:wd} to \ref{sec:local:shapes}), never a conductor; $e$ is a ramification index, except in \Cref{sec:local:shapes},
where it is also the paper's basis vector $e=e_1\otimes e_2^\vee$ of $\ad T$, and we say which. Quadratic characters are $\chi_d$, indexed
by the fundamental discriminant $d$ ($\chi_{-3}$, $\chi_{-4}$, $\chi_8$, $\chi_{q^*}$); a letter subscript $\chi_\ell$ is the
cyclotomic character.
\end{enumerate}
\end{convention}

\subsection{Absolute values}

\begin{definition}[absolute values]\label{def:local:absval}
An \emph{absolute value} on a field $F$ is a map $\abs{\cdot}\colon F\to\R_{\ge0}$ with $\abs x=0$ iff $x=0$, $\abs{xy}=\abs x\abs y$ and
$\abs{x+y}\le\abs x+\abs y$. It is \emph{nonarchimedean} if $\abs{x+y}\le\max(\abs x,\abs y)$ for all $x,y$ (the ultrametric
inequality), and \emph{trivial} if $\abs x=1$ for all $x\ne0$. Two absolute values are \emph{equivalent} if they define the same
topology. For a nonarchimedean $\abs{\cdot}$ the \emph{valuation ring} is $\mathcal O=\{\abs x\le1\}$, its \emph{maximal ideal}
$\mathfrak m=\{\abs x<1\}$, and $k=\mathcal O/\mathfrak m$ the \emph{residue field}. The absolute value is \emph{discrete} if
$\abs{F^\times}$ is a discrete subgroup of $\R_{>0}$.
\end{definition}

The definition abstracts the one example that matters here. On $\Q$, $\abs x_p=p^{-v_p(x)}$ (\Cref{ch:rings}) measures how divisible $x$ is
by $p$: $\abs{12}_2=\frac14$, $\abs{1/12}_2=4$, $\abs{12}_3=\frac13$ and $\abs{12}_5=1$. It is nonarchimedean, because a sum of two multiples of $p^k$ is
again a multiple of $p^k$, that is $v_p(x+y)\ge\min(v_p(x),v_p(y))$; and the inequality can be strict, as $\abs{2+6}_2=\frac18<\frac12$ shows. Its
value group is $p^{\Z}$, so it is discrete; its valuation ring is $\Z_{(p)}$, the fractions $a/b$ with $p\nmid b$, with maximal ideal $p\Z_{(p)}$
and residue field $\F_p$. Two numbers are $p$-adically close when their difference is divisible by a high power of $p$; completing $\Q$ for
this notion of closeness is what produces $\Q_p$ below.

\begin{proposition}\label{prop:local:absval}
Let $\abs{\cdot}$ be an absolute value on a field $F$.
\begin{enumerate}[label=(\roman*)]
\item $\abs{\cdot}$ is nonarchimedean iff $\abs{n\cdot1}\le1$ for all $n\in\Z$.
\item If $\abs{\cdot}$ is nonarchimedean: $\abs{x+y}=\max(\abs x,\abs y)$ whenever $\abs x\ne\abs y$; $\mathcal O$ is a subring of $F$,
$\mathfrak m$ its unique maximal ideal and $\mathcal O^\times=\{\abs x=1\}$; every point of an open ball $B(a,r)=\{\abs{x-a}<r\}$ is a
centre of it, and open balls are closed. If $F$ is complete, a series $\sum a_n$ converges iff $a_n\to0$.
\item If $\abs{\cdot}$ is nonarchimedean, discrete and nontrivial, then $\abs{F^\times}=c^{\Z}$ for a unique $c\in(0,1)$, there is a
valuation $v\colon F^\times\to\Z$ with $\abs x=c^{v(x)}$, and $\mathcal O$ is a discrete valuation ring with maximal ideal $\varpi\mathcal O$
for any $\varpi$ with $v(\varpi)=1$.
\item Two nontrivial absolute values $\abs{\cdot}_1$, $\abs{\cdot}_2$ are equivalent iff $\abs{\cdot}_1=\abs{\cdot}_2^{s}$ for some $s>0$.
\end{enumerate}
\end{proposition}

\begin{proof}
(i) If the ultrametric inequality holds, $\abs{n\cdot1}=\abs{1+\dots+1}\le1$ by induction. Conversely, if $\abs{n\cdot1}\le1$ for all $n$, then
for every $k\ge1$
\[
\abs{x+y}^k=\Bigl\lvert\sum_{i=0}^k\tbinom ki x^iy^{k-i}\Bigr\rvert\le(k+1)\max(\abs x,\abs y)^k,
\]
since each binomial coefficient is an integer, hence of absolute value at most $1$, and each $\abs{x^iy^{k-i}}$ is at most $\max(\abs x,\abs y)^k$.
Taking $k$-th roots gives $\abs{x+y}\le(k+1)^{1/k}\max(\abs x,\abs y)$, and $(k+1)^{1/k}\to1$ as $k\to\infty$, so $\abs{x+y}\le\max(\abs x,\abs y)$.

(ii) If $\abs x<\abs y$ then $\abs y=\abs{(x+y)-x}\le\max(\abs{x+y},\abs x)$ forces $\abs y\le\abs{x+y}$, and $\abs{x+y}\le\abs y$ always. The
ultrametric inequality and multiplicativity show that $\mathcal O$ is a ring and $\mathfrak m$ an ideal; an element of $\mathcal O\setminus\mathfrak m$ has
absolute value $1$, so its inverse lies in $\mathcal O$; hence $\mathfrak m$ is the set of nonunits, the unique maximal ideal. If
$\abs{b-a}<r$ and $\abs{x-a}<r$ then $\abs{x-b}<r$, so $B(a,r)=B(b,r)$. If $c\notin B(a,r)$, then $B(c,r)\cap B(a,r)=\emptyset$ (a common point
would make both balls equal), so the complement of $B(a,r)$ is open. Finally $\abs{\sum_{i=n+1}^ma_i}\le\max_{n<i\le m}\abs{a_i}$, so the partial
sums are Cauchy iff $a_i\to0$; in a complete field Cauchy sequences converge, and conversely if the partial sums $s_n$ converge then
$a_n=s_n-s_{n-1}\to0$. (Over $\R$ the harmonic series shows that $a_n\to0$ is not enough; the ultrametric inequality is what makes it so.)

(iii) A nontrivial discrete subgroup of $\R_{>0}$ (isomorphic by $\log$ to a discrete subgroup of $\R$) is infinite cyclic; let $c<1$ be
its generator below $1$ and $v(x)$ the integer with $\abs x=c^{v(x)}$. Then $v$ is a valuation: $v(xy)=v(x)+v(y)$ is multiplicativity, and
$v(x+y)\ge\min(v(x),v(y))$ is the ultrametric inequality, $c$ being less than $1$. Moreover $\mathcal O=\{v\ge0\}$, and every nonzero
$x\in\mathcal O$ is $\varpi^{v(x)}u$ with $u=x\varpi^{-v(x)}$ of absolute value $1$, a unit. If $I\ne0$ is an ideal and $n$ the least value of
$v$ on $I\setminus\{0\}$, then $I$ contains some $\varpi^nu$, hence $\varpi^n$; and every $x\in I$ is $\varpi^n\cdot x\varpi^{-n}$ with
$x\varpi^{-n}\in\mathcal O$. So $I=\varpi^n\mathcal O$, and $\mathcal O$ is a principal ideal domain with the single nonzero prime ideal
$\varpi\mathcal O$, a discrete valuation ring (\Cref{ch:modules}).

(iv) If $\abs{\cdot}_1=\abs{\cdot}_2^s$ the open balls coincide. Conversely, in the topology of $\abs{\cdot}_i$ one has $\abs x_i<1$ iff
$x^n\to0$; so equivalent absolute values have the same set $\{\abs x<1\}$. Fix $y$ with $\abs y_1>1$ (it exists, the absolute value being
nontrivial), so $\abs y_2>1$ too, and put $s=\log\abs y_1/\log\abs y_2$. For $x\ne0$ write $\abs x_1=\abs y_1^{\alpha}$. For integers $m$, $n>0$
with $m/n>\alpha$ we have $\abs{x^n/y^m}_1<1$, hence $\abs{x^n/y^m}_2<1$, that is $\abs x_2<\abs y_2^{m/n}$; likewise $m/n<\alpha$ gives
$\abs x_2>\abs y_2^{m/n}$. So $\abs x_2=\abs y_2^{\alpha}$ and $\abs x_1=\abs x_2^{s}$.
\end{proof}

\begin{lemma}[summation in a complete nonarchimedean field]\label{lem:local:summation}
Let $F$ be complete for a nonarchimedean absolute value and $(a_i)_{i\in I}$ a countable family in $F$ such that for every $\delta>0$ only
finitely many $\abs{a_i}$ exceed $\delta$. Then $\sum_ia_i$ converges for every enumeration of $I$, to a sum independent of the
enumeration; and for every partition $I=\bigsqcup_jI_j$ the partial sums $s_j=\sum_{i\in I_j}a_i$ satisfy the same hypothesis and
$\sum_js_j=\sum_ia_i$.
\end{lemma}

\begin{proof}
For an enumeration, the terms tend to $0$, so the series converges by \Cref{prop:local:absval}(ii). Given $\delta>0$, let $I_\delta$ be the
finite set of indices with $\abs{a_i}>\delta$. Any partial sum of any enumeration that contains $I_\delta$ differs from
$\sum_{i\in I_\delta}a_i$ by a finite sum of terms of absolute value at most $\delta$, hence by at most $\delta$; so two limits differ by at most
$\delta$, for every $\delta$. For a partition, $\abs{s_j}\le\max_{i\in I_j}\abs{a_i}$, so only finitely many $\abs{s_j}$ exceed $\delta$, and
$\sum_js_j$ differs from $\sum_{i\in I_\delta}a_i$ by at most $\delta$ once the finitely many $j$ meeting $I_\delta$ are included.
\end{proof}

\begin{remark}[Ostrowski, product formula, weak approximation]\label{rem:local:ostrowski-product}
By Ostrowski's theorem (\Cref{ch:rings}) every nontrivial absolute value on $\Q$ is equivalent to $\abs{\cdot}_\infty$ or to $\abs{\cdot}_p$ for
exactly one prime $p$, and $\prod_{v\le\infty}\abs x_v=1$ for $x\in\Q^\times$ (the product formula, also \Cref{ch:rings}). For finitely many
pairwise inequivalent absolute values $\abs{\cdot}_1,\dots,\abs{\cdot}_r$ on a field $F$ and elements $a_1,\dots,a_r$, there is for every
$\delta>0$ an $x\in F$ with $\abs{x-a_i}_i<\delta$ for all $i$ (weak approximation). We do not prove this; see Neukirch
\cite[Chapter II, Section 3]{Neukirch}. It is not used in the book.
\end{remark}

\subsection{Completions}

\begin{theorem}[completion]\label{thm:local:completion}
Let $F$ be a field with an absolute value $\abs{\cdot}$. There is a field $\hat F$ containing $F$, with an absolute value extending
$\abs{\cdot}$, which is complete and in which $F$ is dense; it is unique up to a unique isometric $F$-isomorphism. If $\abs{\cdot}$ is
nonarchimedean, then:
\begin{enumerate}[label=(\roman*)]
\item $\abs{\hat F^\times}=\abs{F^\times}$, and the extended absolute value is nonarchimedean;
\item the valuation ring $\hat{\mathcal O}$ of $\hat F$ is the closure of $\mathcal O$, and if $\abs{\cdot}$ is discrete with uniformiser $\varpi$,
then $\mathcal O/\varpi^n\mathcal O\to\hat{\mathcal O}/\varpi^n\hat{\mathcal O}$ is an isomorphism for every $n\ge1$ (so the residue field does not change);
\item if $\abs{\cdot}$ is discrete, $\hat{\mathcal O}\to\varprojlim_n\mathcal O/\varpi^n\mathcal O$, $x\mapsto(x\bmod\varpi^n)_n$, is an isomorphism of
topological rings, the right side carrying the inverse limit of the discrete topologies.
\end{enumerate}
\end{theorem}

\begin{proof}
Let $C$ be the ring of Cauchy sequences in $F$ (termwise operations) and $\mathfrak n\subset C$ the ideal of sequences tending to $0$. If
$(x_n)\in C\setminus\mathfrak n$, then $\abs{x_n}$ is bounded below by some $\delta>0$ for $n\ge n_0$ (otherwise a subsequence tends to $0$ and so
does the Cauchy sequence); then $(x_n^{-1})_{n\ge n_0}$ is Cauchy, since $\abs{x_n^{-1}-x_m^{-1}}\le\delta^{-2}\abs{x_m-x_n}$. So $\mathfrak n$ is a
maximal ideal and $\hat F=C/\mathfrak n$ is a field; $F$ embeds by constant sequences. Put $\abs{(x_n)}=\lim\abs{x_n}$, which exists since
$\bigl\lvert\abs{x_n}-\abs{x_m}\bigr\rvert\le\abs{x_n-x_m}$, and depends only on the class modulo $\mathfrak n$. It is an absolute value extending
$\abs{\cdot}$: multiplicativity and the triangle inequality hold term by term, $\abs{x_ny_n}=\abs{x_n}\abs{y_n}$ and $\abs{x_n+y_n}\le\abs{x_n}+\abs{y_n}$,
and survive the limit; $\lim\abs{x_n}=0$ means $(x_n)\in\mathfrak n$, the zero class; and a constant sequence has the absolute value of its term. The class of $(x_n)$ is the limit of the constants $x_n$, so $F$ is dense; and a Cauchy sequence of classes converges to the
class of a diagonal sequence of approximating elements of $F$, so $\hat F$ is complete. If $F'$ is another such field, the identity of $F$
extends uniquely to an isometry $\hat F\to F'$ by continuity (send $\lim x_n$ to $\lim x_n$), which is a field isomorphism.

(i) Let $\hat x=\lim x_n\ne0$ with $x_n\in F$. For $n$ large, $\abs{x_n-\hat x}<\abs{\hat x}$, so $\abs{x_n}=\abs{\hat x}$ by
\Cref{prop:local:absval}(ii) applied in $\hat F$, where the ultrametric inequality holds by passing to the limit. Hence the value
groups agree.

(ii) If $\abs{\hat x}\le1$ and $x_n\to\hat x$ with $x_n\in F$, then $\abs{x_n}\le1$ for $n$ large; so $\hat{\mathcal O}$ is the closure of $\mathcal O$,
and $\hat{\mathcal O}\cap F=\mathcal O$. By (i), $\varpi$ is a uniformiser of $\hat F$. The map $\mathcal O\to\hat{\mathcal O}/\varpi^n\hat{\mathcal O}$ is
surjective, because every $\hat x\in\hat{\mathcal O}$ has an $x\in F$ with $\abs{x-\hat x}\le\abs{\varpi}^n$, and then $x\in\mathcal O$; its kernel is
$\mathcal O\cap\varpi^n\hat{\mathcal O}=\varpi^n\mathcal O$.

(iii) By (ii), $\hat{\mathcal O}/\varpi^n\hat{\mathcal O}=\mathcal O/\varpi^n\mathcal O$. The map is injective because $\bigcap_n\varpi^n\hat{\mathcal O}=0$. It is
surjective: given a compatible family $(x_n\bmod\varpi^n)$ with $x_n\in\mathcal O$, the sequence $(x_n)$ is Cauchy, since $x_{m}-x_n\in\varpi^n\mathcal O$
for $m\ge n$, and its limit maps to the family. The ideals $\varpi^n\hat{\mathcal O}$ are a basis of neighbourhoods of $0$ on the left and
correspond to the kernels of the projections on the right, so the bijection is a homeomorphism.
\end{proof}

\subsection{The field \texorpdfstring{$\Q_p$}{Qp}}

\begin{definition}\label{def:local:qp}
$\Q_p$ is the completion of $\Q$ for $\abs{\cdot}_p$, and $\Z_p=\{x\in\Q_p:\abs x_p\le1\}$ is the ring of \emph{$p$-adic integers}. The
valuation $v_p$ extends to $\Q_p^\times$ by $\abs x_p=p^{-v_p(x)}$ (\Cref{thm:local:completion}(i)).
\end{definition}

\begin{theorem}\label{thm:local:qp}
\begin{enumerate}[label=(\roman*)]
\item $\Z_p$ is the closure of $\Z$, and of $\Z_{(p)}$, in $\Q_p$. It is a complete discrete valuation ring with uniformiser $p$, residue field
$\F_p$ and nonzero ideals $p^n\Z_p$.
\item The map $\Z_p\to\varprojlim_n\Z/p^n\Z$, $x\mapsto(x\bmod p^n)_n$, is an isomorphism of topological rings. So the $\Z_p$ of this chapter is the
ring $\Z_\ell$ built algebraically in \Cref{ch:modules} (at $\ell=p$), and the profinite group of \Cref{ch:groups}.
\item $\Z_p^\times=\{\abs x_p=1\}$, and $x\in\Z_p$ is a unit iff $x\not\equiv0\pmod p$. Every $x\in\Q_p^\times$ is uniquely $x=p^{v_p(x)}u$ with
$u\in\Z_p^\times$; so $\Q_p^\times=p^{\Z}\times\Z_p^\times$ and $\Q_p=\Z_p[1/p]=\Frac\Z_p$.
\item $\Z_p$ is compact; $\Q_p$ is locally compact and totally disconnected.
\item For a prime $q\ne p$, $q\in\Z_p^\times$, so $q^{-1}\in\Z_p$, and $q^{a}$ is defined for every $a\in\Z_p$ when $q\in1+p\Z_p$
(\Cref{thm:local:units-structure}).
\end{enumerate}
\end{theorem}

\begin{proof}
(i) The absolute value $\abs{\cdot}_p$ on $\Q$ is discrete with value group $p^{\Z}$ and valuation ring $\Z_{(p)}$, whose residue field is $\F_p$
(\Cref{ch:rings}). By \Cref{thm:local:completion}, $\Z_p$ is the closure of $\Z_{(p)}$ and a complete discrete valuation ring with uniformiser $p$
and residue field $\Z_{(p)}/p\Z_{(p)}=\F_p$, and its nonzero ideals are $p^n\Z_p$ (\Cref{prop:local:absval}(iii)). $\Z$ is dense in $\Z_{(p)}$:
for $a/b\in\Z_{(p)}$ and $n\ge1$, $b$ is invertible modulo $p^n$, and an integer $c$ with $bc\equiv a\pmod{p^n}$ has $\abs{a/b-c}_p\le p^{-n}$.
(ii) is \Cref{thm:local:completion}(iii) with $\mathcal O=\Z_{(p)}$, $\varpi=p$, since $\Z_{(p)}/p^n\Z_{(p)}=\Z/p^n\Z$.
(iii) Units of the valuation ring are the elements of absolute value $1$ (\Cref{prop:local:absval}(ii)), that is, those not in $p\Z_p$.
For $x\ne0$, $u=xp^{-v_p(x)}$ has absolute value $1$. (iv) By (ii), $\Z_p$ is homeomorphic to a closed subset of the product
$\prod_n\Z/p^n$ of finite discrete sets, which is compact (Tychonoff; \Cref{ch:groups}). $\Z_p$ is an open neighbourhood of $0$, and
$x+\Z_p$ of $x$, so $\Q_p$ is locally compact. Balls are open and closed (\Cref{prop:local:absval}(ii)) and arbitrarily small, so
the connected components are points: if a subset $S$ contains $x\ne y$, the ball $B=B(x,\abs{x-y}_p)$ contains $x$ but not $y$, and
$S=(S\cap B)\sqcup(S\setminus B)$ writes $S$ as a disjoint union of two nonempty relatively open sets. (v) $q$ is an integer not divisible by $p$.
\end{proof}

\begin{warning}\label{warn:local:analytic}
\Cref{ch:modules} defines $\Z_\ell=\varprojlim\Z/\ell^n$ and proves that it is a complete discrete valuation ring; \Cref{thm:local:qp}(ii) shows
that this is the closure of $\Z$ in the completion $\Q_\ell$ of $\Q$. Every statement about $\Z_\ell$ in either chapter therefore applies to the
other construction. The two topologies agree as well, so ``continuous'' means the same thing in both.
\end{warning}

\begin{proposition}[$p$-adic expansions]\label{prop:local:expansions}
Every $x\in\Q_p^\times$ has a unique expansion
\[
x=\sum_{i\ge v}a_ip^i,\qquad v=v_p(x),\quad a_i\in\{0,1,\dots,p-1\},\quad a_v\ne0,
\]
converging in $\Q_p$. $x\in\Z_p$ iff $v\ge0$, and for $x\in\Z_p$ the partial sum $\sum_{i<n}a_ip^i$ is the integer in $[0,p^n)$ congruent to
$x$ modulo $p^n$. The number $x$ lies in $\Q$ iff its digit sequence $(a_i)$ is eventually periodic.
\end{proposition}

\begin{proof}
Multiplying by a power of $p$ we may assume $x\in\Z_p^\times$. Let $s_n\in[0,p^n)$ be the integer congruent to $x$ modulo $p^n$
(\Cref{thm:local:qp}(ii)). Then $s_{n+1}\equiv s_n\pmod{p^n}$, so $s_{n+1}=s_n+a_np^n$ with $a_n\in\{0,\dots,p-1\}$, and $s_n\to x$; this gives
the expansion, and uniqueness holds because the partial sums of any expansion are the $s_n$. A sum $\sum a_ip^i$ with integer digits
converges by \Cref{prop:local:absval}(ii).
If the digits are periodic of period $k$ from the index $m$ on, then $x=A+p^mB/(1-p^k)$ with integers $A$, $B$, a rational number.
Conversely let $x=a/b\in\Z_p\cap\Q$ with $b\ge1$, $p\nmid b$, and put $x_n=(x-s_n)/p^n$, the number whose digits are $a_n,a_{n+1},\dots$. Since
$bs_n\equiv a\pmod{p^n}$, $bx_n=(a-bs_n)/p^n$ is an integer, and as $0\le s_n<p^n$ it lies in the interval $(a/p^n-b,\,a/p^n]$, which for all $n$
is contained in $(-\abs a-b,\abs a]$. So the $x_n$ take finitely many values, $x_n=x_m$ for some $n<m$, and the digit sequence is periodic from
$n$ on, with period $m-n$.
\end{proof}

\begin{example}\label{ex:local:expansions}
$-1=\sum_{i\ge0}(p-1)p^i$, since $s_n=p^n-1$; and $1/(1-p)=\sum_{i\ge0}p^i$. For readers who know Laurent series from Riley, Hobson and Bence: a
$p$-adic number is a ``Laurent series in $p$'' with digits in $\{0,\dots,p-1\}$, but arithmetic carries from one digit to the next, so the
digits of a sum are not the sums of the digits. In $\Z_5$, the digits of $\sqrt{-1}$ (a root of $X^2+1$ congruent to $2$ modulo $5$,
\Cref{ex:local:hensel-examples}) begin $2,1,2,1$: that root is $\equiv182\pmod{5^4}$.
\end{example}

\section{Hensel's lemma; units, roots of unity and powers}\label{sec:local:hensel}

This section answers the questions the paper asks about elements of $\Q_q$ and $\Z_\ell$: when a polynomial has a root
(Hensel's lemma), which roots of unity lie in $\Q_p$, what $\Z_p^\times$ and $\Q_p^\times$ look like, and when an element is an
$\ell^k$-th power. The paper uses the last answer in Lemma lem:tamfactor, where ``the reduction of $u$ is an $\ell^k$-th power in
$\F_q^\times$'' is equivalent to $u$ being an $\ell^k$-th power in $\Z_q^\times$ (\Cref{thm:local:powers}(ii)), and in Lemma
lem:tamagawa-exp, through the units of $\Q_q^{\ur}$ (\Cref{prop:local:qpur}(iv)). The rings $\Z_\ell[\zeta_3]$ and $\Z_\ell[i]$ of Lemma
lem:euler are decided here too (\Cref{ex:local:hensel-examples}(g)).

\subsection{Hensel's lemma}

\begin{theorem}[Hensel's lemma]\label{thm:local:hensel}
Let $K$ be a field complete for a nonarchimedean absolute value, with valuation ring $\mathcal O$ and maximal ideal $\mathfrak m$, and let
$f\in\mathcal O[X]$.
\begin{enumerate}[label=(\alph*)]
\item (Newton form.) If $a_0\in\mathcal O$ satisfies $\abs{f(a_0)}<\abs{f'(a_0)}^2$, then the Newton sequence $a_{n+1}=a_n-f(a_n)/f'(a_n)$ is
defined, lies in $\mathcal O$ and converges to a root $a$ of $f$. It is the unique root of $f$ with $\abs{a-a_0}<\abs{f'(a_0)}$; moreover
$\abs{a-a_0}=\abs{f(a_0)/f'(a_0)}$ and $\abs{f'(a)}=\abs{f'(a_0)}$.
\item (Simple-root form.) If $a_0\in\mathcal O$ satisfies $f(a_0)\equiv0\pmod{\mathfrak m}$ and $f'(a_0)\not\equiv0\pmod{\mathfrak m}$, there is a
unique $a\in\mathcal O$ with $f(a)=0$ and $a\equiv a_0\pmod{\mathfrak m}$.
\end{enumerate}
\end{theorem}

\begin{proof}
(a) Taylor's formula in $\mathcal O[x,h]$ reads $f(x+h)=f(x)+f'(x)h+h^2g(x,h)$ and $f'(x+h)=f'(x)+hg_1(x,h)$ with $g,g_1\in\mathcal O[x,h]$: for
$f=\sum c_kX^k$, expand each $(x+h)^k=x^k+kx^{k-1}h+\sum_{j\ge2}\binom kjx^{k-j}h^j$ by the binomial theorem; the integer binomial coefficients keep all
coefficients in $\mathcal O$, and no division by $j!$ occurs. The idea of the proof is the one of Newton's method over $\R$: the error is
squared at each step. Put
$c=\abs{f'(a_0)}$ and $\delta=\abs{f(a_0)}/c^2<1$. We show by induction that $a_n\in\mathcal O$, $\abs{f'(a_n)}=c$ and $\abs{f(a_n)}\le c^2\delta^{2^n}$. This
holds for $n=0$. Given it for $n$, the step $h_n=-f(a_n)/f'(a_n)$ has $\abs{h_n}=\abs{f(a_n)}/c\le c\,\delta^{2^n}<c\le1$, so $a_{n+1}=a_n+h_n\in\mathcal O$;
$\abs{f'(a_{n+1})-f'(a_n)}=\abs{h_ng_1(a_n,h_n)}\le\abs{h_n}<c$ gives $\abs{f'(a_{n+1})}=c$ (\Cref{prop:local:absval}(ii)); and since $h_n$ was chosen
so that $f(a_n)+f'(a_n)h_n=0$, Taylor's formula leaves $f(a_{n+1})=h_n^2g(a_n,h_n)$, which gives
$\abs{f(a_{n+1})}\le\abs{h_n}^2\le c^2\delta^{2^{n+1}}$. Hence $\abs{a_{n+1}-a_n}=\abs{h_n}\to0$, the sequence converges to some $a\in\mathcal O$, and
$f(a)=\lim f(a_n)=0$, $\abs{f'(a)}=c$. If $f(a_0)=0$, every $h_n$ is $0$, $a=a_0$ and $\abs{a-a_0}=0=\abs{f(a_0)/f'(a_0)}$; so assume $\delta>0$. For $n\ge1$, $\abs{h_n}\le c\delta^{2^n}<c\delta=\abs{h_0}$, so $\abs{a-a_0}=\abs{h_0}=\abs{f(a_0)/f'(a_0)}<c$. If $b$ is a
root with $\abs{b-a_0}<c$, then $h=b-a$ has $\abs h<c$, and $0=f(b)=f'(a)h+h^2g(a,h)$; if $h\ne0$ this gives $c=\abs{f'(a)}=\abs{hg(a,h)}\le\abs h<c$,
which is absurd. (b) is (a) with $\abs{f'(a_0)}=1$, since $\abs{f(a_0)}<1$ and $\abs{a-a_0}<1$ means $a\equiv a_0\pmod{\mathfrak m}$.
\end{proof}

\begin{warning}\label{warn:local:hensel}
The simple-root form needs $f'(a_0)$ to be a unit. $X^2-p$ has the root $0$ modulo $p$ but no root in $\Q_p$, since $v_p(x^2)$ is even. On the
other hand $X^2+7$ has no simple root modulo $2$ (its derivative $2X$ vanishes there), yet $-7$ is a square in $\Z_2$: with $a_0=1$,
$\abs{f(1)}_2=\abs8_2=\frac18<\frac14=\abs{f'(1)}_2^2$, so the Newton form applies, and the root is $\equiv1\pmod4$. The iteration can be followed by
hand: $a_1=1-8/2=-3$ with $f(-3)=16$, then $a_2=-3-16/(-6)=-\frac13$ with $f(-\frac13)=\frac{64}9$. So $\abs{f(a_n)}_2=\frac18,\frac1{16},\frac1{64}$,
exactly the bound $c^2\delta^{2^n}$ of the proof with $c=\frac12$, $\delta=\frac12$. The root is $\equiv181\pmod{2^8}$, and indeed
$-\frac13\equiv85\pmod{2^8}$ agrees with it modulo $2^5$, as $\abs{a-a_2}_2=\abs{f(a_2)/f'(a_2)}_2=2^{-5}$ predicts.
\end{warning}

\begin{theorem}[Hensel's lemma, factorisation form]\label{thm:local:hensel-factor}
Let $\mathcal O$ be a complete discrete valuation ring with uniformiser $\varpi$ and residue field $k$, and let $f\in\mathcal O[X]$. Suppose that its
reduction is $\bar f=\bar g\bar h$ in $k[X]$ with $\bar g$ monic and $\bar g$, $\bar h$ coprime. Then $f=gh$ with $g,h\in\mathcal O[X]$, $g$ monic of
degree $\deg\bar g$, and $g$, $h$ reducing to $\bar g$, $\bar h$.
\end{theorem}

\begin{proof}
Choose $a,b\in\mathcal O[X]$ with $\bar a\bar g+\bar b\bar h=1$, a monic lift $g_1$ of $\bar g$ of degree $m=\deg\bar g$, and a lift $h_1$ of $\bar h$. We
construct $g_k$ monic of degree $m$ and $h_k$ with $f\equiv g_kh_k\pmod{\varpi^k}$, $g_{k+1}\equiv g_k$ and $h_{k+1}\equiv h_k\pmod{\varpi^k}$. Given $g_k$,
$h_k$, write $f-g_kh_k=\varpi^kr$. In $k[X]$, divide $\bar r\bar b$ by the monic $\bar g$: $\bar r\bar b=\bar s\bar g+\bar u$ with $\deg\bar u<m$; then
$\bar r=\bar r(\bar a\bar g+\bar b\bar h)=\bar g(\bar r\bar a+\bar s\bar h)+\bar h\bar u$. Lift $\bar u$ to $u$ of degree $<m$ and $\bar r\bar a+\bar s\bar h$ to $w$,
and put $g_{k+1}=g_k+\varpi^ku$, $h_{k+1}=h_k+\varpi^kw$. Then $g_{k+1}$ is monic of degree $m$ and
$f-g_{k+1}h_{k+1}\equiv\varpi^k(r-g_kw-h_ku)\equiv0\pmod{\varpi^{k+1}}$, since $\bar g_k=\bar g$, $\bar h_k=\bar h$ and $\bar r=\bar g\bar w+\bar h\bar u$ (the
term $\varpi^{2k}uw$ is divisible by $\varpi^{k+1}$). As $g_{k+1}\equiv g_k\pmod{\varpi^k}$, the coefficients of the $g_k$ are Cauchy sequences, so they converge, and
$g=\lim g_k$ is monic of degree $m$. Divide $f$ by $g$ in $\mathcal O[X]$: $f=gh+\rho$ with $\deg\rho<m$. Since $g\equiv g_k\pmod{\varpi^k}$, we have
$gh_k\equiv f\pmod{\varpi^k}$, so $f-gh_k=\varpi^kt$ for some $t\in\mathcal O[X]$; dividing $t$ by the monic $g$ shows that the remainder of $f$ on
division by $g$, which is $\rho$, is $\equiv0\pmod{\varpi^k}$ (uniqueness of division by a monic polynomial). This holds for every $k$, so $\rho=0$
and $f=gh$. Finally $\bar g\bar h=\bar f=\bar g\,\overline{h}'$ for the reduction $\overline h'$ of $h$, and $k[X]$ is a domain, so $h$ reduces to $\bar h$.
\end{proof}

\begin{corollary}\label{cor:local:irreducible}
Let $K$ be complete for a discrete valuation, with valuation ring $\mathcal O$, and let $f=a_nX^n+\dots+a_0\in K[X]$ be irreducible with $a_0a_n\ne0$.
Then $\max_i\abs{a_i}=\max(\abs{a_0},\abs{a_n})$. In particular a monic irreducible polynomial over $K$ with constant term in $\mathcal O$ lies in
$\mathcal O[X]$.
\end{corollary}

\begin{proof}
Scale $f$ so that $\max_i\abs{a_i}=1$, and suppose $\abs{a_0}<1$ and $\abs{a_n}<1$. Let $r$ be the least index with $\abs{a_r}=1$; then $0<r<n$ and
$\bar f=X^r\bar h$ with $\bar h(0)=\bar a_r\ne0$, so $X^r$ and $\bar h$ are coprime. \Cref{thm:local:hensel-factor} gives $f=gh$ with $g$ monic of
degree $r$; then $\deg h=n-r\ge1$, contradicting irreducibility. For the last statement, $\abs{a_n}=1\ge\abs{a_0}$, so
$\max_i\abs{a_i}=\max(\abs{a_0},1)=1$, that is, every coefficient lies in $\mathcal O$.
\end{proof}

Read in reverse, the corollary is a reducibility test. For $f=X^2+p^{-1}X+1$ over $\Q_p$ the middle coefficient has $\abs{p^{-1}}_p=p$, larger
than $\abs{a_0}=\abs{a_2}=1$; so $f$ is reducible and has two roots in $\Q_p$. Directly: $pf=pX^2+X+p$ reduces to $X$ modulo $p$, and
\Cref{thm:local:hensel-factor} splits off a linear factor. The roots have valuations $1$ and $-1$ (their product is $1$ and their sum $-p^{-1}$);
for $p=3$, PARI's \texttt{polrootspadic} gives $2\cdot3+2\cdot3^2+\dots$ and $2\cdot3^{-1}+2+\dots$.

\begin{example}[Hensel's lemma at work]\label{ex:local:hensel-examples}
All residues below were computed by Newton's iteration and checked with PARI/GP (\texttt{polrootspadic}).
\begin{enumerate}[label=(\alph*)]
\item $X^2+1$ over $\Z_5$ has the simple roots $2,3$ modulo $5$; the two roots in $\Z_5$ are $\equiv182$ and $\equiv443\pmod{5^4}$, and
$182^2+1=53\cdot5^4$.
\item $\sqrt{-3}\in\Z_7$: the roots of $X^2+3$ are $\equiv37,306\pmod{7^3}$. $\sqrt2\in\Z_7$: the roots of $X^2-2$ are $\equiv108,235\pmod{7^3}$.
\item $X^2+X+1$ has the simple roots $2,4$ modulo $7$, and its roots in $\Z_7$ are $\equiv30,18\pmod{49}$. For instance, one Newton step from $a_0=2$,
where $f(2)=7$ and $f'(2)=5$, gives $a_1=2-7\cdot5^{-1}\equiv2-7\cdot10=-68\equiv30\pmod{49}$, since $5\cdot10\equiv1\pmod{49}$; and
$30^2+30+1=931=19\cdot49$. Hence
$\Z_7[\zeta_3]:=\Z_7[X]/(X^2+X+1)\simeq\Z_7\times\Z_7$ by the Chinese remainder theorem, the two roots differing by a unit.
\item (Squares in $\Z_2$.) $u\in\Z_2^\times$ is a square iff $u\equiv1\pmod8$. If $u\equiv1\pmod8$, apply \Cref{thm:local:hensel}(a) to $X^2-u$ with
$a_0=1$: $\abs{1-u}_2\le\frac18<\frac14=\abs{2}_2^2$. Conversely an odd square is $\equiv1\pmod8$. So $-7$, $17$ are squares in $\Z_2$, and $-1$, $3$, $5$ are
not.
\item (Squares in $\Z_p$, $p$ odd.) $u\in\Z_p^\times$ is a square iff $\bar u$ is a square in $\F_p$: the roots of $X^2-\bar u$ are simple as $2\bar u\ne0$.
\item (Roots of $X^m-u$.) More generally, for $p\nmid m$ and $u\in\Z_p^\times$, every root of $X^m-\bar u$ in $\F_p$ is simple and lifts uniquely.
\item (The rings $\Z_\ell[\zeta_3]$ and $\Z_\ell[i]$.) Let $g=X^2+X+1$ (resp. $X^2+1$) and $R=\Z_\ell[X]/(g)$. If $\ell\equiv1\pmod3$ (resp. $\pmod4$), $\bar g$
has two simple roots in $\F_\ell$, they lift by (b) of the theorem, and $R\simeq\Z_\ell\times\Z_\ell$, which is not a domain. If $\ell\equiv2\pmod3$
(resp. $\ell\equiv3\pmod4$), $\bar g$ is irreducible over $\F_\ell$, so $g$ is irreducible over $\Q_\ell$ (a factorisation over $\Q_\ell$ would be into
monic factors in $\Z_\ell[X]$ by \Cref{cor:local:irreducible}, and would reduce to one of $\bar g$); then $R$ is a domain, $R/\ell R=\F_\ell[X]/(\bar g)
\simeq\F_{\ell^2}$ is a field, and an element of $R$ outside $\ell R$ is a unit (it is invertible modulo $\ell$, and $R$ is $\ell$-adically complete), so $R$
is a discrete valuation ring with uniformiser $\ell$ and residue field $\F_{\ell^2}$, the ring of integers of the unramified quadratic extension of
$\Q_\ell$ (\Cref{thm:local:unramified}). For $\ell=3$ (resp. $\ell=2$), $R$ is the ring of integers of a ramified quadratic extension:
of $\Q_3(\sqrt{-3})$ by \Cref{prop:local:z3zeta3}(ii), and of $\Q_2(i)$ because $1+i$ is a root of the Eisenstein polynomial $X^2-2X+2$, so that
$\Z_2[i]=\Z_2[1+i]$ is that ring of integers by \Cref{thm:local:eisenstein} (\Cref{ex:local:ramgroups}(a)). This is the trichotomy that \Cref{ch:ant} stated for $\Z_\ell\otimes\Z[\zeta_3]$ and
$\Z_\ell\otimes\Z[i]$, now with the local input proved; it is why the paper's $\End_{I_q}(T)=\Z_\ell\otimes\Z[\zeta_3]$ is not a field for
$\ell\equiv1\pmod3$.
\end{enumerate}
\end{example}

\subsection{Teichm\"uller lifts and the structure of the units}

\begin{theorem}[Teichm\"uller lifts]\label{thm:local:teichmuller}
Let $K$ be complete for a discrete valuation with finite residue field $k=\F_q$, $q=p^f$.
\begin{enumerate}[label=(\roman*)]
\item For $a\in k$ and any lift $\tilde a\in\mathcal O$, the limit $[a]=\lim_{n\to\infty}\tilde a^{q^n}$ exists and does not depend on $\tilde a$; $[a]$ is the
unique lift of $a$ with $[a]^q=[a]$.
\item $a\mapsto[a]$ is multiplicative, $[0]=0$, and it restricts to an isomorphism $k^\times\to\mu_{q-1}(K)$ inverse to reduction. It is not
additive.
\item $\mathcal O^\times=\mu_{q-1}(K)\times U^{(1)}$ with $U^{(1)}=1+\mathfrak m$, the projection to the first factor being $u\mapsto[\bar u]$.
\end{enumerate}
\end{theorem}

\begin{proof}
First, if $x\equiv y\pmod{\mathfrak m^j}$ with $j\ge1$, then $x^p\equiv y^p\pmod{\mathfrak m^{j+1}}$: indeed $x^p-y^p=(x-y)\sum_{i=0}^{p-1}x^iy^{p-1-i}$ and
the sum is $\equiv py^{p-1}\equiv0\pmod{\mathfrak m}$. Iterating, $x^q\equiv y^q\pmod{\mathfrak m^{j+1}}$. (i) Since $a^q=a$ in $k$, $\tilde a^q\equiv\tilde a
\pmod{\mathfrak m}$, hence $\tilde a^{q^{n+1}}\equiv\tilde a^{q^n}\pmod{\mathfrak m^{n+1}}$; the sequence is Cauchy, its limit $[a]$ satisfies $[a]^q=[a]$ and
reduces to $a$. Two lifts are congruent modulo $\mathfrak m$, so their $q^n$-th powers are congruent modulo $\mathfrak m^{n+1}$, and the limits agree.
The polynomial $X^q-X$ has derivative $qX^{q-1}-1\equiv-1$, a unit, so by \Cref{thm:local:hensel}(b) it has a unique root lifting $a$.
(ii) $[a][b]$ is a root of $X^q-X$ lifting $ab$, hence equals $[ab]$. For $a\ne0$, $[a]^{q-1}=1$; conversely every $\zeta\in\mu_{q-1}(K)$ is a root
of $X^q-X$, so $\zeta=[\bar\zeta]$. In $\Z_5$, $[1]+[1]=2\ne[2]\equiv182\pmod{5^4}$. (iii) For $u\in\mathcal O^\times$, $u=[\bar u]\cdot u[\bar u]^{-1}$ with the
second factor in $U^{(1)}$, and $\mu_{q-1}(K)\cap U^{(1)}=\{1\}$ by (ii).
\end{proof}

For example, in $\Z_5$ one has $[1],[2],[3],[4]\equiv1,182,443,624\pmod{5^4}$; so $[2]$ and $[3]$ are the two square roots of $-1$ of
\Cref{ex:local:hensel-examples}(a), and $[4]=-1$.

\begin{theorem}[structure of the units]\label{thm:local:units-structure}
\begin{enumerate}[label=(\roman*)]
\item Let $p$ be odd. The series
\[
\exp(x)=\sum_{n\ge0}\frac{x^n}{n!},\qquad\log(1+x)=\sum_{n\ge1}(-1)^{n+1}\frac{x^n}n
\]
converge for $x\in p\Z_p$ and give mutually inverse isomorphisms of topological groups $p\Z_p\simeq1+p\Z_p$, carrying $p^n\Z_p$ onto $1+p^n\Z_p$
for every $n\ge1$. Hence $1+p\Z_p\simeq\Z_p$, topologically generated by $1+p$, $\Z_p^\times=\mu_{p-1}\times(1+p\Z_p)$, and
\[
\Q_p^\times=p^{\Z}\times\mu_{p-1}\times(1+p\Z_p)\simeq\Z\times\Z/(p-1)\times\Z_p .
\]
\item For $p=2$ the same series give $4\Z_2\simeq1+4\Z_2$, carrying $2^n\Z_2$ onto $1+2^n\Z_2$ for $n\ge2$; $1+4\Z_2\simeq\Z_2$ is topologically
generated by $5$, $\Z_2^\times=\{\pm1\}\times(1+4\Z_2)$ and $\Q_2^\times\simeq\Z\times\Z/2\times\Z_2$.
\item These are the inverse limits of the finite statements of \Cref{ch:rings}: $(\Z/p^k)^\times$ is cyclic for $p$ odd, and
$(\Z/2^k)^\times=\{\pm1\}\times\langle5\rangle$ for $k\ge2$.
\item Let $K$ be complete for a discrete valuation with residue field $\F_q$. Then $K^\times=\varpi^{\Z}\times\mu_{q-1}(K)\times U^{(1)}$; the
quotients $U^{(i)}/U^{(i+1)}$, $U^{(i)}=1+\mathfrak m^i$, are isomorphic to the additive group of $\F_q$; $U^{(1)}$ is an abelian pro-$p$ group,
hence a $\Z_p$-module; and $U^{(1)}$ is uniquely $m$-divisible for every $m$ prime to $p$.
\end{enumerate}
\end{theorem}

\begin{proof}
We use $v=v_p$ and two estimates for $n\ge1$: $v(n!)\le(n-1)/(p-1)$ (Legendre's formula $v(n!)=(n-s_p(n))/(p-1)$, $s_p(n)\ge1$ the digit sum,
\Cref{ch:rings}) and $v(n)\le(n-1)/(p-1)$ (if $v(n)=k$ then $n\ge p^k$ and $p^k-1=(p-1)(1+p+\dots+p^{k-1})\ge(p-1)k$). Let $c=1$ if $p$ is odd and $c=2$
if $p=2$, so that $c-1/(p-1)>0$, and let $v(x)\ge c$. Then for $n\ge2$
\[
v\Bigl(\frac{x^n}{n!}\Bigr)-v(x)\ge(n-1)\Bigl(v(x)-\frac1{p-1}\Bigr)>0,\qquad v\Bigl(\frac{x^n}n\Bigr)-v(x)\ge(n-1)\Bigl(v(x)-\frac1{p-1}\Bigr)>0 .
\]
So both series converge, $v(\exp(x)-1)=v(x)$ and $v(\log(1+x))=v(x)$; $\exp$ maps $p^n\Z_p$ into $1+p^n\Z_p$ and $\log$ maps $1+p^n\Z_p$ into
$p^n\Z_p$ for $n\ge c$, and both preserve $v$ of the deviation from $0$, resp. $1$.

The identities $\exp(X+Y)=\exp(X)\exp(Y)$, $\log((1+X)(1+Y))=\log(1+X)+\log(1+Y)$, $\exp(\log(1+X))=1+X$ and $\log(\exp X)=X$ hold in
$\Q[[X,Y]]$. They hold numerically for $v(x),v(y)\ge c$ by \Cref{lem:local:summation}. Take $\exp(\log(1+x))$: expanding
$\sum_m\frac1{m!}\bigl(\sum_n(-1)^{n+1}x^n/n\bigr)^m$ gives a family of terms $\pm x^{N}/(n_1\cdots n_m\,m!)$ with $N=n_1+\dots+n_m\ge m$, of valuation
\[
\ge Nc-\frac{N-m}{p-1}-\frac{m-1}{p-1}=N\Bigl(c-\frac1{p-1}\Bigr)+\frac1{p-1},
\]
and for each $N$ there are finitely many terms; so the family is summable. Grouping by $m$ gives $\exp(\log(1+x))$, and grouping by $N$ gives
$\sum_N\alpha_Nx^N$ where $\alpha_N$ is the coefficient of $X^N$ in $\exp(\log(1+X))=1+X$. The other three identities are proved in the same way
(the terms of $\log(\exp x)$ are $\pm\prod x^{n_i}/(n_i!\,m)$, with the same bound; those of $\log((1+x)(1+y))$ come from $(x+y+xy)^n/n$; those of
$\exp(x+y)$ are $x^ay^b/(a!\,b!)$). Hence $\exp$ and $\log$ are inverse continuous homomorphisms $p^c\Z_p\leftrightarrows1+p^c\Z_p$, matching
$p^n\Z_p$ and $1+p^n\Z_p$ for $n\ge c$.

(i) $p$ odd: $1+p\Z_p\simeq p\Z_p\simeq\Z_p$ ($x\mapsto x/p$); since $v(\log(1+p))=1$, $\log(1+p)$ generates $p\Z_p$ as a $\Z_p$-module, so $1+p$
generates $1+p\Z_p$ topologically. With \Cref{thm:local:teichmuller}(iii) and \Cref{thm:local:qp}(iii) this gives the structure of $\Q_p^\times$.
(ii) $p=2$: every $u\in\Z_2^\times$ is $\equiv\pm1\pmod4$, so exactly one of $\pm u$ lies in $1+4\Z_2$; and $v(\log5)=v(4)=2$, so $5$ generates.
(iii) Reduce modulo $p^k$: $(\Z/p^k)^\times=\Z_p^\times/(1+p^k\Z_p)\simeq\mu_{p-1}\times\Z_p/p^{k-1}\Z_p$, cyclic for $p$ odd, and $\{\pm1\}\times\Z_2/2^{k-2}$ for $p=2$.
(iv) The first statement is \Cref{thm:local:teichmuller}(iii) together with $K^\times=\varpi^{\Z}\times\mathcal O^\times$. The map
$1+\varpi^ia\mapsto\bar a$ is a homomorphism $U^{(i)}\to(\F_q,+)$ with kernel $U^{(i+1)}$, since $(1+\varpi^ia)(1+\varpi^ib)=1+\varpi^i(a+b)+\varpi^{2i}ab$; so
$U^{(1)}/U^{(n)}$ is a finite abelian $p$-group, and $U^{(1)}=\varprojlim_nU^{(1)}/U^{(n)}$ ($U^{(1)}$ is compact and the $U^{(n)}$ form a basis of open
neighbourhoods of $1$). An abelian pro-$p$ group is a $\Z_p$-module ($u^a=\lim u^{a_k}$ for integers $a_k\to a$; the limit exists because
$u^{p^N}\in U^{(n)}$ for $N$ large). For $p\nmid m$, multiplication by $m$ is an automorphism of each finite $p$-group $U^{(1)}/U^{(n)}$, compatibly in
$n$, hence of $U^{(1)}$.
\end{proof}

\begin{corollary}[roots of unity]\label{cor:local:roots-of-unity}
$\mu(\Q_p)=\mu_{p-1}$ for $p$ odd and $\mu(\Q_2)=\{\pm1\}$. For $K$ as in \Cref{thm:local:units-structure}(iv) and $m$ prime to $p$,
$\mu_m(K)=\mu_{\gcd(m,q-1)}$, and reduction is injective on it. In particular, for primes $\ell\ne p$: $\mu_{\ell^k}\subset\Q_p$ iff
$\ell^k\mid p-1$, $\#\mu_{\ell^\infty}(\Q_p)=\ell^{v_\ell(p-1)}$, and $i\in\Q_\ell$ iff $\ell\equiv1\pmod4$ (so $\Q_\ell(i)$ is a field iff
$\ell\not\equiv1\pmod4$). Also $\mu_{p^n}\subset\Q_p$ only for $p^n\in\{1,2\}$.
\end{corollary}

\begin{proof}
The torsion of $\Z\times\Z/(p-1)\times\Z_p$ is $\Z/(p-1)$, and that of $\Z\times\Z/2\times\Z_2$ is $\Z/2$. For $K$: an element of $U^{(1)}$ of order
prime to $p$ is trivial by unique divisibility, so $\mu_m(K)$ injects into $\mu_{q-1}(K)\simeq\F_q^\times$, a cyclic group of order $q-1$, whose
elements of order dividing $m$ form $\mu_{\gcd(m,q-1)}$. The remaining statements are special cases ($\ell=2$: $\mu_2\subset\Q_p$ always;
$i\in\Q_\ell$ iff $4\mid\ell-1$ for $\ell$ odd, and $i\notin\Q_2$).
\end{proof}

\subsection{Powers}

\begin{theorem}[$\ell^k$-th powers, $\ell\ne p$]\label{thm:local:powers}
Let $\ell\ne p$ be primes, $k\ge1$, and $K$ complete for a discrete valuation with residue field $\F_q$ of characteristic $p$ and uniformiser
$\varpi$ (for $K=\Q_p$: $q=p$, $\varpi=p$). Put $v=v_\ell(q-1)$.
\begin{enumerate}[label=(\roman*)]
\item $U^{(1)}$ is uniquely $\ell^k$-divisible.
\item $u\in\mathcal O_K^\times$ is an $\ell^k$-th power in $K$ iff its reduction $\bar u$ is an $\ell^k$-th power in $\F_q^\times$; this holds iff
$\bar u^{(q-1)/d}=1$, where $d=\gcd(\ell^k,q-1)$.
\item $x=\varpi^au\in K^\times$ ($u$ a unit) is an $\ell^k$-th power iff $\ell^k\mid a$ and $\bar u\in(\F_q^\times)^{\ell^k}$.
\item There is an exact sequence $1\to\F_q^\times/(\F_q^\times)^{\ell^k}\to K^\times/(K^\times)^{\ell^k}\xrightarrow{v_K}\Z/\ell^k\to0$, split by $\varpi$, so
\[
K^\times/(K^\times)^{\ell^k}\simeq\Z/\ell^k\times\Z/\ell^{\min(k,v)},\qquad\#K^\times/(K^\times)^{\ell^k}=\ell^k\cdot\#\mu_{\ell^k}(K).
\]
\item In the limit, $\widehat{K^\times}:=\varprojlim_kK^\times/(K^\times)^{\ell^k}$ sits in an exact sequence
$0\to\F_q^\times\otimes\Z_\ell\to\widehat{K^\times}\xrightarrow{v_K}\Z_\ell\to0$, split by $\varpi$; so $\widehat{K^\times}\simeq\Z_\ell\times\Z/\ell^{v}$, and
$\F_q^\times\otimes\Z_\ell=\F_q^\times/(\F_q^\times)^{\ell^{v}}\simeq\mu_{\ell^v}(\F_q)$, via $y\mapsto y^{(q-1)/\ell^v}$.
\item ($\ell=2$, $K=\Q_p$.) For $p$ odd, $\Q_p^\times/(\Q_p^\times)^2$ has order $4$, with representatives $1,u_0,p,u_0p$ for a non-square unit $u_0$;
$\Q_2^\times/(\Q_2^\times)^2$ has order $8$, with representatives $\pm1,\pm5,\pm2,\pm10$.
\end{enumerate}
\end{theorem}

\begin{proof}
(i) is \Cref{thm:local:units-structure}(iv); alternatively, \Cref{thm:local:hensel}(b) applies to $X^{\ell^k}-u$ at $a_0=1$, since $\ell^k$ is a unit.
(ii) Write $u=[\bar u]u_1$ with $u_1\in U^{(1)}$ (\Cref{thm:local:teichmuller}). If $\bar u=\bar b^{\ell^k}$, then $[\bar u]=[\bar b]^{\ell^k}$ and $u_1$ is an
$\ell^k$-th power by (i). Conversely, reduce an equation $u=y^{\ell^k}$ (then $y$ is a unit). The last equivalence is the criterion for powers in the
cyclic group $\F_q^\times$ (\Cref{ch:rings}). (iii) If $x=y^{\ell^k}$ then $a=\ell^kv_K(y)$ and $u=(y\varpi^{-v_K(y)})^{\ell^k}$; conversely apply (ii).
(iv) By (iii) the kernel of $v_K\bmod\ell^k$ is the image of $\mathcal O_K^\times$, and $\mathcal O_K^\times\cap(K^\times)^{\ell^k}=(\mathcal O_K^\times)^{\ell^k}$; by (ii)
$\mathcal O_K^\times/(\mathcal O_K^\times)^{\ell^k}\simeq\F_q^\times/(\F_q^\times)^{\ell^k}$, cyclic of order $\gcd(\ell^k,q-1)=\ell^{\min(k,v)}=\#\mu_{\ell^k}(K)$
(\Cref{cor:local:roots-of-unity}). (v) The sequences of (iv) form an inverse system of finite groups, and inverse limits of finite groups are exact
(\Cref{ch:groups}). For $k\ge v$ the left term is $\F_q^\times/(\F_q^\times)^{\ell^v}$, which is $\F_q^\times\otimes\Z_\ell$; and $y\mapsto y^{(q-1)/\ell^v}$
maps the cyclic group $\F_q^\times$ onto its subgroup of order $\ell^v$ with kernel $(\F_q^\times)^{\ell^v}$. (vi) For $p$ odd, (iv) with $\ell=2$ gives
order $4$. For $p=2$: $\Q_2^\times=2^{\Z}\times\{\pm1\}\times(1+4\Z_2)$, and the squares of $1+4\Z_2$ form $1+8\Z_2$ (\Cref{ex:local:hensel-examples}(d)),
of index $2$ with representatives $1,5$.
\end{proof}

\begin{paperlink}\label{pl:local:powers}
Lemma lem:tamfactor defines $c$ as ``the class of the reduction of $u$ in $\F_q^\times/(\F_q^\times)^{\ell^{v}}$'' and uses that $v_\ell(c)\ge k$ iff the
reduction of $u$ is an $\ell^k$-th power in $\F_q^\times$ ($k\le v$). By \Cref{thm:local:powers}(ii) this is the same as $u\in(\Z_q^\times)^{\ell^k}$, since
$\ell\ne q$. The original version wrote ``$u$ is an $\ell^k$-th power in $\F_q^\times$'' for $u\in\Z_q^\times$; it must be read through the reduction, and
the equivalence fails for $\ell=q$ (for $q$ odd, $1+q$ is a $q$-th power in $\F_q$ but not in $\Z_q$, since $1+q\Z_q\simeq\Z_q$ carries the
$q$-th powers onto $1+q^2\Z_q$). The group $\varprojlim\Q_q^\times/(\Q_q^\times)^{\ell^k}\simeq\Z_\ell\times\Z/\ell^{v}$ of (v) is the paper's
$H^1(G,\Z_\ell(1))$ (\Cref{prop:local:kummer-qp}), with torsion $\F_q^\times\otimes\Z_\ell\simeq\Z_\ell(1)/(q-1)$.
\end{paperlink}

\begin{example}\label{ex:local:powers-examples}
\begin{enumerate}[label=(\alph*)]
\item $\Q_7^\times/(\Q_7^\times)^3\simeq\Z/3\times\Z/3$. The cubes of $\F_7^\times$ are $\{1,6\}$; so $6$ is a cube in $\Q_7$, while $2$ is not, $14$ is not
($v_7(14)=1$), and $7^3\cdot6$ is.
\item $\Q_{19}^\times/(\Q_{19}^\times)^9\simeq\Z/9\times\Z/9$, as $v_3(18)=2$. Since $2$ generates $\F_{19}^\times$, the classes of $2$, $8=2^3$ and
$-1=2^9$ in $\F_{19}^\times/(\F_{19}^\times)^9\simeq\Z/9$ are $1$, $3$ and $0$, of $3$-adic valuations $0$, $1$ and (by the convention of Lemma
lem:tamfactor, $v_\ell(0)=v$) $2$.
\item For the curve \curve{11a1}, $j=-2^{12}31^3/11^5$, and the Tate parameter $q_E=j^{-1}+744j^{-2}+750420j^{-3}+\dots$ (\Cref{ch:reduction})
satisfies $v_{11}(q_E)=5$ and $q_E11^{-5}\equiv-1\pmod{11}$ (by hand: $q_E11^{-5}\equiv(j11^5)^{-1}=-(2^{12}31^3)^{-1}$, with $2^{12}\equiv4$, $31^3\equiv3$, $4\cdot3\equiv1$; PARI gives
$q_E=10\cdot11^5+2\cdot11^6+\dots$). The fifth powers of $\F_{11}^\times$ are $\{\pm1\}$, so $q_E\in(\Q_{11}^\times)^5$.
\item $\Q_5^\times/(\Q_5^\times)^3\simeq\Z/3$: every unit of $\Z_5$ is a cube, as $3\nmid4$.
\end{enumerate}
\end{example}

\section{Finite extensions of \texorpdfstring{$\Q_p$}{Qp}}\label{sec:local:ext}

The coefficient rings of the paper ($\Z_\ell[\zeta_3]$, the completions $\mathcal O_\lambda$ of rings of integers at primes $\lambda\mid\ell$) and
the fields over which the paper's curves acquire good or multiplicative reduction are finite extensions of $\Q_\ell$ or $\Q_q$. This section
extends the absolute value to them, defines the ramification index and residue degree, identifies totally ramified extensions with
Eisenstein polynomials, compares completions of number fields with the global theory of \Cref{ch:ant}, and proves that the
decomposition group of a prime of $\Qbar$ above $p$ is $G_{\Q_p}$. In the paper these facts are used in Lemma lem:IV ($\Z_3[\zeta_3]$ as a
discrete valuation ring, now the ring of integers of $\Q_3(\sqrt{-3})$), in Lemma lem:crit (``The same statements hold with
$\Z$ replaced by $\Z_\ell$ and $r$ by its $\ell$-part'') and Lemma lem:old, and in every local argument at a bad prime $q$, which is an argument about $G_{\Q_q}\subset G_\Q$.

Throughout, $K$ is complete for a discrete valuation with finite residue field, as in \Cref{conv:local:conventions}.

\begin{lemma}[equivalence of norms]\label{lem:local:norms}
Let $K$ be complete for an absolute value and $V$ a finite-dimensional $K$-vector space. A \emph{norm} on $V$ is a map $\norm{\cdot}\colon V\to\R_{\ge0}$ with
$\norm x=0$ iff $x=0$, $\norm{ax}=\abs a\norm x$ and $\norm{x+y}\le\norm x+\norm y$. Any two norms $\norm{\cdot}_1,\norm{\cdot}_2$ on $V$ are equivalent:
$c\norm x_1\le\norm x_2\le C\norm x_1$ for constants $c,C>0$. In particular $V$ is complete for every norm.
\end{lemma}

\begin{proof}
Fix a basis $e_1,\dots,e_n$ and the sup norm $\norm{\sum x_ie_i}_\infty=\max\abs{x_i}$; it suffices to compare a norm $\norm{\cdot}$ with it. Clearly
$\norm x\le C\norm x_\infty$ with $C=\sum\norm{e_i}$, since $\norm{\sum x_ie_i}\le\sum\abs{x_i}\norm{e_i}\le\max_i\abs{x_i}\cdot\sum_i\norm{e_i}$. For the other
inequality we argue by induction on $n$; for $n=1$, $\norm{x_1e_1}=\abs{x_1}\norm{e_1}$, so $c=\norm{e_1}$ works. If no $c>0$
works, there are $v_k=\sum_ix_i^{(k)}e_i$ with $\norm{v_k}_\infty=1$ and $\norm{v_k}\to0$. For each $k$ some coordinate has absolute value $1$, and one index occurs
for infinitely many $k$; passing to a subsequence and renumbering, $\abs{x_n^{(k)}}=1$
for all $k$. Put $w_k=v_k/x_n^{(k)}=u_k+e_n$ with $u_k\in W=\langle e_1,\dots,e_{n-1}\rangle$. Then $\norm{w_k}=\norm{v_k}\to0$, so $(u_k)$ is Cauchy for $\norm{\cdot}$ on
$W$ (as $\norm{u_k-u_j}=\norm{w_k-w_j}\le\norm{w_k}+\norm{w_j}$), hence for $\norm{\cdot}_\infty$ by induction, and converges in the sup norm (as $K$ is complete) to some $u\in W$, hence also for $\norm{\cdot}$.
Then $\norm{u+e_n}=\lim\norm{w_k}=0$, so $e_n=-u\in W$, which is absurd. Completeness follows: a $\norm{\cdot}$-Cauchy sequence is
$\norm{\cdot}_\infty$-Cauchy, and converges coordinatewise.
\end{proof}

\begin{theorem}[extension of the absolute value]\label{thm:local:extension-absval}
Let $L/K$ be a finite extension of degree $n$. Then
\[
\abs x_L=\abs{N_{L/K}(x)}^{1/n}
\]
is the unique absolute value on $L$ extending $\abs{\cdot}$; it is discrete, $L$ is complete for it, and $\mathcal O_L=\{\abs x_L\le1\}$ is the
integral closure of $\mathcal O_K$ in $L$. Consequently every algebraic extension of $K$ (for instance $\Qbar_p$) carries a unique absolute value
extending $\abs{\cdot}$; automorphisms over $K$ are isometries, and conjugate elements over $K$ have the same absolute value.
\end{theorem}

\begin{proof}
Let $x\in L$ have minimal polynomial $m$ over $K$, of degree $d$; then the norm of $x$ is
$N_{L/K}(x)=((-1)^dm(0))^{n/d}$. We claim: $x$ is integral over $\mathcal O_K$ iff
$m\in\mathcal O_K[X]$ iff $\abs{N_{L/K}(x)}\le1$. If $x$ is integral, so are its conjugates in a splitting field, and $m(0)$, a product of them up
to sign, is an element of $K$ integral over the integrally closed $\mathcal O_K$, so $m(0)\in\mathcal O_K$, and $m\in\mathcal O_K[X]$ by
\Cref{cor:local:irreducible}. If $m\in\mathcal O_K[X]$, then $x$ is integral (it is a root of the monic $m$), and
$\abs{N_{L/K}(x)}=\abs{m(0)}^{n/d}\le1$. Finally, if $\abs{N_{L/K}(x)}\le1$, then $\abs{m(0)}\le1$, and the monic irreducible $m$ lies in
$\mathcal O_K[X]$ by the last statement of \Cref{cor:local:irreducible}. So $\{\abs x_L\le1\}$ is the integral closure, a ring. Now $\abs{\cdot}_L$ is
multiplicative (the norm is), vanishes only at $0$, and extends $\abs{\cdot}$ since $N_{L/K}(x)=x^n$ for $x\in K$. For the ultrametric inequality, let
$x,y\in L$ with, say, $\abs y_L\le\abs x_L$ and $x\ne0$; then $\abs{x+y}_L=\abs x_L\abs{1+y/x}_L$, so it suffices
to show $\abs{1+z}_L\le1$ when $\abs z_L\le1$; and this holds because $z$ and $1$ lie in the integral closure, which is a ring.
The values lie in $\abs{K^\times}^{1/n}$, a discrete group. Uniqueness: two extensions are norms on the $K$-vector space $L$, hence equivalent
by \Cref{lem:local:norms}: $c\abs x_1\le\abs x_2\le C\abs x_1$; applied to $x^k$ this gives $c^{1/k}\abs x_1\le\abs x_2\le C^{1/k}\abs x_1$ for all $k$, so
$\abs x_1=\abs x_2$. Completeness is \Cref{lem:local:norms}. An algebraic extension is the union of its finite subextensions, on which the
extensions are compatible by uniqueness. If $\sigma$ is a $K$-automorphism, $x\mapsto\abs{\sigma x}$ is another extension, hence equals $\abs x$;
conjugates are exchanged by such automorphisms of a normal closure.
\end{proof}

The field $\Qbar_p$ is not complete; its completion $\C_p$ is algebraically closed, but we shall not need it.

\begin{theorem}[ramification index and residue degree]\label{thm:local:ef}
Let $L/K$ be finite of degree $n$, with normalised valuation $v_L$ and residue field $k_L$. Put $e=e(L/K)=v_L(\varpi_K)=[\abs{L^\times}:\abs{K^\times}]$ and
$f=f(L/K)=[k_L:k_K]$. Let $\alpha_1,\dots,\alpha_f\in\mathcal O_L$ lift a basis of $k_L$ over $k_K$ and let $\varpi_L$ be a uniformiser of $L$. Then
$\{\alpha_i\varpi_L^j:1\le i\le f,\ 0\le j<e\}$ is a basis of $\mathcal O_L$ over $\mathcal O_K$; so $\mathcal O_L$ is free of rank $n$ and $ef=n$. Both $e$
and $f$ are multiplicative in towers.
\end{theorem}

\begin{proof}
First, for $c_1,\dots,c_f\in K$ not all $0$, $v_L(\sum_ic_i\alpha_i)=e\min_iv_K(c_i)$: dividing by $c_{i_0}$ with $v_K(c_{i_0})$ minimal, the sum becomes
an element of $\mathcal O_L$ whose reduction is a nontrivial $k_K$-combination of the $\bar\alpha_i$, hence a unit. Consequently, for
$c_{ij}\in K$ not all $0$,
\[
v_L\Bigl(\sum_{i,j}c_{ij}\alpha_i\varpi_L^j\Bigr)=\min_j\Bigl(e\min_iv_K(c_{ij})+j\Bigr),
\]
because the inner sums $s_j=\sum_ic_{ij}\alpha_i\varpi_L^j$ have valuations in distinct classes modulo $e$ ($0\le j<e$), so the minimum is attained
once. Hence the $\alpha_i\varpi_L^j$ are linearly independent over $K$ (in particular $f\le n$ and $e\le n$ are finite), and an element of their
$\mathcal O_K$-span $M$ that lies in $\varpi_K^k\mathcal O_L$ has all coefficients in $\varpi_K^k\mathcal O_K$ (if $e\,v_K(c_{ij})+j\ge ek$ with $j<e$, then
$v_K(c_{ij})\ge k$). Next, $\mathcal O_L=M+\varpi_K\mathcal O_L$: the quotient $\mathcal O_L/\varpi_K\mathcal O_L=\mathcal O_L/\mathfrak m_L^e$ is filtered by the
$\mathfrak m_L^j/\mathfrak m_L^{j+1}$, $0\le j<e$, each a one-dimensional $k_L$-space spanned by $\varpi_L^j$, hence a $k_K$-space with basis the
$\bar\alpha_i\varpi_L^j$. Iterating, $\mathcal O_L=M+\varpi_K^k\mathcal O_L$ for all $k$. Given $x\in\mathcal O_L$, write $x=m_k+\varpi_K^ky_k$ with $m_k\in M$;
then $m_{k+1}-m_k\in M\cap\varpi_K^k\mathcal O_L$ has coefficients in $\varpi_K^k\mathcal O_K$, so the coefficients converge in the complete $\mathcal O_K$, and
$m_k\to m\in M$ with $x-m\in\bigcap_k\varpi_K^k\mathcal O_L=0$. So $\mathcal O_L=M$ is free with the stated basis. Every $x\in L$ has $\varpi_K^kx\in\mathcal O_L$
for $k$ large, so the basis is also a $K$-basis of $L$, and $ef=n$. Multiplicativity in towers $K\subseteq L\subseteq M$: writing
$\varpi_K=\varpi_L^{e(L/K)}u$ with $u\in\mathcal O_L^\times$ and applying $v_M$, which is $e(M/L)$ on $\varpi_L$ and $0$ on units, gives
$e(M/K)=e(M/L)e(L/K)$; and $[k_M:k_K]=[k_M:k_L][k_L:k_K]$ is the tower law for the residue fields.
\end{proof}

The two numbers measure different things: $e$ counts how much the value group grows, $f$ how much the residue field grows. In the two
extreme cases one of them is $1$: $\Q_3(\sqrt{-3})$ has $e=2$, $f=1$ (\Cref{prop:local:z3zeta3}(ii)), and $\Q_2(\sqrt5)$ has $e=1$, $f=2$
(\Cref{prop:local:quadratic-ext}(iii)); the next definition names these cases.

\begin{definition}\label{def:local:ram-types}
A finite extension $L/K$ is \emph{unramified} if $e(L/K)=1$, \emph{totally ramified} if $f(L/K)=1$, \emph{tamely ramified} if $p\nmid e(L/K)$ and
\emph{wildly ramified} if $p\mid e(L/K)$, $p$ the residue characteristic. (Residue fields are finite, hence perfect, so no separability
condition on $k_L/k_K$ is needed.) For $K$ a completion of a number field, this is the global notion of \Cref{ch:ant} for the one prime above
$\mathfrak p$ in the completion (\Cref{thm:local:completion-numberfield}).
\end{definition}

\begin{theorem}[Eisenstein polynomials]\label{thm:local:eisenstein}
An extension $L/K$ is totally ramified of degree $e$ iff $L=K(\pi)$ for a root $\pi$ of an Eisenstein polynomial $X^e+c_{e-1}X^{e-1}+\dots+c_0$
($c_i\in\mathfrak m_K$, $c_0\notin\mathfrak m_K^2$). Then the polynomial is irreducible, $\pi$ is a uniformiser of $L$ and $\mathcal O_L=\mathcal O_K[\pi]$.
Conversely the minimal polynomial over $K$ of any uniformiser of a totally ramified $L$ is Eisenstein.
\end{theorem}

\begin{proof}
Let $\pi$ be a root, $L=K(\pi)$, $d=[L:K]\le e$, and $w=v_L$. As $\pi$ is integral, $w(\pi)\ge0$; if $w(\pi)=0$, then $\pi^e=-\sum_{i<e}c_i\pi^i$ has positive
valuation, which is absurd; so $w(\pi)\ge1$. On $K$, $w=e(L/K)\,v_K$; as $v_K(c_0)=1$ and $v_K(c_i)\ge1$, we get $w(c_0)=e(L/K)$ and
$w(c_i\pi^i)\ge e(L/K)+i>e(L/K)$ for $i\ge1$. In $\pi^e=-c_0-\sum_{1\le i<e}c_i\pi^i$ the term $c_0$ therefore has strictly the smallest valuation,
so $ew(\pi)=w(\pi^e)=e(L/K)\le d\le e$ (\Cref{prop:local:absval}(ii) and \Cref{thm:local:ef}). Hence
$w(\pi)=1$, $e(L/K)=d=e$, the polynomial is irreducible and $L/K$ is totally ramified; and by \Cref{thm:local:ef} (with $f=1$, $\alpha_1=1$)
$1,\pi,\dots,\pi^{e-1}$ is a basis of $\mathcal O_L$. Conversely, if $L/K$ is totally ramified of degree $e$ and $\pi$ is a uniformiser, $1,\pi,\dots,\pi^{e-1}$
is a basis (\Cref{thm:local:ef}), so $L=K(\pi)$ has degree $e$. The conjugates of $\pi$ in a normal closure have absolute value $\abs\pi<1$
(\Cref{thm:local:extension-absval}), so the non-leading coefficients of the minimal polynomial, which are sums of products of conjugates, lie in
$\mathfrak m_K$; and $v_K(c_0)=v_K(N_{L/K}\pi)=v_L(\pi)=1$, because $\abs{N_{L/K}x}=\abs x_L^e$ and $v_L=e\,v_K$ on $K$.
\end{proof}

\begin{proposition}[examples of totally ramified extensions]\label{prop:local:z3zeta3}
\begin{enumerate}[label=(\roman*)]
\item For $n\ge1$, $\Phi_{p^n}(X+1)$ is Eisenstein, so $\Q_p(\zeta_{p^n})/\Q_p$ is totally ramified of degree $\varphi(p^n)$, with uniformiser
$1-\zeta_{p^n}$, ring of integers $\Z_p[\zeta_{p^n}]$, and $\Gal(\Q_p(\zeta_{p^n})/\Q_p)\simeq(\Z/p^n)^\times$; the residue field is $\F_p$, so the whole group acts
trivially on it (it is all inertia, in the language of \Cref{sec:local:ram}).
\item $\Q_3(\zeta_3)=\Q_3(\sqrt{-3})$ is a totally ramified quadratic extension, tamely ramified ($3\nmid2$). Its ring of integers is
$\Z_3[\zeta_3]\simeq\Z_3[X]/(X^2+X+1)\simeq\Z_3\otimes_\Z\Z[\zeta_3]$, a discrete valuation ring with uniformiser $\pi=1-\zeta_3$, residue field $\F_3$
and $\pi^2=-3\zeta_3$; it is the $\pi$-adic completion of $\Z[\zeta_3]$.
\item For $m\ge1$, $X^m-p$ is Eisenstein, so $\Q_p(p^{1/m})/\Q_p$ is totally ramified of degree $m$.
\item For $p$ odd and $p^*=(-1)^{(p-1)/2}p$, $\Q_p(\sqrt{p^*})/\Q_p$ is totally and tamely ramified of degree $2$.
\end{enumerate}
\end{proposition}

\begin{proof}
(i) $\Phi_{p^n}(X)=\sum_{i=0}^{p-1}X^{ip^{n-1}}$, so modulo $p$, $\Phi_{p^n}(X+1)=((X+1)^{p^n}-1)/((X+1)^{p^{n-1}}-1)\equiv X^{p^n}/X^{p^{n-1}}=X^{\varphi(p^n)}$, and
its constant term is $\Phi_{p^n}(1)=p$. The root $\zeta_{p^n}-1$ is then a uniformiser and $\mathcal O=\Z_p[\zeta_{p^n}-1]$ by \Cref{thm:local:eisenstein}. The
Galois group embeds in $(\Z/p^n)^\times$ (\Cref{ch:galois}) and has order $\varphi(p^n)$; as $f=1$ the residue field is $\F_p$, on which every automorphism is trivial (so, once \Cref{sec:local:ram} defines inertia,
the whole group is the inertia group, \Cref{prop:local:ramgroups}(i)).
(ii) $\Phi_3(X+1)=(X+1)^2+(X+1)+1=X^2+3X+3$ is Eisenstein at $3$, and $\sqrt{-3}=2\zeta_3+1$, because
$(2\zeta_3+1)^2=4(\zeta_3^2+\zeta_3)+1=-4+1=-3$. So $\Q_3(\zeta_3)=\Q_3(\sqrt{-3})$, of degree $2$, $X^2+X+1$ is irreducible
over $\Q_3$, the ring of integers is $\Z_3[\zeta_3-1]=\Z_3[\zeta_3]\simeq\Z_3[X]/(X^2+X+1)$, with uniformiser $\zeta_3-1=-\pi$; and
$\pi^2=1-2\zeta_3+\zeta_3^2=-3\zeta_3$. Right exactness of $\otimes$ gives $\Z_3\otimes\Z[\zeta_3]\simeq\Z_3[X]/(X^2+X+1)$, and
\Cref{thm:local:completion-numberfield}(ii) identifies this with the $\pi$-adic completion of $\Z[\zeta_3]$, $(\pi)$ being the only prime above $3$.
(iii), (iv) are \Cref{thm:local:eisenstein}: all non-leading coefficients of $X^m-p$ and of $X^2-p^*$ are divisible by $p$, and the constant
terms $-p$ and $-p^*=\mp p$ are not divisible by $p^2$. In (iv) the ramification index $2$ is prime to the odd $p$, so the extension is tame.
\end{proof}

\begin{paperlink}\label{pl:local:z3zeta3}
The paragraph after Lemma lem:IV says that ``$T_3E$ is then free of rank one over $\Z_3[X]/(X^2+X+1)\simeq\Z_3[\zeta_3]$, with $X$ acting as $\sigma$'', and
the proof of Lemma lem:tamadd uses the same ring. \Cref{ch:ant} proved that this ring is a discrete valuation ring and equals $\Z_3\otimes\Z[\zeta_3]$;
\Cref{prop:local:z3zeta3}(ii) identifies it as the ring of integers of $\Q_3(\sqrt{-3})$, a discrete valuation ring with uniformiser $\pi=1-\zeta_3$ and residue
field $\F_3$, and as the completion of $\Z[\zeta_3]$ at $(1-\zeta_3)$, which is what \Cref{ch:ant} stated with a pointer here.
\end{paperlink}

\subsection{Completions of number fields}

\begin{theorem}[completions of a number field]\label{thm:local:completion-numberfield}
Let $F$ be a number field and $p$ a prime with $p\mathcal O_F=\prod_{\mathfrak p\mid p}\mathfrak p^{e_{\mathfrak p}}$ and residue degrees $f_{\mathfrak p}$
(\Cref{ch:ant}).
\begin{enumerate}[label=(\roman*)]
\item The absolute values on $F$ extending $\abs{\cdot}_p$ are exactly the maps $\abs x_{\mathfrak p}=p^{-v_{\mathfrak p}(x)/e_{\mathfrak p}}$, one for each $\mathfrak p\mid p$,
and distinct primes give inequivalent absolute values.
\item The completion $F_{\mathfrak p}$ is a finite extension of $\Q_p$ with $e(F_{\mathfrak p}/\Q_p)=e_{\mathfrak p}$ and $f(F_{\mathfrak p}/\Q_p)=f_{\mathfrak p}$, and
its ring of integers is $\mathcal O_{F_{\mathfrak p}}=\varprojlim_n\mathcal O_F/\mathfrak p^n$.
\item $\mathcal O_F\otimes_\Z\Z_p\simeq\prod_{\mathfrak p\mid p}\mathcal O_{F_{\mathfrak p}}$ and $F\otimes_\Q\Q_p\simeq\prod_{\mathfrak p\mid p}F_{\mathfrak p}$; each
$\mathcal O_{F_{\mathfrak p}}$ is free over $\Z_p$ of rank $e_{\mathfrak p}f_{\mathfrak p}$, and $\sum_{\mathfrak p}e_{\mathfrak p}f_{\mathfrak p}=[F:\Q]$.
\end{enumerate}
\end{theorem}

\begin{proof}
(i) Let $\abs{\cdot}$ extend $\abs{\cdot}_p$. It is nonarchimedean (\Cref{prop:local:absval}(i)). Every $x\in\mathcal O_F$ satisfies a monic
equation over $\Z$, and $\abs x\le1$ follows as in the proof of \Cref{thm:local:extension-absval} (if $\abs x>1$ the leading term of the
equation dominates). So $\mathfrak p=\{x\in\mathcal O_F:\abs x<1\}$ is a prime ideal containing $p$, and elements of $\mathcal O_F\setminus\mathfrak p$ have
absolute value $1$; hence the valuation ring of $\abs{\cdot}$ contains the localisation $\mathcal O_{F,\mathfrak p}$, a discrete valuation ring
(\Cref{ch:ant}). A ring strictly between a discrete valuation ring and its fraction field contains $\varpi^{-1}$ and is the whole field; so the
valuation ring is $\mathcal O_{F,\mathfrak p}$ and $\abs x=c^{v_{\mathfrak p}(x)}$; the normalisation $\abs p=p^{-1}$ and $v_{\mathfrak p}(p)=e_{\mathfrak p}$ give
$c=p^{-1/e_{\mathfrak p}}$. Different primes give different valuation rings, hence inequivalent absolute values.
(ii) By \Cref{thm:local:completion}, $\mathcal O_{F_{\mathfrak p}}=\varprojlim\mathcal O_{F,\mathfrak p}/\mathfrak p^n=\varprojlim\mathcal O_F/\mathfrak p^n$, with residue
field $\mathcal O_F/\mathfrak p$ and $v_{\mathfrak p}(p)=e_{\mathfrak p}$; finiteness over $\Q_p$ follows from (iii).
(iii) As $\mathcal O_F\simeq\Z^n$, $\mathcal O_F\otimes\Z_p=\varprojlim_k\mathcal O_F/p^k\mathcal O_F$, and by the Chinese remainder theorem
$\mathcal O_F/p^k\mathcal O_F\simeq\prod_{\mathfrak p}\mathcal O_F/\mathfrak p^{ke_{\mathfrak p}}$ compatibly in $k$; the limit is $\prod_{\mathfrak p}\mathcal O_{F_{\mathfrak p}}$.
Each factor is a direct summand of the free module $\Z_p^n$, so it is finitely generated and torsion free over the discrete valuation ring
$\Z_p$, hence free (\Cref{ch:modules}); its rank is $e_{\mathfrak p}f_{\mathfrak p}$ by \Cref{thm:local:ef}, and the ranks add up to $n$. Inverting $p$ gives
the statement for $F\otimes\Q_p$.
\end{proof}

For example $\Z[\zeta_3]\otimes\Z_3=\Z_3[\zeta_3]$ (ramified), $\Z[\zeta_3]\otimes\Z_7\simeq\Z_7\times\Z_7$ (split, \Cref{ex:local:hensel-examples}(c)), and
$\Z[\zeta_3]\otimes\Z_2$ is the ring of integers of the unramified quadratic extension of $\Q_2$ (inert). In each case
$\Z[\zeta_3]\otimes\Z_p=\Z_p[X]/(X^2+X+1)$, and the three behaviours are read off from the reduction of $X^2+X+1$ modulo $p$: it is $(X-1)^2$ for $p=3$
(one prime, $e=2$, $f=1$), it has the two distinct roots $2$, $4$ for $p=7$ (two primes, $e=f=1$), and it is irreducible for $p=2$, having no root
in $\F_2$ (one prime, $e=1$, $f=2$, by \Cref{ex:local:hensel-examples}(g)). In all three, $\sum e_{\mathfrak p}f_{\mathfrak p}=2$, as (iii) requires. For the coefficient field of a newform
and a prime $\lambda\mid\ell$ the theorem gives the complete discrete valuation ring $\mathcal O_\lambda$ used in \Cref{ch:hecke,ch:modgalrep}. This is
the result that \Cref{ch:ant} stated and deferred here.

\begin{lemma}[Krasner]\label{lem:local:krasner}
Let $\alpha\in\overline K$ have distinct conjugates $\alpha=\alpha_1,\alpha_2,\dots,\alpha_d$ over $K$.
\begin{enumerate}[label=(\roman*)]
\item If $\beta\in\overline K$ satisfies $\abs{\beta-\alpha}<\abs{\alpha-\alpha_i}$ for $2\le i\le d$, then $K(\alpha)\subseteq K(\beta)$.
\item Let $f$ be the minimal polynomial of $\alpha$. There is $\delta>0$ such that every monic $g\in K[X]$ of degree $d$ whose coefficients differ from
those of $f$ by less than $\delta$ has a root $\beta$ with $K(\beta)=K(\alpha)$.
\end{enumerate}
\end{lemma}

\begin{proof}
(i) Let $\sigma\in\Gal(\overline K/K(\beta))$. Since $\sigma$ is an isometry (\Cref{thm:local:extension-absval}), $\abs{\sigma\alpha-\beta}=\abs{\sigma(\alpha-\beta)}=\abs{\alpha-\beta}$, so
$\abs{\sigma\alpha-\alpha}\le\max(\abs{\sigma\alpha-\beta},\abs{\beta-\alpha})=\abs{\alpha-\beta}<\abs{\alpha-\alpha_i}$ for $i\ge2$; as $\sigma\alpha$ is a conjugate, $\sigma\alpha=\alpha$.
So $\alpha$ is fixed by $\Gal(\overline K/K(\beta))$ and lies in $K(\beta)$ (characteristic $0$). (ii) If the coefficients of $g-f$ are smaller than $\delta$,
then $\abs{g(\alpha)}=\abs{g(\alpha)-f(\alpha)}<\delta\max(1,\abs\alpha)^d$, because $g-f=\sum_{i<d}b_iX^i$ with $\abs{b_i}<\delta$ and each
$\abs{b_i\alpha^i}$ is at most $\abs{b_i}\max(1,\abs\alpha)^d$. Writing $g=\prod_{j=1}^d(X-\beta_j)$ over $\overline K$ we have
$\abs{g(\alpha)}=\prod_j\abs{\alpha-\beta_j}$, so the smallest factor satisfies $\abs{\alpha-\beta_j}\le\abs{g(\alpha)}^{1/d}<\bigl(\delta\max(1,\abs\alpha)^d\bigr)^{1/d}$,
which is less than $\min_{i\ge2}\abs{\alpha-\alpha_i}$ once $\delta$ is small enough. By (i), $K(\alpha)\subseteq K(\beta_j)$, and $[K(\beta_j):K]\le d$ gives equality.
\end{proof}

(Here $\overline K$ denotes $\Qbar_p$, an algebraic closure of $K$.)

\begin{proposition}[decomposition groups are local Galois groups]\label{prop:local:decomposition}
\begin{enumerate}[label=(\roman*)]
\item Let $L/F$ be a finite Galois extension of number fields and $\mathfrak P\mid\mathfrak p$ primes. Then $L_{\mathfrak P}=L\cdot F_{\mathfrak p}$, the extension
$L_{\mathfrak P}/F_{\mathfrak p}$ is Galois, and restriction to $L$ is an isomorphism $\Gal(L_{\mathfrak P}/F_{\mathfrak p})\to D_{\mathfrak P}$, carrying the inertia
subgroup onto $I_{\mathfrak P}$ and Frobenius elements to Frobenius elements.
\item Let $\iota\colon\Qbar\to\Qbar_p$ be an embedding and $\bar{\mathfrak P}=\{x\in\bar\Z:\abs{\iota x}_p<1\}$ the prime of $\bar\Z$ it defines. Then
$\sigma\mapsto\iota^{-1}\sigma\iota$ is an isomorphism of profinite groups $G_{\Q_p}\to D_{\bar{\mathfrak P}}\subseteq G_\Q$, carrying $I_p$ onto $I_{\bar{\mathfrak P}}$ and
arithmetic Frobenius lifts to Frobenius elements. Another embedding $\iota\circ\tau$ ($\tau\in G_\Q$) composes this map with conjugation by $\tau^{-1}$.
\item Hence for a continuous representation $\rho$ of $G_\Q$ the restriction $\rho|_{G_{\Q_p}}$, the image $\rho(I_p)$, the invariants $V^{I_p}$ and, if
$\rho$ is unramified at $p$, $\rho(\Frob_p)$, are well defined up to conjugation.
\end{enumerate}
\end{proposition}

\begin{proof}
(i) The closure of $F$ in $L_{\mathfrak P}$ is a completion of $F$ for $\abs{\cdot}_{\mathfrak p}$ (up to equivalence, by
\Cref{thm:local:completion-numberfield}(i)), so $F_{\mathfrak p}\subseteq L_{\mathfrak P}$. The compositum $LF_{\mathfrak p}$ is finite over $F_{\mathfrak p}$, hence
complete (\Cref{thm:local:extension-absval}) and closed, and contains the dense subfield $L$; so $L_{\mathfrak P}=LF_{\mathfrak p}$. Write $L=F(\theta)$; the
minimal polynomial $g$ of $\theta$ splits in $L$, and $L_{\mathfrak P}=F_{\mathfrak p}(\theta)$ is generated by roots of $g$, so it is a splitting field of $g$
over $F_{\mathfrak p}$, hence Galois. An element $\tau$ of $\Gal(L_{\mathfrak P}/F_{\mathfrak p})$ permutes the roots of $g$, so maps $L$ to $L$; it is an isometry,
so $\tau|_L$ preserves $\mathfrak P=\{x\in\mathcal O_L:\abs x<1\}$, i.e. $\tau|_L\in D_{\mathfrak P}$; and $\tau$ is determined by $\tau|_L$ by continuity. Conversely,
for $\sigma\in D_{\mathfrak P}$, $x\mapsto\abs{\sigma x}_{\mathfrak P}$ is the absolute value of $\sigma^{-1}\mathfrak P=\mathfrak P$, so $\sigma$ is an isometry of $L$ and
extends by continuity to an automorphism of $L_{\mathfrak P}$ over $F_{\mathfrak p}$. The residue fields of $L$ at $\mathfrak P$ and of $L_{\mathfrak P}$ coincide
(\Cref{thm:local:completion}(ii)), compatibly with the actions; so inertia corresponds to inertia and Frobenius to Frobenius.
(ii) $\iota(\Qbar)$ is the set of elements of $\Qbar_p$ algebraic over $\Q$, so it is stable under $G_{\Q_p}$, and $\sigma\mapsto\iota^{-1}\sigma\iota$ is a
continuous homomorphism into $G_\Q$, with image in $D_{\bar{\mathfrak P}}$ since $\sigma$ is an isometry. It is injective: given $\alpha\in\Qbar_p$ with
minimal polynomial $f$ over $\Q_p$, choose $g\in\Q[X]$ monic, with coefficients close to those of $f$ ($\Q$ is dense in $\Q_p$); by
\Cref{lem:local:krasner}(ii) some root $\beta$ of $g$ has $\Q_p(\beta)=\Q_p(\alpha)$, and $\beta\in\iota(\Qbar)$. So $\Qbar_p$ is generated over $\Q_p$ by
$\iota(\Qbar)$, and an element of $G_{\Q_p}$ trivial on $\iota(\Qbar)$ is trivial. It is surjective: for each finite Galois $L/\Q$, with $\mathfrak P_L$ the
prime below $\bar{\mathfrak P}$, $G_{\Q_p}$ maps onto $\Gal(\iota(L)\Q_p/\Q_p)$, which (i) identifies with $D_{\mathfrak P_L}$; so the image is dense in
$D_{\bar{\mathfrak P}}=\varprojlim D_{\mathfrak P_L}$ (\Cref{ch:ant}), and closed, being compact. A continuous bijection of compact Hausdorff groups is a
homeomorphism. Inertia and Frobenius correspond at every finite level by (i), hence in the limit. Finally $(\iota\tau)^{-1}\sigma(\iota\tau)=
\tau^{-1}(\iota^{-1}\sigma\iota)\tau$. (iii) follows, since the primes above $p$ are conjugate (\Cref{ch:ant}).
\end{proof}

From now on a decomposition group at $p$ is identified with $G_{\Q_p}$ in this way, and we speak of $\rho|_{G_{\Q_p}}$ and $\rho(I_p)$; this is the
statement $D_p\simeq G_{\Q_p}$ that \Cref{ch:ant} deferred here.

\begin{proposition}[tame totally ramified extensions]\label{prop:local:tame-totally-ramified}
Let $L/K$ be totally ramified of degree $e$ with $p\nmid e$. Then $L=K(\alpha)$ with $\alpha^e=\varpi$ for some uniformiser $\varpi$ of $K$. The extension
is Galois iff $\mu_e\subset K$, that is iff $e\mid q-1$; then $\Gal(L/K)\to\mu_e$, $\sigma\mapsto\sigma(\alpha)/\alpha$, is an isomorphism, and
$\Gal(L/K)$ is cyclic.
\end{proposition}

\begin{proof}
Let $\pi_L$ be a uniformiser of $L$; then $\pi_L^e=\varpi_Kw$ with $w\in\mathcal O_L^\times$. As $k_L=k_K$, choose $w_0\in\mathcal O_K^\times$ with the same
reduction as $w$; then $w_1=w/w_0\in U^{(1)}_L$, which is uniquely $e$-divisible (\Cref{thm:local:units-structure}(iv)). So
$\alpha=\pi_Lw_1^{-1/e}\in L$ is a uniformiser of $L$ with $\alpha^e=\varpi_Kw_0=:\varpi$, a uniformiser of $K$. By \Cref{thm:local:eisenstein},
$[K(\alpha):K]=e=[L:K]$, so $L=K(\alpha)$. The conjugates of $\alpha$ are the $\zeta\alpha$, $\zeta\in\mu_e$. If $L/K$ is Galois they all lie in $L$, so
$\mu_e\subset L$, and then $e\mid\#k_L^\times=q-1$ (\Cref{cor:local:roots-of-unity}) and $\mu_e\subset K$. Conversely, if $\mu_e\subset K$, $L$ is the splitting
field of $X^e-\varpi$, and $\sigma\mapsto\sigma(\alpha)/\alpha$ is an injective homomorphism into $\mu_e$ (as $\sigma$ fixes $\mu_e$) between groups of order $e$.
\end{proof}

\section{Unramified extensions, \texorpdfstring{$\Q_p^{\mathrm{ur}}$}{Qp-ur} and Frobenius}\label{sec:local:unram}

The unramified extensions of a local field are generated by roots of unity of order prime to $p$, are cyclic, and are classified by
their residue fields; their union $K^{\ur}$ has Galois group $\hat\Z$, generated by the Frobenius. This is where the paper's Frobenius lives,
and where its two conventions (arithmetic and geometric) must be fixed once and for all (\Cref{conv:local:frobenius}). The paper uses the
section in Lemmas lem:tamagawa-exp and lem:tamfactor (``Over the unramified quadratic extension of $\Q_q$ (which does not change the inertia
group) $E$ is a Tate curve''; Kummer theory over $\Q_q^{\ur}$, whose units are $\ell^t$-th powers), in Lemma lem:twist(i) (invariance under unramified base change),
and in Section 1.1, where $\Frob_p$ and its inverse appear.

\begin{theorem}[unramified extensions]\label{thm:local:unramified}
Let $K$ have residue field $\F_q$.
\begin{enumerate}[label=(\roman*)]
\item For each $f\ge1$ there is exactly one unramified extension $K_f$ of $K$ of degree $f$ inside $\Qbar_p$, namely $K_f=K(\mu_{q^f-1})$.
\item $K_f/K$ is Galois, and reduction is an isomorphism $\Gal(K_f/K)\to\Gal(\F_{q^f}/\F_q)$. So $\Gal(K_f/K)$ is cyclic of order $f$, generated by
the unique $\Frob_{K_f/K}$ with $\Frob(x)\equiv x^q\pmod{\mathfrak m_{K_f}}$ for all $x\in\mathcal O_{K_f}$; on $\mu_{q^f-1}$ it is $\zeta\mapsto\zeta^q$.
\item A finite $L/K$ is unramified iff $L=K_f$ for some $f$. Composita and subextensions of unramified extensions are unramified. Every
finite $L/K$ has a largest unramified subextension $L_0$; $L_0=K_{f(L/K)}$ and $L/L_0$ is totally ramified.
\item If $L/K$ is unramified and $\alpha\in\mathcal O_L$ reduces to a generator of $k_L$ over $k_K$, then $L=K(\alpha)$.
\end{enumerate}
\end{theorem}

\begin{proof}
(i) Let $\zeta\in\Qbar_p$ be a primitive $(q^f-1)$-th root of unity and $L=K(\zeta)$. Reduction is injective on roots of unity of order prime to $p$
(\Cref{cor:local:roots-of-unity}, applied in $L$), so $\bar\zeta$ has order $q^f-1$ and generates $\F_{q^f}$ over $\F_q$; hence $f(L/K)\ge f$. The minimal
polynomial $g$ of $\zeta$ over $K$ lies in $\mathcal O_K[X]$ (\Cref{cor:local:irreducible}) and divides $X^{q^f-1}-1$, so $\bar g$ divides the separable
polynomial $X^{q^f-1}-1$ and is separable. Let $\bar h$ be the minimal polynomial of $\bar\zeta$ over $\F_q$, of degree $f$; it divides $\bar g$. If
$\bar g\ne\bar h$, then $\bar g=\bar h\bar r$ with $\bar h$, $\bar r$ coprime, and \Cref{thm:local:hensel-factor} would factor $g$. So $\bar g=\bar h$,
$[L:K]=f\le f(L/K)$, and therefore $f(L/K)=f$, $e(L/K)=1$. Uniqueness: if $L'/K$ is unramified of degree $f$, then $k_{L'}=\F_{q^f}$, and the
Teichm\"uller lifts (\Cref{thm:local:teichmuller}, applied in $L'$) give $\mu_{q^f-1}\subset L'$; so $K_f\subseteq L'$, and the degrees agree.
(ii) $K_f$ is the splitting field of $X^{q^f-1}-1$, so it is Galois. Every $\sigma\in\Gal(K_f/K)$ is an isometry, preserves $\mathcal O_{K_f}$ and
$\mathfrak m_{K_f}$, and induces an automorphism $\bar\sigma$ of $\F_{q^f}$ over $\F_q$. If $\bar\sigma=1$, then $\sigma$ fixes $\mu_{q^f-1}$ (by injectivity of
reduction on it), so $\sigma=1$. Both groups have order $f$, so $\sigma\mapsto\bar\sigma$ is an isomorphism, and $\Frob$ is the preimage of $x\mapsto x^q$;
it maps $\zeta\in\mu_{q^f-1}$ to the root of unity reducing to $\bar\zeta^q$, which is $\zeta^q$.
(iii) If $L$ is unramified with $f(L/K)=f$, the argument for uniqueness in (i) shows $L=K_f$. Subextensions are unramified since $e$ is
multiplicative; and $K_f,K_{f'}\subseteq K_{ff'}$, so composita are unramified. For any finite $L$, the Teichm\"uller lifts give
$\mu_{q^{f(L/K)}-1}\subset L$, so $K_{f(L/K)}\subseteq L$; if $K_{f'}\subseteq L$ then $\F_{q^{f'}}\subseteq k_L$, so $f'\mid f(L/K)$ and
$\mu_{q^{f'}-1}\subseteq\mu_{q^{f(L/K)}-1}$; hence $L_0=K_{f(L/K)}$, and $f(L/L_0)=1$. (iv) Let $f=[L:K]=f(L/K)$. The residue field of $K(\alpha)$
contains $\bar\alpha$, hence $k_K(\bar\alpha)=k_L$, so $[K(\alpha):K]\ge f(K(\alpha)/K)\ge[k_K(\bar\alpha):k_K]=f=[L:K]$ (the first inequality is $ef=n$);
as $K(\alpha)\subseteq L$, equality holds and $K(\alpha)=L$.
\end{proof}

\begin{definition}[maximal unramified extension]\label{def:local:qpur}
$K^{\ur}=\bigcup_{f\ge1}K_f\subseteq\Qbar_p$. Equivalently $K^{\ur}=K(\mu_m:p\nmid m)$, since $\mu_m\subset K_f$ when $m\mid q^f-1$, which holds for $f$ the
order of $q$ modulo $m$. For $K=\Q_p$ we write $\Q_p^{\ur}$.
\end{definition}

The point of passing to $K^{\ur}$ is that it removes the residue field from the picture: in $K^{\ur}$ every unit is an $m$-th power for every $m$
prime to $p$ (\Cref{prop:local:qpur}(iv) below), so the only obstruction left to extracting roots is the valuation, and what remains of the Galois
group is inertia. For example $\Q_7^{\ur}$ contains $\sqrt3$, since $3$ is not a square modulo $7$ but is one in $\F_{49}$, so $\Q_7(\sqrt3)$ is the
quadratic unramified extension $K_2$ of $\Q_7$; and $\Q_2^{\ur}$ contains $\zeta_3$, with $f=2$ the order of $2$ modulo $3$, so $\Q_2(\zeta_3)=\Q_2(\sqrt{-3})$ is
the unramified quadratic extension of $\Q_2$ (\Cref{prop:local:quadratic-ext}(iii)).

\begin{proposition}\label{prop:local:qpur}
\begin{enumerate}[label=(\roman*)]
\item The residue field of $K^{\ur}$ is $\overline\F_q$, and restriction gives isomorphisms $\Gal(K^{\ur}/K)\simeq\varprojlim_f\Gal(K_f/K)\simeq\Gal(\overline\F_q/\F_q)\simeq\hat\Z$,
topologically generated by the arithmetic Frobenius.
\item The absolute value of $K^{\ur}$ is discrete with the same value group as $K$, and $\varpi_K$ is a uniformiser; but $K^{\ur}$ is not complete.
\item Hensel's lemma (\Cref{thm:local:hensel}) holds for polynomials over $\mathcal O_{K^{\ur}}$.
\item For every $m$ prime to $p$: $\mu_m\subset K^{\ur}$; every unit of $\mathcal O_{K^{\ur}}$ is an $m$-th power in $K^{\ur}$; and $x\in(K^{\ur})^\times$ is an
$m$-th power in $K^{\ur}$ iff $m\mid v_K(x)$. In particular $\mu_{\ell^\infty}\subset\Q_q^{\ur}$ for every prime $\ell\ne q$.
\end{enumerate}
\end{proposition}

\begin{proof}
(i) The residue field is $\bigcup_f\F_{q^f}=\overline\F_q$. Restriction $\Gal(K^{\ur}/K)\to\varprojlim\Gal(K_f/K)$ is an isomorphism by infinite Galois theory
(\Cref{ch:galois}), and the isomorphisms of \Cref{thm:local:unramified}(ii) are compatible; $\Gal(\overline\F_q/\F_q)\simeq\hat\Z$ with generator
$x\mapsto x^q$ (\Cref{ch:galois}). (ii) Each $K_f$ is unramified, so has the value group of $K$. Not complete: choose $\omega_n\in\mu_{q^{n!}-1}$ whose
reduction generates $\F_{q^{n!}}$, and let $x=\sum_n\omega_n\varpi_K^n$ in the completion. If $x$ lay in some $K_f$, its expansion with Teichm\"uller digits in
$K_f$ (obtained as in \Cref{prop:local:expansions}, the digit of $y\in\mathcal O_{K_f}$ being $[\bar y]$) would be the expansion above, by uniqueness of
digits; but $\omega_n\notin K_f$ once $n!\nmid f$. (iii) The coefficients of $f$ and the approximate root $a_0$ lie in some $K_f$, which is
complete; the Newton sequence stays in $K_f$, and so does the root. (iv) $\mu_m\subset K^{\ur}$ by definition. For a unit $u$, $\bar u\in\F_{q^f}^\times$ for
some $f$; choose $f'$ with $m(q^f-1)\mid q^{f'}-1$ (possible since $p\nmid m(q^f-1)$); in the cyclic group $\F_{q^{f'}}^\times$ of order $q^{f'}-1$, the element
$\bar u$, of order dividing $q^f-1$, is an $m$-th power. The root of $X^m-\bar u$ so found is simple ($m$ is a unit), and lifts by (iii). Finally
$x=\varpi_K^{am}u$ is then an $m$-th power, and conversely $v_K(y^m)=mv_K(y)$.
\end{proof}

\begin{convention}[Frobenius]\label{conv:local:frobenius}
$\Frob=\Frob_K$ denotes an \emph{arithmetic Frobenius} of $K$: any element of $G_K$ whose restriction to $K^{\ur}$ is the generator of
\Cref{prop:local:qpur}(i), that is, which induces $x\mapsto x^q$ on $\overline\F_q$. It is well defined modulo $I_K$ (and, in a global setting, up to
conjugation). The \emph{geometric Frobenius} is $\Phi=\Frob^{-1}$. The book's $\Frob_p$ is always arithmetic (Appendix B, convention B-C2).
\begin{enumerate}[label=(\alph*)]
\item $\Frob$ acts on $\mu_m$ ($p\nmid m$) by $\zeta\mapsto\zeta^q$, so $\chi_\ell(\Frob)=q$ and $\chi_\ell(\Phi)=q^{-1}$ for $\ell\ne p$
(\Cref{prop:local:cyclotomic-local}).
\item Since $\Frob$ is defined only modulo $I_K$, a statement about $\rho(\Frob)$ for a ramified $\rho$ must be checked to be independent of the
lift, as in \Cref{cor:local:finite-tame-image} and \Cref{thm:local:kummer-inertia}(iv).
\item An eigenvalue $\lambda$ of $\Frob$ is an eigenvalue $\lambda^{-1}$ of $\Phi$. When the multiset of eigenvalues on the relevant space is stable
under inversion (as on $\ad V$ at a good prime, or on a line where Frobenius acts by $\pm1$), the two conventions give the same $L$-factor;
otherwise they do not (the Steinberg case, \Cref{prop:local:steinberg-shape}).
\item Group-theoretic statements (the tame relation, conjugation actions, the matrices in the proof of Lemma lem:tamfactor, Frobenius on Kummer
classes) use $\Frob$; local $L$-factors use $\Phi$ (\Cref{def:local:L-factor}). The revised paper says so in Section 1.1. The original version used
the arithmetic $\Frob$ in Lemma lem:tamfactor, correctly, but wrote $\Frob_p$ in Section 1.1 where its table, Lemma lem:euler and Lemma lem:fpr require
the geometric Frobenius.
\end{enumerate}
\end{convention}

\begin{proposition}[unramified base change]\label{prop:local:unram-basechange}
Let $K/\Q_p$ be finite unramified of degree $f$. Then $K^{\ur}=\Q_p^{\ur}$ and $K^{\tame}=\Q_p^{\tame}$, so $I_K=I_p$ and $P_K=P_p$ as subgroups of
$G_K\subseteq G_{\Q_p}$, and $\Frob_K\equiv\Frob_{\Q_p}^f\pmod{I_p}$. Consequently, for a representation $\rho$ of $G_{\Q_p}$, the image $\rho(I_p)$, the
invariants $V^{I_p}$ and the property of being tamely or wildly ramified are unchanged by restriction to $G_K$; so are the monodromy operator and the
Swan and Artin conductors, which are defined later (forward references: \Cref{thm:local:grothendieck}, \Cref{def:local:swan}, \Cref{def:local:conductor}).
For an arbitrary finite $K/\Q_p$, $K^{\ur}=K\Q_p^{\ur}$ and $I_K=I_p\cap G_K$ has index $e(K/\Q_p)$ in $I_p$.
\end{proposition}

\begin{proof}
For an arbitrary finite $K$: $\Q_p^{\ur}=\Q_p(\mu_m:p\nmid m)\subseteq K^{\ur}$, and $K^{\ur}=K(\mu_m:p\nmid m)\subseteq K\Q_p^{\ur}$; so $K^{\ur}=K\Q_p^{\ur}$ and
$I_K=\Gal(\Qbar_p/K\Q_p^{\ur})=G_K\cap I_p$. Its index in $I_p$ is $[K\Q_p^{\ur}:\Q_p^{\ur}]$ (\Cref{ch:galois}, translation theorem). Let $K_0$ be the largest
unramified subextension, of degree $f(K/\Q_p)$, so that $K=K_0(\pi)$ with $\pi$ a root of an Eisenstein polynomial of degree $e=e(K/\Q_p)$ over $K_0$
(\Cref{thm:local:unramified}(iii), \Cref{thm:local:eisenstein}). That polynomial stays Eisenstein over every finite unramified extension of
$K_0$, hence irreducible over $K_0^{\ur}=\Q_p^{\ur}$; so the index is $e$. If $K$ is unramified, $K=K_0\subseteq\Q_p^{\ur}$, so $K^{\ur}=\Q_p^{\ur}$. A finite
extension of $K$ is tamely ramified over $K$ iff it is over $\Q_p$, $e$ being multiplicative and $e(K/\Q_p)=1$; so $K^{\tame}=\Q_p^{\tame}$
(\Cref{thm:local:inertia-structure}(ii)). $\Frob_K$ induces $x\mapsto x^{p^f}$ on the residue field. The image $\rho(I_p)$, the invariants $V^{I_p}$ and tameness ($\rho(P_p)=1$) depend
only on the subgroups $I_p$, $P_p$, which do not change. The monodromy operator is determined by $\rho$ on an open subgroup of $I_p$ and by the
tame character, which is the same for $K$ and $\Q_p$ because $p$ is a uniformiser of both (\Cref{thm:local:tame-character}(i)). For the
conductors see \Cref{thm:local:swan}(iii). (These two forward references do not use the present proposition, so there is no circularity.)
\end{proof}

\begin{proposition}[the cyclotomic character locally]\label{prop:local:cyclotomic-local}
\begin{enumerate}[label=(\roman*)]
\item For a prime $\ell\ne p$, $\Q_p(\mu_{\ell^k})/\Q_p$ is unramified of degree the order of $p$ in $(\Z/\ell^k)^\times$, and $\Frob$ acts on $\mu_{\ell^\infty}$ by
$\zeta\mapsto\zeta^p$. So $\chi_\ell|_{G_{\Q_p}}$ is unramified with $\chi_\ell(\Frob)=p$; more generally $\chi_\ell(\Frob_K)=q$ for $K$ with residue field $\F_q$.
\item For $\ell=p$, $\Q_p(\mu_{p^n})/\Q_p$ is totally ramified with group $(\Z/p^n)^\times$, so $\chi_p(I_p)=\Z_p^\times$, and
$\Q_p(\mu_{p^\infty})\cap\Q_p^{\ur}=\Q_p$.
\item Globally, through \Cref{prop:local:decomposition}: $\chi_\ell(\Frob_p)=p$ for every prime $p\ne\ell$.
\end{enumerate}
\end{proposition}

\begin{proof}
(i) Write $\Q_{p,f}$ for the unramified extension of degree $f$. Its residue field is $\F_{p^f}$, so $\mu_{\ell^k}\subset\Q_{p,f}$ iff $\ell^k\mid p^f-1$
(\Cref{cor:local:roots-of-unity}), that is iff $f$ is a multiple of the order $f_0$ of $p$ in $(\Z/\ell^k)^\times$. The field $\Q_p(\mu_{\ell^k})$ is contained in
$\Q_{p,f_0}=\Q_p(\mu_{p^{f_0}-1})$, hence unramified (\Cref{thm:local:unramified}(iii)), so it is some $\Q_{p,f}$ with $f\mid f_0$; it contains $\mu_{\ell^k}$, so
$f_0\mid f$, and $f=f_0$. By \Cref{thm:local:unramified}(ii) $\Frob$ acts on roots of unity of order prime to $p$ by $\zeta\mapsto\zeta^p$; as
$\chi_\ell$ is defined by $\sigma(\zeta)=\zeta^{\chi_\ell(\sigma)}$ on $\mu_{\ell^\infty}$, $\chi_\ell(\Frob)=p$, and $\chi_\ell$ is trivial on $I_p$, which fixes
$\Q_p^{\ur}\supset\mu_{\ell^\infty}$. For $K$ with residue field $\F_q$ the same argument gives $\zeta\mapsto\zeta^q$. (ii) \Cref{prop:local:z3zeta3}(i): the image of $I_p$ in
$\Gal(\Q_p(\mu_{p^n})/\Q_p)$ is the inertia group, which is everything; so $\chi_p(I_p)$ is dense in $\Z_p^\times$, and closed. An extension that is both
unramified and totally ramified is trivial. (iii) follows from (i).
\end{proof}

\begin{paperlink}\label{pl:local:tamagawa-exp-ur}
The proof of Lemma lem:tamagawa-exp passes to the unramified quadratic extension $K_2$ of $\Q_q$, ``which does not change the inertia group'':
this is \Cref{prop:local:unram-basechange}. It then uses that $\Q_q^{\ur}$ contains $\mu_{\ell^t}$ ($q\ne\ell$) and that ``the units of its valuation
ring are $\ell^t$-th powers by Hensel's lemma'': \Cref{prop:local:qpur}(iv), where Hensel's lemma is applied in a finite, complete, subextension
because $\Q_q^{\ur}$ itself is not complete. The conclusion, that $q_E$ is an $\ell^t$-th power in $\Q_q^{\ur}$ if and only if $\ell^t\mid n_q$, is the last statement of
\Cref{prop:local:qpur}(iv); \Cref{thm:local:kummer-inertia}(v) puts it in terms of inertia.
\end{paperlink}

\section{Ramification groups, inertia and the tame quotient}\label{sec:local:ram}

This section is the heart of the chapter for the paper. It shows that the inertia group $I_K$ has a normal pro-$p$ subgroup $P_K$ with
procyclic quotient $\prod_{\ell\ne p}\Z_\ell(1)$, defines the tame character $t_\ell\colon I_K\to\Z_\ell(1)$, and proves the tame relation
$\Frob\,t\,\Frob^{-1}=t^q$. The paper uses these facts in Lemma lem:tamfactor (inertia acts through its maximal pro-$\ell$ quotient, generated by
$t$, and $\Frob^{-1}t\Frob=t^{q^{-1}}$), in Lemmas lem:IV and lem:euler and Proposition prop:shapes (a finite image of inertia of order prime to
$q$ is cyclic, and Frobenius acts on it by the $q$-th power), in Lemma lem:tamadd (the maximal pro-$\ell$ quotient of an open subgroup of
inertia), and in Remark rem:wild (wild inertia at $q=2,3$). \Cref{ch:ant} proved the finite-level relation for tamely ramified extensions of
number fields modulo one step, that the kernel of $\tau\mapsto\tau(\varpi)/\varpi$ is a $p$-group; that step is \Cref{prop:local:ramgroups}(iii).

\subsection{Ramification groups of a finite Galois extension}

\begin{definition}[lower ramification groups]\label{def:local:ramgroups}
Let $L/K$ be finite Galois with group $G$. For an integer $i\ge-1$ put
\[
G_i=\{\sigma\in G:\ v_L(\sigma(x)-x)\ge i+1\text{ for all }x\in\mathcal O_L\}.
\]
$G_0=I(L/K)$ is the \emph{inertia group} and $G_1=P(L/K)$ the \emph{wild inertia group} of $L/K$.
\end{definition}

The index $i$ measures how close $\sigma$ is to the identity: $\sigma\in G_i$ means that $\sigma$ moves every integer of $L$ by something divisible by
$\pi^{i+1}$. Every $\sigma$ lies in $G_{-1}$; $\sigma\in G_0$ means that $\sigma$ acts trivially on the residue field $\mathcal O_L/\mathfrak m_L$; and the higher
$G_i$ refine this. For $L=\Q_3(\sqrt{-3})$ and the nontrivial $\sigma$, the uniformiser $\pi=\sqrt{-3}$ is moved by $\sigma\pi-\pi=-2\sqrt{-3}$, of
valuation $1$; by \Cref{prop:local:ramgroups}(i) below this places $\sigma$ in $G_0$ but not in $G_1$. Wild ramification is the case in which some
$\sigma\ne1$ moves everything by at least $\pi^2$, and it can happen only when $p$ divides $e$.

\begin{proposition}\label{prop:local:ramgroups}
Let $L/K$ be finite Galois with group $G$, $\pi$ a uniformiser of $L$, $L_0$ the largest unramified subextension and $e=e(L/K)$.
\begin{enumerate}[label=(\roman*)]
\item $G_{-1}=G$; the $G_i$ are normal in $G$, decreasing, and $G_i=1$ for $i$ large. $G_0$ is the kernel of the surjection $G\to\Gal(k_L/k_K)$,
$G_0=\Gal(L/L_0)$ and $\#G_0=e$. For $\sigma\in G_0$ and $i\ge0$: $\sigma\in G_i$ iff $v_L(\sigma(\pi)-\pi)\ge i+1$. For a subgroup $H$, $H_i=H\cap G_i$.
\item $\theta_0(\sigma)=\sigma(\pi)/\pi\bmod\mathfrak m_L$ defines a homomorphism $G_0\to k_L^\times$, independent of $\pi$, with kernel $G_1$. So $G_0/G_1$ is cyclic
of order prime to $p$.
\item For $i\ge1$, $\sigma\mapsto\sigma(\pi)/\pi-1\bmod\mathfrak m_L^{i+1}$ defines a homomorphism $G_i\to\mathfrak m_L^i/\mathfrak m_L^{i+1}\simeq(k_L,+)$ with kernel
$G_{i+1}$. So $G_1$ is a $p$-group, the unique Sylow $p$-subgroup of $G_0$, and $L/K$ is tamely ramified iff $G_1=1$.
\item (Stated.) Put $G_t=G_{\lceil t\rceil}$ for real $t\ge-1$ and $\varphi_{L/K}(u)=\int_0^u dt/[G_0:G_t]$, with inverse $\psi_{L/K}$; the \emph{upper numbering}
is $G^v=G_{\psi_{L/K}(v)}$. If $H$ is normal in $G$ then $(G/H)^v=G^vH/H$ (Herbrand), and if $G$ is abelian the jumps of the upper filtration are
integers (Hasse-Arf). We do not prove this; see Serre \cite[Chapter IV, Section 3, and Chapter V, Section 7]{SerreLocalFields}.
\end{enumerate}
\end{proposition}

\begin{proof}
(i) $G_{-1}=G$ trivially. For $\tau\in G$, $v_L(\tau\sigma\tau^{-1}x-x)=v_L(\sigma(\tau^{-1}x)-\tau^{-1}x)$ since $\tau$ is an isometry, and $\tau^{-1}$ permutes $\mathcal O_L$;
so $G_i$ is normal. Every element of $L$ is $x/\varpi_K^k$ with $x\in\mathcal O_L$, so each $\sigma\ne1$ moves some $x\in\mathcal O_L$, and $G_i=1$ once $i+1$
exceeds the finitely many values $v_L(\sigma x-x)$ so obtained. By definition $G_0$ is the kernel of the action on $k_L$. An element acts trivially on
$k_L$ iff it fixes the Teichm\"uller lifts $\mu_{q^{f}-1}\subset L$, $f=f(L/K)$ (reduction is injective on them), iff it fixes $L_0=K(\mu_{q^f-1})$
(\Cref{thm:local:unramified}); so $G_0=\Gal(L/L_0)$, of order $[L:L_0]=e$, and $G/G_0=\Gal(L_0/K)\simeq\Gal(k_L/k_K)$. As $L/L_0$ is totally ramified,
$\mathcal O_L=\mathcal O_{L_0}[\pi]$ (\Cref{thm:local:eisenstein}); for $\sigma\in G_0$, which fixes $\mathcal O_{L_0}$, $\sigma(\sum c_j\pi^j)-\sum c_j\pi^j=\sum c_j(\sigma(\pi)^j-\pi^j)$ is
divisible by $\sigma(\pi)-\pi$, so $\min_xv_L(\sigma x-x)=v_L(\sigma\pi-\pi)$. The last assertion is immediate from the definition.
(ii) For $\sigma\in G_0$, $\sigma(\pi)$ is a uniformiser, so $\sigma(\pi)/\pi$ is a unit. If $\pi'=u\pi$, then $\sigma(\pi')/\pi'=(\sigma(\pi)/\pi)(\sigma(u)/u)$ and
$\sigma(u)\equiv u\pmod{\mathfrak m_L}$, so $\theta_0$ does not depend on $\pi$. Then $\sigma\tau(\pi)/\pi=\bigl(\sigma(\tau\pi)/\tau\pi\bigr)\bigl(\tau\pi/\pi\bigr)$, and the first
factor is $\theta_0(\sigma)$ computed with the uniformiser $\tau\pi$; so $\theta_0$ is a homomorphism. Its kernel is $\{v_L(\sigma\pi-\pi)\ge2\}=G_1$ by (i). As $k_L^\times$ is
cyclic of order $\#k_L-1$, prime to $p$, so is $G_0/G_1$.
(iii) For $\sigma\in G_i$ ($i\ge1$) write $\sigma(\pi)=\pi(1+a_\sigma)$, $v_L(a_\sigma)\ge i$. First, $\sigma\in G_i$ and $x\in\mathfrak m_L^j$ give $v_L(\sigma x-x)\ge i+j$: write $x=\pi^jy$, so
$\sigma x-x=\sigma(\pi^j)(\sigma y-y)+y(\sigma(\pi^j)-\pi^j)$, with $v_L(\sigma y-y)\ge i+1$ and $\sigma(\pi^j)-\pi^j=\pi^j((1+a_\sigma)^j-1)$ of valuation $\ge i+j$. Now for
$\sigma,\tau\in G_i$, $1+a_{\sigma\tau}=\frac{\sigma(\tau\pi)}{\tau\pi}\cdot\frac{\tau\pi}{\pi}$ and $\tau\pi=\pi u$ with $u=1+a_\tau$; then $\sigma(\tau\pi)/\tau\pi=(1+a_\sigma)\sigma(u)/u$ and
$v_L(\sigma(u)-u)=v_L(\sigma a_\tau-a_\tau)\ge2i\ge i+1$. Hence $1+a_{\sigma\tau}\equiv(1+a_\sigma)(1+a_\tau)$ and $a_{\sigma\tau}\equiv a_\sigma+a_\tau\pmod{\mathfrak m_L^{i+1}}$. The kernel is
$\{v_L(a_\sigma)\ge i+1\}=G_{i+1}$ by (i). The quotient $\mathfrak m_L^i/\mathfrak m_L^{i+1}$ is a one-dimensional $k_L$-vector space, of exponent $p$; so each
$G_i/G_{i+1}$ ($i\ge1$) is an elementary abelian $p$-group, and $G_1$ is a $p$-group. It is normal in $G_0$ with quotient of order prime to $p$ by
(ii), so it is the unique Sylow $p$-subgroup. If $p\nmid e=\#G_0$ then $G_1=1$; conversely, if $G_1=1$ then $\#G_0$ divides $\#k_L^\times$, prime to
$p$.
\end{proof}

\begin{example}[ramification groups of quadratic extensions]\label{ex:local:ramgroups}
Let $L/K$ be a ramified quadratic extension with $\mathcal O_L=\mathcal O_K[\pi]$, $G=\{1,\sigma\}$. By \Cref{prop:local:ramgroups}(i), $\sigma\in G_i$ iff
$v_L(\sigma\pi-\pi)\ge i+1$. In (a) to (c), $K=\Q_2$ and $e=2$, so $v_L=2v_2$ on $\Q_2$; in particular $v_L(2)=2$, and
$v_L(\sqrt d)=\frac12v_L(d)=v_2(d)$: $\sqrt d$ is a unit if $d$ is odd and a uniformiser if $v_2(d)=1$.
\begin{enumerate}[label=(\alph*)]
\item $\Q_2(i)$: $\pi=1+i$ (a root of the Eisenstein polynomial $(X-1)^2+1=X^2-2X+2$), $\sigma\pi-\pi=(1-i)-(1+i)=-2i$ has $v_L=v_L(2)=2$, $i$ being a unit; so $G_0=G_1=G$, $G_2=1$.
\item $\Q_2(\sqrt3)=\Q_2(\sqrt{-5})$ (since $-5/3=-15/9$ and $-15\equiv1\pmod8$ is a square in $\Q_2$, \Cref{ex:local:hensel-examples}(d)): $\pi=1+\sqrt3$ (root of
$(X-1)^2-3=X^2-2X-2$), $\sigma\pi-\pi=-2\sqrt3$ has $v_L=2$; again $G_1=G$, $G_2=1$.
\item $\Q_2(\sqrt d)$ for $d\in\{\pm2,\pm10\}$: $\pi=\sqrt d$ (root of the Eisenstein polynomial $X^2-d$, as $v_2(d)=1$), $\sigma\pi-\pi=-2\sqrt d$ has $v_L=2+1=3$; so $G_0=G_1=G_2=G$, $G_3=1$.
\item $\Q_p(\sqrt{p^*})$, $p$ odd, and $\Q_3(\sqrt{-3})$: $e=2$ is prime to $p$, so $G_1=1$. Likewise $\Q_p(\zeta_p)/\Q_p$ ($e=p-1$) is tame.
\end{enumerate}
These are the wildly ramified quadratic extensions of $\Q_2$ (\Cref{prop:local:quadratic-ext}(iii)); their conductor exponents are computed in
\Cref{thm:local:swan}(iv).
\end{example}

\begin{remark}[the tame relation of \Cref{ch:ant}]\label{rem:local:ant-tame}
Let $L/F$ be a finite Galois extension of number fields and $\mathfrak P\mid\mathfrak p$. By \Cref{prop:local:decomposition}(i), $I_{\mathfrak P}$ is the inertia
group $G_0$ of $L_{\mathfrak P}/F_{\mathfrak p}$, and the map $\tau\mapsto\tau(\varpi)/\varpi\bmod\mathfrak P$ used in \Cref{ch:ant} is $\theta_0$ (a uniformiser of the
localisation $\mathcal O_{L,\mathfrak P}$ is one of the completion). By \Cref{prop:local:ramgroups}(ii), (iii), its kernel is $G_1$, a $p$-group. This is the step
that \Cref{ch:ant} left unproved; so the tame relation $F\tau F^{-1}=\tau^{N\mathfrak p}$ for tamely ramified primes is now proved in full.
\end{remark}

\subsection{The inertia group of \texorpdfstring{$G_K$}{GK} and the tame quotient}

\begin{theorem}[the filtration $G_K\supset I_K\supset P_K$]\label{thm:local:inertia-structure}
\begin{enumerate}[label=(\roman*)]
\item $I_K$ is a closed normal subgroup of $G_K$ with $G_K/I_K\simeq\Gal(\overline\F_q/\F_q)\simeq\hat\Z$; for every finite Galois $L/K$, the image of $I_K$ in $\Gal(L/K)$
is $G_0=I(L/K)$.
\item Let $\varpi$ be a uniformiser of $K$ and fix a compatible system of roots $\varpi^{1/m}\in\Qbar_p$ ($p\nmid m$). A finite extension $L/K$ is tamely
ramified iff $L\subseteq K^{\ur}(\varpi^{1/m})$ for some $m$ prime to $p$. Hence the union $K^{\tame}$ of all finite tamely ramified extensions of $K$ in $\Qbar_p$
is the field $\bigcup_{p\nmid m}K^{\ur}(\varpi^{1/m})$, Galois over $K$.
\item $P_K=\Gal(\Qbar_p/K^{\tame})$ is a closed normal subgroup of $G_K$ contained in $I_K$; its image in $\Gal(L/K)$ is $G_1(L/K)$ for every finite Galois
$L/K$; it is a pro-$p$ group, and every pro-$p$ subgroup of $I_K$ is contained in it (it is the unique pro-$p$ Sylow subgroup of $I_K$).
\item (Stated.) The upper numbering of \Cref{prop:local:ramgroups}(iv) is compatible with quotients, so it defines a decreasing filtration
$(G_K^u)_{u\ge-1}$ by closed normal subgroups, with $G_K^{-1}=G_K$, $G_K^0=I_K$ and $P_K$ the closure of $\bigcup_{u>0}G_K^u$. We do not prove this;
see Serre \cite[Chapter IV, Section 3]{SerreLocalFields}.
\end{enumerate}
\end{theorem}

\begin{proof}
(i) $I_K$ is the kernel of the continuous restriction map $G_K\to\Gal(K^{\ur}/K)\simeq\hat\Z$ (\Cref{prop:local:qpur}(i)). Its image in $\Gal(L/K)$ is
$\Gal(L/L\cap K^{\ur})$ (\Cref{ch:galois}); $L\cap K^{\ur}$ is unramified and contains every unramified subextension, so it is $L_0$, and
$\Gal(L/L_0)=G_0$ by \Cref{prop:local:ramgroups}(i).
(ii) Suppose $p\nmid e(L/K)$. $L/L_0$ is totally and tamely ramified of degree $e=e(L/K)$, so $L=L_0(\alpha)$ with $\alpha^e=\varpi_0$ a uniformiser of $L_0$
(\Cref{prop:local:tame-totally-ramified}). Write $\varpi_0=\varpi u$ with $u\in\mathcal O_{L_0}^\times\subseteq\mathcal O_{K^{\ur}}^\times$; then $u=\beta^e$ with
$\beta\in K^{\ur}$ (\Cref{prop:local:qpur}(iv)), so $(\alpha/\beta)^e=\varpi$ and $\alpha/\beta=\zeta\varpi^{1/e}$ with $\zeta\in\mu_e\subset K^{\ur}$. Hence $L\subseteq
K^{\ur}(\varpi^{1/e})$. Conversely, if $L\subseteq K^{\ur}(\varpi^{1/m})$ then $L\subseteq M=K_f(\varpi^{1/m})$ for some $f$; $X^m-\varpi$ is Eisenstein over
$K_f$, so $e(M/K)=m$ is prime to $p$, and $e(L/K)\mid e(M/K)$. The fields $K^{\ur}(\varpi^{1/m})$ form a directed family ($K^{\ur}(\varpi^{1/m})$ and
$K^{\ur}(\varpi^{1/m'})$ lie in $K^{\ur}(\varpi^{1/mm'})$), so their union is a field, equal to $K^{\tame}$; it is normal over $K$ because the conjugates of
$\varpi^{1/m}$ are the $\zeta\varpi^{1/m}$, $\zeta\in\mu_m\subset K^{\ur}$.
(iii) $P_K$ is closed and normal since $K^{\tame}/K$ is Galois, and $P_K\subseteq I_K$ since $K^{\ur}\subseteq K^{\tame}$. Its image in $\Gal(L/K)$ is
$\Gal(L/L_t)$ with $L_t=L\cap K^{\tame}$, the largest tamely ramified subextension. Let $L_1=L^{G_1}$. The inertia group of $L_1/K$ is the image $G_0/G_1$
of $G_0$ (by (i), applied to $L$ and to $L_1$), of order prime to $p$; so $e(L_1/K)$ is prime to $p$ and $L_1\subseteq L_t$. Conversely the image of the
$p$-group $G_1$ in $\Gal(L_t/K)$ lies in the inertia group of $L_t/K$, of order $e(L_t/K)$ prime to $p$; so it is trivial, and $L_t\subseteq L_1$. Hence
the image of $P_K$ is $G_1$, a $p$-group, and $P_K=\varprojlim_LG_1(L/K)$ is pro-$p$. If $Q\subseteq I_K$ is pro-$p$, its image in each $\Gal(L/K)$ is a
$p$-subgroup of $G_0$, hence lies in the unique Sylow $p$-subgroup $G_1$, the image of $P_K$; as $P_K$ is closed, $Q\subseteq P_K$.
\end{proof}

\begin{corollary}[$G_K$ is a semidirect product]\label{cor:local:split}
Let $\Frob\in G_K$ be an arithmetic Frobenius lift and $Z$ the closed subgroup it generates. Then $Z\simeq\hat\Z$, $Z\cap I_K=1$ and
$G_K=I_K\rtimes Z$. Consequently a continuous character $\psi$ of $I_K$ with $\psi(\Frob\,\sigma\,\Frob^{-1})=\psi(\sigma)$ for all $\sigma\in I_K$ extends to a continuous
character $\chi$ of $G_K$, by $\chi(\sigma z)=\psi(\sigma)$ ($\sigma\in I_K$, $z\in Z$).
\end{corollary}

\begin{proof}
$Z$ is procyclic, so $n\mapsto\Frob^n$ extends to a continuous surjection $\hat\Z\to Z$ (each finite quotient of $Z$ is cyclic, generated by the
image of $\Frob$). The composite $\hat\Z\to Z\to G_K/I_K\simeq\hat\Z$ sends $1$ to $1$, hence is the identity; so $\hat\Z\to Z$ is injective, $Z\simeq\hat\Z$,
$Z\cap I_K=1$, and $Z$ maps onto $G_K/I_K$, so $G_K=I_KZ$. The map $I_K\times Z\to G_K$ is a continuous bijection of compact spaces, hence a homeomorphism.
For the extension, $(\sigma z)(\sigma'z')=\sigma(z\sigma'z^{-1})zz'$ and $\psi(z\sigma'z^{-1})=\psi(\sigma')$ ($z$ is a limit of powers of $\Frob$ and $\psi$ is continuous), so
$\chi$ is a homomorphism, continuous by the homeomorphism.
\end{proof}

\begin{theorem}[the tame character]\label{thm:local:tame-character}
Let $m$ be prime to $p$ and $\varpi$ a uniformiser of $K$.
\begin{enumerate}[label=(\roman*)]
\item $\theta_m\colon I_K\to\mu_m$, $\sigma\mapsto\sigma(\varpi^{1/m})/\varpi^{1/m}$, is a continuous surjective homomorphism. It does not depend on the root
$\varpi^{1/m}$ or on the uniformiser $\varpi$, and it does not change when $K$ is replaced by a finite unramified extension.
\item For compatible roots, $\theta_{mm'}^{m'}=\theta_m$. The limit $\theta\colon I_K\to\varprojlim_{p\nmid m}\mu_m=\prod_{\ell\ne p}\Z_\ell(1)$ is surjective with kernel
$P_K$:
\[
I_K/P_K\ \simeq\ \prod_{\ell\ne p}\Z_\ell(1).
\]
This group is abelian and procyclic, and each of its finite quotients has order prime to $p$.
\item For a prime $\ell\ne p$, the \emph{tame character} $t_\ell\colon I_K\to\Z_\ell(1)$, $\sigma\mapsto\bigl(\sigma(\varpi^{1/\ell^k})/\varpi^{1/\ell^k}\bigr)_k$, is the
$\ell$-component of $\theta$. It is surjective, every finite quotient of its kernel has order prime to $\ell$, and $t_\ell$ identifies the
maximal pro-$\ell$ quotient $I_K^{(\ell)}$ of $I_K$ with $\Z_\ell(1)$. A topological generator $t$ of $I_K^{(\ell)}$ (or of $I_K/P_K$) is called a
\emph{generator of tame inertia}.
\item Every continuous homomorphism from $I_K$ to a pro-$\ell$ group (for instance to $1+\ell\Mat_n(\Z_\ell)$, or to a group of unipotent matrices
over $\Z_\ell$) factors through $t_\ell$.
\item If $L/K$ is finite with ramification index $e$, then $I_L\subseteq I_K$ and $\theta_{m,K}|_{I_L}=\theta_{m,L}^{\,e}$; so $t_{\ell,K}|_{I_L}=e\,t_{\ell,L}$, and the two
have the same kernel.
\end{enumerate}
\end{theorem}

\begin{proof}
(i) $\mu_m\subset K^{\ur}$ is fixed by $I_K$. So for another root $\zeta\varpi^{1/m}$ the quotient is unchanged, and
$\theta_m(\sigma\tau)=\sigma(\theta_m(\tau)\varpi^{1/m})/\varpi^{1/m}=\theta_m(\tau)\theta_m(\sigma)$. If $\varpi'=u\varpi$ with $u\in\mathcal O_K^\times$, then $u^{1/m}$ can be chosen
in $K^{\ur}$ (\Cref{prop:local:qpur}(iv)), where it is fixed by $I_K$, and $(\varpi')^{1/m}=\zeta u^{1/m}\varpi^{1/m}$ gives the same character. A finite
unramified extension $K'$ has $I_{K'}=I_K$ and $\varpi$ as a uniformiser (\Cref{prop:local:unram-basechange}). $\theta_m$ factors through
$\Gal(K^{\ur}(\varpi^{1/m})/K^{\ur})$, so it is continuous. Since $X^m-\varpi$ is Eisenstein over every $K_f$, it is irreducible over $K^{\ur}$, so that Galois
group has order $m$, and it maps injectively to $\mu_m$ by $\theta_m$; as $I_K$ maps onto it, $\theta_m$ is surjective.
(ii) $(\varpi^{1/mm'})^{m'}=\varpi^{1/m}$ gives the compatibility. The kernel of $\theta$ is the subgroup fixing every $\varpi^{1/m}$, that is
$\Gal(\Qbar_p/K^{\tame})=P_K$ by \Cref{thm:local:inertia-structure}(ii). The image is closed (compact) and surjects onto every $\mu_m$, so it is
everything. By the Chinese remainder theorem $\mu_m=\prod_{\ell}\mu_{\ell^{v_\ell(m)}}$, so the limit is $\prod_{\ell\ne p}\Z_\ell(1)\simeq\prod_{\ell\ne p}\Z_\ell$, which is
procyclic. An open subgroup contains the kernel of some projection to $\mu_m$, so finite quotients are quotients of some $\mu_m$, of order prime to
$p$.
(iii) Surjectivity is (ii). The kernel of $t_\ell$ is an extension of $\prod_{\ell'\ne\ell,p}\Z_{\ell'}(1)$ by $P_K$; a finite quotient of it is an extension of a
quotient of the former (order prime to $\ell$) by a quotient of $P_K$ (a $p$-group), so it has order prime to $\ell$. If $f\colon I_K\to H$ is a continuous
homomorphism to a pro-$\ell$ group, the image of $\ker t_\ell$ in each finite quotient of $H$ is an $\ell$-group of order prime to $\ell$, hence trivial;
so $f$ kills $\ker t_\ell$ and factors through $I_K/\ker t_\ell\simeq\Z_\ell(1)$, which is itself pro-$\ell$. This proves (iii) and (iv);
$1+\ell\Mat_n(\Z_\ell)$ is pro-$\ell$ because its quotients by $1+\ell^k\Mat_n(\Z_\ell)$ have order a power of $\ell$ (\Cref{ch:groups}).
(v) $L^{\ur}=LK^{\ur}\supseteq K^{\ur}$ (as in \Cref{prop:local:unram-basechange}), so $I_L\subseteq I_K$. Write $\varpi_K=\varpi_L^eu$ with $u\in\mathcal O_L^\times$ and choose
$u^{1/m}\in L^{\ur}$; then $\varpi_K^{1/m}=\zeta(\varpi_L^{1/m})^eu^{1/m}$ for some $\zeta\in\mu_m$, and $\sigma\in I_L$ fixes $\zeta$ and $u^{1/m}$. Since $\Z_\ell(1)$ is
torsion free, $e\,t_{\ell,L}$ and $t_{\ell,L}$ have the same kernel.
\end{proof}

\begin{corollary}[open subgroups of inertia]\label{cor:local:open-inertia}
Every open subgroup $U$ of $I_K$ is $I_L$ for some finite extension $L/K$. Its maximal pro-$\ell$ quotient ($\ell\ne p$) is $\Z_\ell(1)$, through
$t_{\ell,L}$; equivalently, every continuous homomorphism from $U$ to a pro-$\ell$ group factors through $t_{\ell,K}|_U$.
\end{corollary}

\begin{proof}
The fixed field $F$ of $U$ is a finite extension of $K^{\ur}$ (\Cref{ch:galois}), generated by finitely many elements, each algebraic over $K$; let $L$
be the extension of $K$ they generate. Then $L^{\ur}=LK^{\ur}=F$ (\Cref{thm:local:tame-character}(v)), so $I_L=\Gal(\Qbar_p/F)=U$. Now apply
\Cref{thm:local:tame-character}(iii), (iv) to $L$ and (v).
\end{proof}

\begin{corollary}[a Frobenius fixing the roots of a uniformiser]\label{cor:local:frob-fixing-roots}
Let $\varpi$ be a uniformiser of $K$, $(\varpi^{1/m})_{p\nmid m}$ a compatible system of roots and $K_\varpi=K(\varpi^{1/m}:p\nmid m)$. There is an arithmetic
Frobenius lift $\Frob\in G_K$ fixing $K_\varpi$; two such lifts differ by an element of $P_K$.
\end{corollary}

\begin{proof}
Each $K(\varpi^{1/m})$ is totally ramified over $K$ (Eisenstein), so $K_\varpi\cap K^{\ur}=K$. By the translation theorem (\Cref{ch:galois}), restriction
$G_{K_\varpi}\to\Gal(K^{\ur}/K)$ is surjective, and any preimage of the Frobenius generator is the required lift. Two lifts differ by an element of
$G_{K_\varpi}\cap I_K$, which fixes every $\varpi^{1/m}$, so has $\theta=1$ and lies in $P_K$.
\end{proof}

\begin{theorem}[the tame relation]\label{thm:local:tame-relation}
\begin{enumerate}[label=(\roman*)]
\item For $\tau\in G_K$ and $\sigma\in I_K$: $\theta_m(\tau\sigma\tau^{-1})=\tau(\theta_m(\sigma))$, and $t_\ell(\tau\sigma\tau^{-1})=\chi_\ell(\tau)\,t_\ell(\sigma)$ (in the additive
notation for $\Z_\ell(1)$).
\item For every arithmetic Frobenius lift $\Frob\in G_K$ and every $t\in I_K/P_K$,
\[
\Frob\,t\,\Frob^{-1}=t^{q},\qquad\text{equivalently}\qquad\Frob^{-1}\,t\,\Frob=t^{q^{-1}},
\]
where $q^{-1}$ is the inverse of $q$ in the ring $\prod_{\ell\ne p}\Z_\ell$ of exponents of the procyclic group $I_K/P_K$. For the geometric Frobenius,
$\Phi\,t\,\Phi^{-1}=t^{q^{-1}}$. The same relations hold in $I_K^{(\ell)}\simeq\Z_\ell(1)$, with $q^{-1}\in\Z_\ell$.
\item Let $\Frob\in G_K$ be any arithmetic Frobenius lift. Then $\Gal(K^{\tame}/K)=(I_K/P_K)\rtimes\overline{\langle\Frob\rangle}\simeq\bigl(\prod_{\ell\ne p}\Z_\ell(1)\bigr)\rtimes\hat\Z$,
the generator $1\in\hat\Z$ acting by multiplication by $q$. (Stated.) As a profinite group it is generated by $t$ and $\Frob$ subject to the single
relation $\Frob\,t\,\Frob^{-1}=t^q$; see \cite[Chapter VII, Section 5]{NSW}.
\end{enumerate}
\end{theorem}

\begin{proof}
(i) Put $\alpha=\varpi^{1/m}$. Since $\tau(\alpha)$ is a root of $X^m-\varpi$, $\tau(\alpha)=\zeta_\tau^{-1}\alpha$ for some $\zeta_\tau\in\mu_m$, i.e. $\tau^{-1}(\alpha)=\tau^{-1}(\zeta_\tau)\alpha$. As
$\sigma$ fixes $\mu_m\subset K^{\ur}$,
\[
\tau\sigma\tau^{-1}(\alpha)=\tau\sigma\bigl(\tau^{-1}(\zeta_\tau)\alpha\bigr)=\tau\bigl(\tau^{-1}(\zeta_\tau)\theta_m(\sigma)\alpha\bigr)=\zeta_\tau\,\tau(\theta_m(\sigma))\,\zeta_\tau^{-1}\alpha=\tau(\theta_m(\sigma))\,\alpha .
\]
Passing to the limit over $m=\ell^k$ gives the formula for $t_\ell$, since $\tau$ acts on $\Z_\ell(1)$ through $\chi_\ell$.
(ii) $\Frob$ acts on $\mu_m$ by $\zeta\mapsto\zeta^q$ (\Cref{conv:local:frobenius}(a)), so $\theta(\Frob\,\sigma\,\Frob^{-1})=\theta(\sigma)^q$; and $\Frob^{-1}$ acts by the
inverse map $\zeta\mapsto\zeta^{q'}$, $qq'\equiv1\pmod m$, so $\theta(\Frob^{-1}\sigma\Frob)=\theta(\sigma)^{q^{-1}}$. As $\theta$ identifies $I_K/P_K$ with $\prod\Z_\ell(1)$, these are
the stated relations; $\Phi=\Frob^{-1}$ gives the last one, and projecting to the $\ell$-component gives the relations in $I_K^{(\ell)}$.
(iii) The proof of \Cref{cor:local:split} applies verbatim to the profinite group $\Gal(K^{\tame}/K)=G_K/P_K$, its closed normal subgroup
$\Gal(K^{\tame}/K^{\ur})=I_K/P_K$ (with quotient $G_K/I_K\simeq\hat\Z$) and the image of $\Frob$: the closure of $\langle\Frob\rangle$ in $G_K/P_K$ is isomorphic to $\hat\Z$, meets
$I_K/P_K$ trivially, and $\Gal(K^{\tame}/K)$ is the semidirect product of $I_K/P_K$ by it; the action is (ii).
\end{proof}

\begin{warning}[the sign in the tame relation]\label{warn:local:tame-sign}
With the arithmetic Frobenius, conjugation by $\Frob$ raises tame inertia to the power $q$; with the geometric one, $\Phi\,t\,\Phi^{-1}=t^{q^{-1}}$. The
paper writes ``$\Frob^{-1}t\Frob=t^{q^{-1}}$'' in the proof of Lemma lem:tamfactor, with $\Frob$ arithmetic: this is the second form in
\Cref{thm:local:tame-relation}(ii), and it is correct. Every $q$ prime to $m\in\{3,4,6\}$ satisfies $q\equiv\pm1\pmod m$, hence $q^{-1}\equiv q\pmod m$; so the
distinction is invisible on inertia images of order $3$, $4$ or $6$ (\Cref{cor:local:finite-tame-image}), but not on $\Z_\ell(1)$, where it decides
whether $\Frob$ acts on $H^1(I_q,\cdot)$ through $q$ or through $q^{-1}$.
\end{warning}

\begin{corollary}[finite tame images of inertia]\label{cor:local:finite-tame-image}
Let $\rho\colon G_K\to H$ be a continuous homomorphism to a Hausdorff topological group with $\rho(I_K)$ finite of order prime to $p$ (for example
$\rho\colon G_K\to\GL_n(\Q_\ell)$ continuous with such an image of inertia). Then $\rho(P_K)=1$; $\rho|_{I_K}$ factors through $I_K/P_K$, so $\rho(I_K)$ is cyclic,
generated by $\rho(t)$ for any generator $t$ of $I_K/P_K$; and for every arithmetic Frobenius lift and every $h\in\rho(I_K)$,
\[
\rho(\Frob)\,h\,\rho(\Frob)^{-1}=h^q .
\]
If $\rho(I_K)=\langle g\rangle$ has order $m$, then $\rho(\Frob)$ commutes with $g$ iff $q\equiv1\pmod m$, and it conjugates $g$ to $g^{-1}$ iff
$q\equiv-1\pmod m$.
\end{corollary}

\begin{proof}
The kernel of $\rho|_{I_K}$ is closed of finite index, hence open, so $N=\ker\rho|_{I_K}$ is an open normal subgroup of $I_K$. The image of $P_K$ in the
finite group $I_K/N$ is a quotient of $P_K$ by an open subgroup, a $p$-group (\Cref{thm:local:inertia-structure}(iii)); its order divides $\#\rho(I_K)$,
which is prime to $p$; so $\rho(P_K)=1$. Then $\rho(I_K)$ is a quotient of the procyclic $I_K/P_K$, and the image of the dense subgroup $\langle t\rangle$ is a
dense subgroup of a finite discrete group, hence everything. For $h=\rho(\sigma)$, $\rho(\Frob)h\rho(\Frob)^{-1}=\rho(\Frob\,\sigma\,\Frob^{-1})=\rho(\sigma^q)=h^q$ by
\Cref{thm:local:tame-relation}(ii) (the relation holds modulo $P_K$, which $\rho$ kills). Another lift is $\Frob\sigma'$ with $\sigma'\in I_K$, and $\rho(\sigma')$ commutes with
$h$, $\rho(I_K)$ being abelian. Finally $g^q=g$ iff $m\mid q-1$, and $g^q=g^{-1}$ iff $m\mid q+1$.
\end{proof}

\begin{proposition}[the image of wild inertia]\label{prop:local:wild-image}
Let $\ell\ne p$, $E_\lambda$ a finite extension of $\Q_\ell$ with ring of integers $\mathcal O$ and uniformiser $\varpi_\lambda$, and $\rho\colon G_K\to\GL_n(\mathcal O)$
continuous. Then $\rho(P_K)$ is a finite $p$-group, mapping injectively to $\GL_n(\mathcal O/\varpi_\lambda)$. Moreover:
\begin{enumerate}[label=(\roman*)]
\item $I_1=I_K\cap\rho^{-1}(1+\ell^2\Mat_n(\mathcal O))$ is an open subgroup of $I_K$, normal in $G_K$, and $\rho|_{I_1}$ factors through $t_\ell|_{I_1}$;
\item $\rho$ is tamely ramified (that is, $\rho(P_K)=1$) iff $\rho(I_K)$ contains no element of order $p$; in particular if $\rho(I_K)$ is finite of order
prime to $p$.
\end{enumerate}
\end{proposition}

\begin{proof}
$\Gamma_1=1+\varpi_\lambda\Mat_n(\mathcal O)$ is an open normal pro-$\ell$ subgroup of $\GL_n(\mathcal O)$: its quotients by $1+\varpi_\lambda^k\Mat_n(\mathcal O)$ are
successive extensions of copies of the additive group $\Mat_n(\mathcal O/\varpi_\lambda)$. The group $P'=P_K\cap\rho^{-1}(\Gamma_1)$ is open in $P_K$, hence pro-$p$,
and $\rho(P')$ is a compact subgroup of $\Gamma_1$, so pro-$\ell$, and a continuous image of a pro-$p$ group, so pro-$p$ (its finite quotients are
quotients of finite quotients of $P'$). Hence $\rho(P')=1$, $\rho(P_K)\cap\Gamma_1=1$, and $\rho(P_K)$ embeds in $\GL_n(\mathcal O/\varpi_\lambda)$; being a finite
continuous image of a pro-$p$ group, it is a $p$-group. (i) $1+\ell^2\Mat_n(\mathcal O)$ is open, normal and pro-$\ell$, so $I_1$ is open in $I_K$ and normal in
$G_K$, and $I_1=I_L$ for a finite $L/K$ (\Cref{cor:local:open-inertia}); $\rho|_{I_1}$ is a homomorphism to a pro-$\ell$ group, so it factors through $t_{\ell,K}|_{I_1}$ by
\Cref{cor:local:open-inertia}. (ii) If $\rho(P_K)\ne1$ it is a nontrivial $p$-group, which contains an element of order $p$. Conversely, if $\rho(P_K)=1$ then
$\rho(I_K)$ is a continuous image of $I_K/P_K$, a profinite abelian group all of whose finite quotients have order prime to $p$; so is $\rho(I_K)$, and an
element of order $p$ would generate a finite subgroup of order $p$, whose image in a suitable finite quotient of $\rho(I_K)$ is nontrivial.
\end{proof}

\begin{remark}\label{rem:local:wild-ec}
Nothing in \Cref{prop:local:wild-image} forbids wild ramification for a general $\rho$ at $p\ge5$: $\GL_2(\Z_\ell)$ contains elements of order $5$ when
$\ell\equiv1\pmod5$. For $\rho=\rho_{E,\ell}$ with potentially good reduction, the characteristic polynomials of the elements of $\rho(I_K)$ have integer
coefficients independent of $\ell$, so their orders lie in $\{1,2,3,4,6\}$, and the image is tame for $p\ge5$ (Serre and Tate \cite[Section 2]{SerreTate68};
\Cref{ch:reduction}). At $p=2,3$ the image can be wild: $Q_8$ or $\SL_2(\F_3)$ at $p=2$, cyclic of order $3$ or $6$ at $p=3$; an image of order $3$
is wild at $p=3$, because $I_K/P_K$ has no quotient of order $3$, and tame at $p=2$. These are the facts behind Remark rem:wild.
\end{remark}

\begin{example}[the tame relation at a finite level]\label{ex:local:tame-explicit}
Let $\ell\ne q$ be primes, $m=\ell^k$, and $L=\Q_q(\mu_m,q^{1/m})$. Fix a primitive $\zeta\in\mu_m$. An element $\sigma\in\Gal(L/\Q_q)$ is determined by $(a,b)$ with
$\sigma(\zeta)=\zeta^a$, $\sigma(q^{1/m})=\zeta^bq^{1/m}$, and $\sigma_{a,b}\sigma_{a',b'}=\sigma_{aa',ab'+b}$, because
$\sigma_{a,b}(\zeta^{b'}q^{1/m})=\zeta^{ab'}\zeta^bq^{1/m}$: the group is a subgroup of the affine group of $\Z/m$
(\Cref{ch:galois}). Since $\Q_q(\mu_m)$ is unramified with group generated by $\zeta\mapsto\zeta^q$, the values of $a$ are exactly the powers of $q$; and
$X^m-q$ is Eisenstein over $\Q_q(\mu_m)$, hence irreducible, so each automorphism of $\Q_q(\mu_m)$ has $m$ extensions to $L$, one for each value of $b$. So the group is
$\{(a,b):a\in\langle q\rangle\subseteq(\Z/m)^\times,\ b\in\Z/m\}$, of order $m\cdot\#\langle q\rangle$. The inertia group fixes the unramified field
$\Q_q(\mu_m)$, so it is contained in $\{a=1\}$, and it has order $e=m$ (\Cref{prop:local:ramgroups}(i)); so it is $\{a=1\}$, generated by $t=\sigma_{1,1}$, and
$t_\ell\bmod m$ is $\sigma_{1,b}\mapsto\zeta^b$ (as $\sigma_{1,b}^n=\sigma_{1,nb}$). The element $\Frob=\sigma_{q,0}$ is a Frobenius fixing $q^{1/m}$. Then
\[
\Frob\,t\,\Frob^{-1}=\sigma_{q,q}\sigma_{q',0}=\sigma_{1,q}=t^q,\qquad\Frob^{-1}t\,\Frob=\sigma_{q',q'}\sigma_{q,0}=\sigma_{1,q'}=t^{q'},
\]
where $qq'\equiv1\pmod m$; here $\Frob^{-1}=\sigma_{q',0}$, since $\sigma_{q,0}\sigma_{q',0}=\sigma_{1,0}$, and the first product is
$\Frob\,t=\sigma_{q,0}\sigma_{1,1}=\sigma_{q,q}$ followed by $\sigma_{q,q}\sigma_{q',0}=\sigma_{qq',q\cdot0+q}=\sigma_{1,q}$.
For $(q,\ell,k)=(7,3,2)$ these are $t^7$ and $t^4$; for $(13,3,2)$, $t^4$ and $t^7$; for $(2,3,2)$, $t^2$ and $t^5$ (all checked by composing the pairs).
\end{example}

\begin{paperlink}\label{pl:local:tamfactor-tame}
In the proof of Lemma lem:tamfactor: ``The inertia group $I=I_q$ acts through its maximal pro-$\ell$ tame quotient, generated by an element $t$'' is
\Cref{thm:local:tame-character}(iii), (iv), since $I_q$ acts on $T_\ell E$ through unipotent matrices over $\Z_\ell$ (a pro-$\ell$ group); ``the kernel of
$I\to\langle t\rangle$ has trivial pro-$\ell$ quotient'' is the statement that every finite quotient of $\ker t_\ell$ has order prime to $\ell$; and
``$\Frob^{-1}t\Frob=t^{q^{-1}}$'' is \Cref{thm:local:tame-relation}(ii) for the arithmetic Frobenius. In Lemmas lem:IV and lem:euler, ``the order $3$ is prime
to $q$, so the image is a quotient of the tame inertia group'' and ``Frobenius centralises $\sigma$ if $q\equiv1\pmod3$ and inverts it if $q\equiv2$'' are
\Cref{cor:local:finite-tame-image}, which needs only $q\nmid\#\rho(I_q)$, not $q\ge5$. In Proposition prop:shapes(ii), the extension of the
$G_{\Q_q}$-invariant quadratic character $\psi$ of $I_q$ to $G_{\Q_q}=I_q\rtimes\langle\phi\rangle$ is \Cref{cor:local:split}. In Lemma lem:tamadd, the fact that the maximal
pro-$\ell$ abelian quotient of the open subgroup $K\subseteq I_q$ is $\Z_\ell(1)$ is \Cref{cor:local:open-inertia}.
\end{paperlink}

\section{Kummer theory over local fields; quadratic characters}\label{sec:local:kummer}

The Tate curve makes the Kummer extension of a single element $q_E\in\Q_q^\times$ appear inside $T_\ell E$, and the paper reads off the action of
inertia and Frobenius on $T_\ell E$ from the Kummer cocycle of $q_E$. This section computes that cocycle on inertia ($v_K(x)$ times the tame
character) and on Frobenius (the class of the reduction of the unit part), with all the choices made explicit. It then describes the
quadratic characters of $G_{\Q_p}$ and the local behaviour of the global characters $\chi_d$, and states local class field theory in the form
needed for conductors. The paper uses this in Lemmas lem:tamagawa-exp and lem:tamfactor (Kummer classes), in Lemma lem:twist and Theorem
thm:twist (the twist by $\chi_{q^*}$ is unramified outside $q$ and tame at odd $q$), in Lemma lem:euler and Proposition prop:shapes ($V=V'\otimes\chi$
locally) and in Theorem thm:qual. The cohomological statements (Hilbert 90, $H^1(K,\Z_\ell(1))=\widehat{K^\times}$) are proved in \Cref{ch:galcoh};
here we prove the field-theoretic half.

\begin{proposition}[Kummer theory over $K$]\label{prop:local:kummer-qp}
Let $\ell\ne p$ and $x\in K^\times$.
\begin{enumerate}[label=(\roman*)]
\item (\Cref{ch:galois}, recalled.) For a compatible system of roots $(x^{1/\ell^k})_k$ in $\Qbar_p$, the \emph{Kummer cocycle}
$\kappa_x(\sigma)=\bigl(\sigma(x^{1/\ell^k})/x^{1/\ell^k}\bigr)_k\in\Z_\ell(1)$ is continuous and satisfies $\kappa_x(\sigma\tau)=\kappa_x(\sigma)+\chi_\ell(\sigma)\kappa_x(\tau)$. Changing the
system changes $\kappa_x$ by the coboundary $\sigma\mapsto(\chi_\ell(\sigma)-1)z$, $z\in\Z_\ell(1)$; $\kappa_{xy}=\kappa_x+\kappa_y$ for the product systems; and for any
extension $K'\supseteq K(\mu_{\ell^k})$, $\kappa_x\bmod\ell^k$ vanishes on $G_{K'}$ iff $x\in(K'^\times)^{\ell^k}$.
\item (Stated.) The Kummer maps $K^\times/(K^\times)^{\ell^k}\to H^1(K,\mu_{\ell^k})$ and $\widehat{K^\times}=\varprojlim_kK^\times/(K^\times)^{\ell^k}\to H^1(K,\Z_\ell(1))$ are
isomorphisms. We do not prove this here; it is proved in \Cref{ch:galcoh} from Hilbert's Theorem 90 (see also \cite[Chapter XIV]{SerreLocalFields}).
\item With \Cref{thm:local:powers}: $H^1(K,\Z_\ell(1))\simeq\Z_\ell\times\Z/\ell^{v}$, $v=v_\ell(q-1)$; the projection to $\Z_\ell$ is induced by $v_K$, the torsion is the
image of the units, $\F_q^\times\otimes\Z_\ell\simeq\Z_\ell(1)/(q-1)\Z_\ell(1)$; and $\#H^1(K,\mu_{\ell^k})=\ell^k\cdot\#\mu_{\ell^k}(K)$.
\end{enumerate}
\end{proposition}

\begin{proof}
(i) Write $\alpha_k=x^{1/\ell^k}$. Then $\sigma(\alpha_k)/\alpha_k\in\mu_{\ell^k}$, compatible in $k$, and
$\sigma\tau(\alpha_k)/\alpha_k=\sigma\bigl(\tau(\alpha_k)/\alpha_k\bigr)\cdot\sigma(\alpha_k)/\alpha_k$, which is the cocycle identity written multiplicatively. Each component factors through
the finite extension $K(\mu_{\ell^k},\alpha_k)$, so $\kappa_x$ is continuous. Another system is $(\zeta_k\alpha_k)$ with $z=(\zeta_k)\in\Z_\ell(1)$, which adds
$\sigma(\zeta_k)/\zeta_k$, that is $(\chi_\ell(\sigma)-1)z$. The product rule is clear. For $\sigma\in G_{K'}$, $\sigma$ fixes $\mu_{\ell^k}$, and $\kappa_x(\sigma)\equiv0\pmod{\ell^k}$ means $\sigma(\alpha_k)=\alpha_k$; this holds for all such
$\sigma$ iff $\alpha_k\in K'$; and if $x=y^{\ell^k}$ with $y\in K'$, then $\alpha_k=\zeta y\in K'$ for some $\zeta\in\mu_{\ell^k}$. (iii) combines (ii) with
\Cref{thm:local:powers}(iv), (v).
\end{proof}

\begin{theorem}[the Kummer cocycle on inertia and on Frobenius]\label{thm:local:kummer-inertia}
Let $\ell\ne p$, $\varpi$ a uniformiser of $K$, $x\in K^\times$, $a=v_K(x)$ and $x=\varpi^au$ with $u\in\mathcal O_K^\times$. Choose compatible systems of $\ell$-power roots of
$\varpi$ and $u$, and give $x$ the product system. Put $v=v_\ell(q-1)$.
\begin{enumerate}[label=(\roman*)]
\item $\kappa_x=a\,\kappa_\varpi+\kappa_u$.
\item $\kappa_\varpi|_{I_K}=t_\ell$ and $\kappa_u|_{I_K}=0$; hence
\[
\kappa_x|_{I_K}=v_K(x)\cdot t_\ell ,
\]
independently of all choices.
\item $\kappa_u$ is unramified: it factors through $G_K/I_K$, so it is determined by $\kappa_u(\Frob)$, which is well defined modulo $(q-1)\Z_\ell(1)$. Under
$\Z_\ell(1)/(q-1)\Z_\ell(1)=\mu_{\ell^v}$ its class is the Teichm\"uller lift of $\bar u^{(q-1)/\ell^v}$; that is, it is the class of $\bar u$ in
$\F_q^\times/(\F_q^\times)^{\ell^v}=\F_q^\times\otimes\Z_\ell$. In particular $\kappa_u$ is a coboundary iff $\bar u\in(\F_q^\times)^{\ell^v}$.
\item If $\Frob$ fixes the chosen roots of $\varpi$ (\Cref{cor:local:frob-fixing-roots}), then $\kappa_\varpi(\Frob)=0$ and $\kappa_x(\Frob)=\kappa_u(\Frob)$. For any other lift
$\Frob\sigma$ ($\sigma\in I_K$), $\kappa_x(\Frob\sigma)=\kappa_x(\Frob)+q\,a\,t_\ell(\sigma)$. Hence, after identifying $\Z_\ell(1)\simeq\Z_\ell$, the value $c'=\kappa_x(\Frob)$ is well
defined modulo the ideal $(a,q-1)\Z_\ell=\ell^{\min(v_\ell(a),v)}\Z_\ell$ when the lift and the roots vary, and $\min(v_\ell(a),v,v_\ell(c'))$ does not depend on any choice.
\item For $k\ge1$: $I_K$ fixes the $\ell^k$-th roots of $x$ iff $\ell^k\mid a$; the image of $I_K$ in $\Gal(K^{\ur}(x^{1/\ell^k})/K^{\ur})\subseteq\mu_{\ell^k}$ has order
$\ell^{k-\min(k,v_\ell(a))}$. Equivalently, $x$ is an $\ell^k$-th power in $K^{\ur}$ iff $\ell^k\mid v_K(x)$.
\end{enumerate}
\end{theorem}

\begin{proof}
(i) is the product rule. (ii) On $I_K$, $\kappa_\varpi$ is by definition $t_\ell$ (\Cref{thm:local:tame-character}(iii)). The $\ell^k$-th roots of $u$ lie in
$K^{\ur}$ (\Cref{prop:local:qpur}(iv), with $\mu_{\ell^k}\subset K^{\ur}$), so $I_K$ fixes them and $\kappa_u|_{I_K}=0$. A cocycle restricted to $I_K$, on which
$\chi_\ell=1$, is a homomorphism, and coboundaries $(\chi_\ell-1)z$ vanish on $I_K$; so the restriction does not depend on the choices.
(iii) For $g\in G_K$ and $\sigma\in I_K$, $\kappa_u(g\sigma)=\kappa_u(g)+\chi_\ell(g)\kappa_u(\sigma)=\kappa_u(g)$, so $\kappa_u$ factors through $G_K/I_K$; there $\Frob$ is a topological
generator, and the coboundaries are $(\chi_\ell(\Frob)-1)z=(q-1)z$. To compute the class, write $u=[\bar u]u_1$ with $u_1\in U^{(1)}$ (\Cref{thm:local:teichmuller}).
Changing roots changes $\kappa_u$ by a coboundary, so we may take the unique roots of $u_1$ in $U^{(1)}$ (\Cref{thm:local:powers}(i)), for which
$\kappa_{u_1}=0$, and roots $\omega_k=[\bar u]^{1/\ell^k}$, which are roots of unity of order prime to $p$. Since $\Frob$ acts on such roots of unity by
$y\mapsto y^q$, $\kappa_u(\Frob)=(\omega_k^{q-1})_k$. The projection $\Z_\ell(1)\to\Z_\ell(1)/(q-1)\Z_\ell(1)=\Z_\ell(1)/\ell^v\Z_\ell(1)=\mu_{\ell^v}$ is the component at
$k=v$, which is $\omega_v^{q-1}=\omega_v^{\ell^vN}=[\bar u]^N$ with $N=(q-1)/\ell^v$. Its reduction is $\bar u^N$. Finally $y\mapsto y^N$ maps the cyclic group
$\F_q^\times$ onto $\mu_{\ell^v}(\F_q)$ with kernel $(\F_q^\times)^{\ell^v}$ (\Cref{thm:local:powers}(v)).
(iv) If $\Frob$ fixes the roots of $\varpi$, $\kappa_\varpi(\Frob)=0$ and (i) gives $\kappa_x(\Frob)=\kappa_u(\Frob)$. The cocycle identity and (ii) give
$\kappa_x(\Frob\sigma)=\kappa_x(\Frob)+\chi_\ell(\Frob)\,a\,t_\ell(\sigma)=\kappa_x(\Frob)+q\,a\,t_\ell(\sigma)$, and $t_\ell$ is surjective; changing roots adds $(q-1)z$. As $q$ is a unit in
$\Z_\ell$, the ambiguity is the ideal $(a,q-1)$, generated by $\ell^{\min(v_\ell(a),v)}$, and changing $c'$ by an element of it does not change
$\min(v_\ell(a),v,v_\ell(c'))$. Changing the identification $\Z_\ell(1)\simeq\Z_\ell$ multiplies $c'$ by a unit.
(v) By (ii), $\sigma\in I_K$ moves $x^{1/\ell^k}$ by $\zeta^{a\,t(\sigma)}$, with $\zeta$ a generator of $\mu_{\ell^k}$ and $t(\sigma)\in\Z/\ell^k$ the image of $t_\ell(\sigma)$; as
$t_\ell$ is surjective, the image of $I_K$ is the subgroup of $\mu_{\ell^k}$ generated by $\zeta^a$, of order $\ell^{k-\min(k,v_\ell(a))}$, which is trivial iff
$\ell^k\mid a$. The last statement is also \Cref{prop:local:qpur}(iv).
\end{proof}

A small case shows (iii) at work. Take $K=\Q_7$, $\ell=3$, so $v=v_3(6)=1$ and $N=(q-1)/\ell^v=2$, and $x=u=2$. Then $\bar u^N=4$, a primitive cube root
of unity in $\F_7$ ($4^3=64\equiv1$), so the class of $\kappa_2$ in $\Z_3(1)/(q-1)\Z_3(1)\simeq\mu_3$ is nontrivial: $\kappa_2$ is not a coboundary, and $2$ is not
a cube in $\Q_7$, in agreement with \Cref{ex:local:powers-examples}(a). For $x=u=6$ instead, $\bar u^N=36\equiv1$, and $\kappa_6$ is a coboundary, matching
the fact that $6\equiv-1=(-1)^3$ is a cube.

\begin{paperlink}\label{pl:local:kummer}
\emph{Lemma lem:tamagawa-exp.} Over $\Q_q^{\ur}$ (or already over the unramified quadratic extension, where the curve is a Tate curve) the
$\ell^t$-torsion is generated by $\mu_{\ell^t}$ and an $\ell^t$-th root of $q_E$, and inertia acts trivially on $E[\ell^t]$ iff it fixes that root, iff
$\ell^t\mid v_q(q_E)=n_q$ by \Cref{thm:local:kummer-inertia}(v). This is the precise form of the paper's statement that $q_E$ is an $\ell^t$-th power in $\Q_q^{\ur}$ if and only if
$\ell^t\mid n_q$. \emph{Lemma lem:tamfactor.} ``The restriction of $\kappa$ to $I$ is $n$ times the tame character'' is
\Cref{thm:local:kummer-inertia}(ii) with $x=q_E$, $a=n$. ``$\kappa=n\kappa_q+\kappa_u$, where $\kappa_u$ is unramified, inflated from the class of $u$'' is (i) and
(iii) with $\varpi=q$. ``$c'=\varepsilon\kappa(\Frob)$ is a cocycle value: it is $\varepsilon$ times the class of $u$ for the Frobenius lift fixing a compatible system of
$\ell$-power roots of $q$, and in general it is well defined only modulo $\ell^{\min(v_\ell(n),v_\ell(q-1))}$'' is (iii) and (iv), together with
\Cref{cor:local:frob-fixing-roots}; the factor $\varepsilon$ is the twist of \Cref{ex:local:tate-shape-wd}. The identification of the class in
$H^1(G/I,\Z_\ell(1))$ with the class of $\bar u$ is (iii); its cohomological meaning is \Cref{ch:galcoh}.
\end{paperlink}

\subsection{Quadratic extensions and quadratic characters}

\begin{proposition}[quadratic extensions of $\Q_p$]\label{prop:local:quadratic-ext}
\begin{enumerate}[label=(\roman*)]
\item Let $p$ be odd and $u_0\in\Z_p^\times$ a non-square. $\Q_p$ has exactly three quadratic extensions: the unramified $\Q_p(\sqrt{u_0})=\Q_p(\mu_{p^2-1})$ and
the ramified $\Q_p(\sqrt p)$, $\Q_p(\sqrt{u_0p})$, both tamely ramified. The field $\Q_p(\sqrt{p^*})$ is $\Q_p(\sqrt p)$ if $p\equiv1\pmod4$ and $\Q_p(\sqrt{-p})=\Q_p(\sqrt{u_0p})$
if $p\equiv3\pmod4$.
\item Let $p$ be odd. The quadratic characters of $G_{\Q_p}$ are $1$, the unramified $\eta_{\ur}$ ($\eta_{\ur}(\Frob)=-1$) and two ramified ones $\chi$, $\chi\eta_{\ur}$. Every
ramified one restricts on $I_p$ to the unique character of order $2$ of $I_p$, which is $\theta_2\colon\sigma\mapsto\sigma(\sqrt p)/\sqrt p$.
\item $\Q_2$ has seven quadratic extensions, $\Q_2(\sqrt d)$ for $d\in\{-1,\pm5,\pm2,\pm10\}$. The unramified one is $\Q_2(\sqrt5)=\Q_2(\sqrt{-3})$; the others are wildly
ramified, $\Q_2(\sqrt{-1})$ and $\Q_2(\sqrt{-5})=\Q_2(\sqrt3)$ with lower ramification jump $1$, and $\Q_2(\sqrt{\pm2})$, $\Q_2(\sqrt{\pm10})$ with jump $2$
(\Cref{ex:local:ramgroups}).
\end{enumerate}
\end{proposition}

\begin{proof}
By Kummer theory for $n=2$ (\Cref{ch:galois}), quadratic extensions of $\Q_p$ correspond to the nontrivial elements of $\Q_p^\times/(\Q_p^\times)^2$, through
$d\mapsto\Q_p(\sqrt d)$, and quadratic characters to the same group (the character of $\Q_p(\sqrt d)$ being $\sigma\mapsto\sigma(\sqrt d)/\sqrt d$). (i) By
\Cref{thm:local:powers}(vi) the nontrivial classes are $u_0$, $p$, $u_0p$. The residue field $\F_{p^2}$ of $\Q_p(\mu_{p^2-1})$ (\Cref{thm:local:unramified}) contains a
square root of $\bar u_0$, so $u_0$ is a square there (\Cref{thm:local:powers}(ii)), and $\Q_p(\sqrt{u_0})=\Q_p(\mu_{p^2-1})$ is unramified. $X^2-p$ and $X^2-u_0p$ are
Eisenstein. If $p\equiv3\pmod4$, $-1$ is not a square modulo $p$, so $-p\in u_0p(\Q_p^\times)^2$. Ramification index $2$ is prime to $p$. (ii) A quadratic
character is unramified iff its field is, and $\eta_{\ur}$ is the character of the unramified extension, on whose group $\Frob$ acts nontrivially.
A ramified quadratic character is nontrivial on $I_p$. A character of order $2$ of $I_p$ is continuous with finite image, so it kills the pro-$p$ group
$P_p$ ($p$ odd), and factors through the procyclic $I_p/P_p\simeq\prod_{\ell\ne p}\Z_\ell(1)$, which has exactly one quotient of order $2$; and
$\theta_2$ (\Cref{thm:local:tame-character}, $m=2$, $\varpi=p$) is such a character. (iii) $\Q_2^\times/(\Q_2^\times)^2$ has order $8$ with representatives $\pm1,\pm5,\pm2,\pm10$
(\Cref{thm:local:powers}(vi)). $\alpha=(1+\sqrt5)/2$ is a root of $X^2-X-1\in\Z_2[X]$, so it lies in the ring of integers of $L=\Q_2(\sqrt5)$, and its reduction is a root of $X^2+X+1$,
which is irreducible over $\F_2$. So the residue field of $L$ contains $\F_4$, $f(L/\Q_2)=2=[L:\Q_2]$, $e(L/\Q_2)=1$ and $L$ is unramified (\Cref{thm:local:ef}); and
$-3\cdot5=-15\equiv1\pmod8$ and $3\cdot(-5)\equiv1\pmod 8$ are squares, so $\Q_2(\sqrt{-3})=\Q_2(\sqrt5)$ and $\Q_2(\sqrt3)=\Q_2(\sqrt{-5})$. There is only one unramified quadratic
extension, so the other six are ramified, of degree $2$, hence wild; the jumps are \Cref{ex:local:ramgroups}.
\end{proof}

\begin{proposition}[global quadratic characters locally]\label{prop:local:global-quadratic-local}
Let $d$ be a fundamental discriminant and $\chi_d$ the character of $\Q(\sqrt d)$ (\Cref{ch:galois}). Through \Cref{prop:local:decomposition}, $\chi_d|_{G_{\Q_p}}$ is the
character of $\Q_p(\sqrt d)$.
\begin{enumerate}[label=(\roman*)]
\item $\chi_d$ is unramified at $p$ iff $p\nmid d$, and then $\chi_d(\Frob_p)=\bigl(\frac dp\bigr)$ (\Cref{ch:ant}).
\item For an odd prime $q$: at every $p\ne q$, including $p=2$, $\chi_{q^*}$ is unramified, with $\chi_{q^*}(\Frob_p)=\bigl(\frac pq\bigr)$; at $q$ it is ramified and tame, and
$\chi_{q^*}|_{I_q}=\theta_2$.
\item $\chi_{-3}$ at $3$ is tamely ramified; $\chi_{-4}$ at $2$ and $\chi_{8}$, $\chi_{-8}$ at $2$ are wildly ramified, with ramification jumps $1$ and $2$.
\item For a finite set $S$ of odd primes, $\chi_S=\prod_{q\in S}\chi_{q^*}$ is unramified outside $S$ and restricts to $\theta_2$ on $I_q$ for each $q\in S$.
\item For any representation $\rho$ of $G_\Q$ and odd $q$: $(\rho\otimes\chi_{q^*})|_{I_p}=\rho|_{I_p}$ for $p\ne q$, and $(\rho\otimes\chi_{q^*})|_{I_q}=\rho|_{I_q}\otimes\theta_2$. So
twisting by $\chi_{q^*}$ changes nothing on inertia at $p\ne q$ (the image of inertia, $V^{I_p}$, the conductor exponent), multiplies Frobenius
there by the sign $\chi_{q^*}(\Frob_p)$, and at $q$ multiplies the action of inertia by the character $\theta_2$.
\item Every quadratic character of $G_{\Q_p}$ is the restriction of a product of characters $\chi_d$: for $p$ odd, of $1$, $\chi_{d_0}$, $\chi_{p^*}$ or
$\chi_{d_0}\chi_{p^*}$, where $d_0$ is a fundamental discriminant prime to $p$ with $\bigl(\frac{d_0}p\bigr)=-1$; for $p=2$, of a product of $\chi_{-4}$, $\chi_8$ and
$\chi_5$.
\end{enumerate}
\end{proposition}

\begin{proof}
The local statement follows from the definitions: $\sigma(\sqrt d)=\chi_d(\sigma)\sqrt d$ for $\sigma\in G_\Q$, and the embedding $\iota$ of \Cref{prop:local:decomposition}
transports this to $\Q_p(\sqrt d)$. (i) For $p$ odd, $d$ is $p$ times a unit or a unit, according as $p\mid d$ or not; in the first case $X^2-d$ is
Eisenstein, in the second $\Q_p(\sqrt d)$ is trivial or unramified (\Cref{prop:local:quadratic-ext}(i)). For $p=2$, a fundamental discriminant is $\equiv1\pmod4$
(odd) or divisible by $4$; if $d\equiv1\pmod8$ the field is trivial, if $d\equiv5\pmod8$ it is $\Q_2(\sqrt5)$, unramified, and if $4\mid d$ then
$d/4\equiv2,3\pmod4$ and $\Q_2(\sqrt d)=\Q_2(\sqrt{d/4})$ is one of the six ramified fields (\Cref{prop:local:quadratic-ext}(iii)). The value at $\Frob_p$ is
\Cref{ch:ant}. (ii) $q^*\equiv1\pmod4$, so (i) applies; $(q^*/p)=(p/q)$ is quadratic reciprocity (\Cref{ch:ant}). At $q$, $\chi_{q^*}$ is ramified and $q$ is odd,
so \Cref{prop:local:quadratic-ext}(ii) gives $\theta_2$. (iii) $\Q_3(\sqrt{-3})$ is tame (\Cref{prop:local:z3zeta3}(ii)); $\Q_2(\sqrt{-1})$, $\Q_2(\sqrt{\pm2})$ are
\Cref{ex:local:ramgroups}. (iv) and (v) follow from (ii), since a character is trivial on $I_p$ where it is unramified. (vi) By \Cref{prop:local:quadratic-ext}
the classes in $\Q_p^\times/(\Q_p^\times)^2$ are represented by $1,d_0,p^*,d_0p^*$ for $p$ odd ($d_0$ is a non-square unit; such a $d_0$ exists: $-4$ if
$p\equiv3\pmod4$, $8$ if $p\equiv5\pmod8$, and in general a prime $d_0\equiv1\pmod4$ that is a quadratic non-residue modulo $p$, which exists by Dirichlet's
theorem, a case of Chebotarev's, \Cref{ch:ant}), and are generated by $-1$, $2$, $5$ for $p=2$;
the characters of $\Q_2(\sqrt{-1})$, $\Q_2(\sqrt2)$, $\Q_2(\sqrt5)$ are the restrictions of $\chi_{-4}$, $\chi_8$, $\chi_5$.
\end{proof}

\begin{theorem}[local class field theory for $\Q_p$]\label{thm:local:lcft}
\begin{enumerate}[label=(\roman*)]
\item (Local Kronecker-Weber.) The maximal abelian extension of $\Q_p$ is $\Q_p^{\mathrm{ab}}=\Q_p^{\ur}\cdot\Q_p(\mu_{p^\infty})$.
\item There is a unique continuous homomorphism $\Art\colon\Q_p^\times\to\Gal(\Q_p^{\mathrm{ab}}/\Q_p)$ such that $\Art(p)$ is the arithmetic Frobenius on $\Q_p^{\ur}$ and the
identity on $\Q_p(\mu_{p^\infty})$, and, for $u\in\Z_p^\times$, $\Art(u)$ is the identity on $\Q_p^{\ur}$ and acts on $\mu_{p^\infty}$ by $\zeta\mapsto\zeta^{u^{-1}}$. It is
injective with dense image, $\Art(\Z_p^\times)$ is the image of $I_p$, and $\Art(1+p^n\Z_p)$ is the image of the upper ramification group $G_{\Q_p}^n$ ($n\ge1$).
\item For a character $\chi$ of $G_{\Q_p}$ of finite order, its Artin conductor exponent $a(\chi)$ (\Cref{def:local:swan}) is the least $n\ge0$ with
$\chi(\Art(U^{(n)}))=1$, where $U^{(0)}=\Z_p^\times$ and $U^{(n)}=1+p^n\Z_p$. For a Dirichlet character $\chi$ (\Cref{ch:ant}) the product of the
$p^{a(\chi|_{G_{\Q_p}})}$ is the conductor of $\chi$.
\end{enumerate}
We do not prove this; see Serre's chapter ``Local class field theory'' in Cassels and Fr\"ohlich \cite{CasselsFrohlich}, Milne \cite[Chapter I]{MilneCFT},
and Neukirch \cite[Chapter V]{Neukirch}.
\end{theorem}

\begin{warning}[normalisations of the Artin map]\label{warn:local:artin}
The book's Artin map sends a uniformiser to an \emph{arithmetic} Frobenius, as in \Cref{ch:ant}; units then act on $\mu_{p^\infty}$ through $u^{-1}$.
Deligne \cite{Deligne73} and many sources on Weil-Deligne representations send uniformisers to geometric Frobenius elements, and units act through
$u$. The book never needs the other normalisation, and uses local class field theory only for the conductors of characters,
\Cref{thm:local:lcft}(iii). For example, $\chi_{-4}$ is trivial on $\Art(1+4\Z_2)$ but not on $\Art(\Z_2^\times)$, so $a(\chi_{-4}|_{G_{\Q_2}})=2$; $\chi_8$ is trivial on
$\Art(1+8\Z_2)$ but not on $\Art(1+4\Z_2)$, so $a=3$; these agree with \Cref{thm:local:swan}(iv) and with the conductors $4$ and $8$.
\end{warning}

\section{The Weil group and Weil-Deligne representations}\label{sec:local:wd}

An $\ell$-adic representation of $G_K$ ($\ell\ne p$) can have infinite image of inertia, as the Tate module of a curve with multiplicative reduction
does; Grothendieck's theorem says that this infinite part is always unipotent, given by a nilpotent ``monodromy'' operator $N$. Recording the
finite part $r$ and $N$ separately gives the Weil-Deligne representation, in terms of which the local invariants ($V^{I}$, $L$-factors, conductors)
are computed in \Cref{sec:local:L,sec:local:shapes}, independently of $\ell$. The paper speaks of ``local components'' and the ``unramified twist of the
Steinberg representation'' in Lemmas lem:twist(iii) and lem:euler; on the Galois side these are the Weil-Deligne representations
$\Sp(2)\otimes\mu$ of this section (the automorphic side is \Cref{ch:modgalrep}).

\begin{convention}[the norm on the Weil group]\label{conv:local:weil}
For $w$ in the Weil group $W_K$ (\Cref{def:local:weil-group}), $d(w)\in\Z$ is defined by $w|_{K^{\ur}}=\Frob^{d(w)}$, and $\norm w=q^{d(w)}$. So $\norm{\Frob}=q$,
$\norm{\Phi}=q^{-1}$, $\norm{\cdot}$ is trivial on $I_K$, and $\norm w=\chi_\ell(w)$ for every $\ell\ne p$. A Weil-Deligne representation satisfies
\[
r(w)\,N\,r(w)^{-1}=\norm w\,N,\qquad\text{that is}\qquad r(\Frob)\,N\,r(\Frob)^{-1}=q\,N,\quad r(\Phi)\,N\,r(\Phi)^{-1}=q^{-1}N .
\]
Some sources write $\Frob\,N\,\Frob^{-1}=q^{-1}N$ with $\Frob$ geometric; here $\Frob$ is arithmetic (Appendix B, convention B-C2(f)). This $\norm{\cdot}$ agrees with
Tate's normalisation in \cite{Tate79}.
\end{convention}

\begin{definition}[the Weil group]\label{def:local:weil-group}
The \emph{Weil group} of $K$ is $W_K=\{w\in G_K:\ w|_{K^{\ur}}\in\Frob^{\Z}\}$, the preimage of the subgroup $\Z\subset\hat\Z\simeq\Gal(K^{\ur}/K)$. It is given the
topology in which $I_K$, with its profinite topology, is an open subgroup.
\end{definition}

\begin{proposition}\label{prop:local:weil-group}
\begin{enumerate}[label=(\roman*)]
\item $W_K$ is a dense subgroup of $G_K$ containing $I_K$, with $W_K/I_K\simeq\Z$ generated by the image of $\Frob$; its topology is finer than the
topology induced from $G_K$, and every $w\in W_K$ is uniquely $w=\Phi^n\sigma$ with $n\in\Z$, $\sigma\in I_K$, for a fixed geometric Frobenius lift $\Phi\in W_K$.
\item $d\colon W_K\to\Z$ is a continuous homomorphism, $\norm{\Frob}=q$, $\norm{\Phi}=q^{-1}$, and $\norm w=\chi_\ell(w)$ for $\ell\ne p$.
\item A homomorphism $\mu\colon W_K\to\Omega^\times$ trivial on $I_K$ (an \emph{unramified character}) is determined by $\mu(\Phi)$, which may be any element of
$\Omega^\times$; $\norm{\cdot}$ is the unramified character with $\Phi\mapsto q^{-1}$. A character of $G_K$ with values in $\Q_\ell^\times$ takes values in $\Z_\ell^\times$, but a
character of $W_K$ need not.
\item (Stated.) The Artin map of \Cref{thm:local:lcft} identifies $\Q_p^\times$ with $W_{\Q_p}^{\mathrm{ab}}$, and $\norm{\Art(x)}=\abs x_p^{-1}$.
\end{enumerate}
\end{proposition}

\begin{proof}
(i) The restriction $G_K\to\hat\Z$ is open and surjective, and $\Z$ is dense in $\hat\Z$, so the preimage $W_K$ is dense; $I_K$ is the kernel. Open subsets of
$I_K$ are open in $W_K$ by definition, and a subset of $W_K$ open in the induced topology is a union of cosets of open subsets of $I_K$; so the
topology is finer. A lift $\Phi$ maps to $-1\in\Z$, a generator. (ii) $d$ is the composite $W_K\to W_K/I_K\simeq\Z$, continuous because $I_K$ is open; $\chi_\ell$ is
unramified with $\chi_\ell(\Frob)=q$ (\Cref{prop:local:cyclotomic-local}). (iii) $W_K/I_K\simeq\Z$ is free on the image of $\Phi$. A continuous character of the
compact group $G_K$ has compact image in $\Q_\ell^\times$; its image under $v_\ell$ is a compact subgroup of $\Z$, hence $0$, so the image lies in $\Z_\ell^\times$.
The unramified character of $W_K$ with $\Phi\mapsto\ell$ does not.
\end{proof}

\begin{lemma}[matrix logarithm]\label{lem:local:matrix-log}
Let $\mathcal O$ be the ring of integers of a finite extension $E_\lambda$ of $\Q_\ell$ and $\mathfrak g=\ell^2\Mat_n(\mathcal O)$. Then $\exp(X)=\sum X^k/k!$ converges for
$X\in\mathfrak g$ and $\log(g)=\sum_{k\ge1}(-1)^{k+1}(g-1)^k/k$ for $g\in1+\mathfrak g$; they are mutually inverse bijections $\mathfrak g\leftrightarrows1+\mathfrak g$;
$\exp(X+Y)=\exp X\exp Y$ when $XY=YX$; and $g^a=\exp(a\log g)$ for $a\in\Z_\ell$, where $g^a=\lim g^{a_k}$ for integers $a_k\to a$. For nilpotent $N$ the
series $\exp(tN)$ is a polynomial in $t$ and $\log\exp(tN)=tN$ for all $t\in E_\lambda$.
\end{lemma}

\begin{proof}
Use $\norm X=\max\abs{x_{ij}}_\ell$, which satisfies $\norm{XY}\le\norm X\norm Y$ (ultrametric inequality), and the estimates in the proof of
\Cref{thm:local:units-structure} with $c=2$: the terms are of valuation $\ge2k-(k-1)=k+1\ge2$, so both series converge into $\mathfrak g$, resp.
$1+\mathfrak g$. The identities hold in $\Q[[X,Y]]$ (for commuting variables) and transfer, entry by entry, by \Cref{lem:local:summation} with the
same bounds, all the matrices involved commuting. For $a\in\Z$, $\exp(a\log g)=\exp(\log g)^a=g^a$; both sides are continuous in $a$ (the left side by
the bounds, the right side because $g^{\ell^k}\to1$ in the pro-$\ell$ group $1+\mathfrak g$), so they agree on $\Z_\ell$. For nilpotent $N$ the series are
finite and the identity is a polynomial identity.
\end{proof}

\begin{theorem}[Grothendieck's $\ell$-adic monodromy theorem]\label{thm:local:grothendieck}
Let $\ell\ne p$, $E_\lambda$ a finite extension of $\Q_\ell$, $V$ a finite-dimensional $E_\lambda$-vector space and $\rho\colon G_K\to\GL(V)$ continuous. Fix a generator
$\zeta$ of $\Z_\ell(1)$ and let $t_\zeta\colon I_K\to\Z_\ell$ be the tame character followed by $\zeta\mapsto1$. There are an open subgroup $I_1\subseteq I_K$, normal in $G_K$,
and a unique nilpotent $N\in\End(V)$ with
\[
\rho(\sigma)=\exp\bigl(t_\zeta(\sigma)N\bigr)\qquad\text{for all }\sigma\in I_1,
\]
and $\rho(w)\,N\,\rho(w)^{-1}=\norm w\,N$ for all $w\in W_K$; in particular $\rho(\Frob)N\rho(\Frob)^{-1}=qN$. Moreover $N=0$ iff $\rho(I_K)$ is finite, and replacing
$\zeta$ by $\zeta^a$ ($a\in\Z_\ell^\times$) replaces $N$ by $aN$.
\end{theorem}

\begin{proof}
$\rho(G_K)$ is compact, so it stabilises a lattice $T$ (the stabiliser of any lattice $T_0$ is open, and the sum of the finitely many translates
$\rho(g)T_0$ is stable); in a basis of $T$, $\rho\colon G_K\to\GL_n(\mathcal O)$. Let $I_1=I_K\cap\rho^{-1}(1+\ell^2\Mat_n(\mathcal O))$, open in $I_K$ and normal in $G_K$.
By \Cref{prop:local:wild-image}(i), $\rho|_{I_1}=\phi\circ t_\zeta|_{I_1}$ for a homomorphism $\phi\colon t_\zeta(I_1)\to1+\ell^2\Mat_n(\mathcal O)$, continuous because
$t_\zeta|_{I_1}$ is a quotient map of compact groups. The image $t_\zeta(I_1)$ is a closed subgroup of finite index of $\Z_\ell$, so $t_\zeta(I_1)=\ell^m\Z_\ell$. Choose
$\sigma_1\in I_1$ with $t_\zeta(\sigma_1)=\ell^m$, put $g=\rho(\sigma_1)$ and $N=\ell^{-m}\log g$ (\Cref{lem:local:matrix-log}). For $\sigma\in I_1$ with $t_\zeta(\sigma)=\ell^mb$, $b\in\Z_\ell$,
continuity of $\phi$ gives $\rho(\sigma)=\phi(\ell^m)^b=g^b=\exp(b\log g)=\exp(t_\zeta(\sigma)N)$. For $w\in W_K$, $w\sigma_1w^{-1}\in I_1$ and $t_\zeta(w\sigma_1w^{-1})=\chi_\ell(w)\ell^m=\norm w\ell^m$
(\Cref{thm:local:tame-relation}(i)); hence
\[
\exp\bigl(\rho(w)\log g\,\rho(w)^{-1}\bigr)=\rho(w)g\rho(w)^{-1}=\rho(w\sigma_1w^{-1})=\exp(\norm w\log g).
\]
Both arguments of $\exp$ lie in $\ell^2\Mat_n(\mathcal O)$ ($\rho(w)\in\GL_n(\mathcal O)$, and $\norm w\in\Z_\ell^\times$), where $\exp$ is injective; so $\rho(w)N\rho(w)^{-1}=\norm wN$. In
particular $N$ is conjugate to $qN$, so the finite set of eigenvalues of $N$ (in an algebraic closure) is stable under multiplication by $q$; as
$q^i\ne q^j$ for $i\ne j$, all eigenvalues are $0$ and $N$ is nilpotent. Uniqueness: if $N'$ is nilpotent with $\rho=\exp(t_\zeta N')$ on an open $I_1'$, pick
$\sigma\in I_1\cap I_1'$ with $t_\zeta(\sigma)=\ell^{m'}\ne0$; then $\ell^{m'}N=\log\rho(\sigma)=\ell^{m'}N'$ by \Cref{lem:local:matrix-log}. If $N=0$ then $\rho(I_1)=1$ and
$\rho(I_K)$ is finite. If $\rho(I_K)$ is finite, $g$ has finite order and $\log(g^k)=k\log g$, so $\log g=0$ and $N=0$. Finally $t_{\zeta^a}=a^{-1}t_\zeta$, so
$\exp(t_\zeta(\sigma)N)=\exp(t_{\zeta^a}(\sigma)\,aN)$.
\end{proof}

The theorem says that $N$ is the logarithm of the monodromy: just as a unipotent matrix $u$ is $\exp(\log u)$ with $\log u$ nilpotent, the action of
an open subgroup of inertia is $\exp$ of a multiple of $N$. The same statement holds for continuous representations of $W_K$ (Deligne
\cite[Section 8]{Deligne73}); we need only $G_K$.

\begin{definition}[Weil-Deligne representations]\label{def:local:wd}
Let $\Omega$ be a field of characteristic $0$. A \emph{Weil-Deligne representation} of $W_K$ over $\Omega$ is a pair $(r,N)$ on a finite-dimensional
$\Omega$-vector space $V$: a homomorphism $r\colon W_K\to\GL(V)$ whose restriction to $I_K$ has open kernel (so $r(I_K)$ is finite), and a nilpotent $N\in\End(V)$
with $r(w)Nr(w)^{-1}=\norm wN$ for all $w\in W_K$. A morphism is a linear map commuting with $r$ and $N$. $(r,N)$ is \emph{Frobenius-semisimple} if $r(\Phi)$ is
semisimple.
\end{definition}

\begin{proposition}[operations on Weil-Deligne representations]\label{prop:local:wd-constructions}
If $(r,N,V)$ and $(r',N',V')$ are Weil-Deligne representations, so are: $V\oplus V'$ with $N\oplus N'$; $V\otimes V'$ with $N\otimes1+1\otimes N'$; the dual $V^\vee$ with
$r^\vee(w)=r(w^{-1})^{\mathsf T}$ and $-N^{\mathsf T}$; $\Hom(V,V')$ with $f\mapsto N'f-fN$; in particular $\End V$ and its trace-zero part $\ad V$ with $N_{\ad}=[N,\cdot\,]$; and the
twist $V\otimes\mu$ by a character $\mu$ of $W_K$ (with $N$ unchanged). The kernel $\ker N$ and the invariants $V^{r(I_K)}$ are stable under $r(W_K)$.
(Stated.) Every $(r,N)$ has a Frobenius-semisimplification, obtained by replacing $r(\Phi^n\sigma)$ by $r(\Phi)_{\mathrm{ss}}^nr(\sigma)$; see \cite[Section 8]{Deligne73},
\cite[Section 4.1]{Tate79}.
\end{proposition}

\begin{proof}
Each relation is a one-line check. For example $(r\otimes r')(w)(N\otimes1+1\otimes N')(r\otimes r')(w)^{-1}=\norm wN\otimes1+1\otimes\norm wN'$; for the dual,
$r(w^{-1})^{\mathsf T}(-N^{\mathsf T})r(w)^{\mathsf T}=-(r(w)Nr(w)^{-1})^{\mathsf T}$; and for $\ad$, $r(w)[N,x]r(w)^{-1}=\norm w[N,r(w)xr(w)^{-1}]$. Nilpotency is clear in each case
($N\otimes1$ and $1\otimes N'$ commute). If $Nv=0$ then $Nr(w)v=\norm w^{-1}r(w)Nv=0$; and $I_K$ is normal in $W_K$.
\end{proof}

\begin{theorem}[the Weil-Deligne representation of an $\ell$-adic representation]\label{thm:local:wd-equivalence}
Let $\rho$, $\zeta$, $N$ be as in \Cref{thm:local:grothendieck}, and fix a geometric Frobenius lift $\Phi\in W_K$. Put
\[
r(\Phi^n\sigma)=\rho(\Phi^n\sigma)\exp\bigl(-t_\zeta(\sigma)N\bigr)\qquad(n\in\Z,\ \sigma\in I_K).
\]
\begin{enumerate}[label=(\roman*)]
\item $\WD(\rho)=(r,N)$ is a Weil-Deligne representation over $E_\lambda$; $r=\rho$ on $\Phi$, and $r$ is trivial on $I_1$.
\item Its isomorphism class does not depend on $\Phi$ or on $\zeta$.
\item (Stated.) $\rho\mapsto\WD(\rho)$ is an equivalence between continuous representations of $W_K$ on finite-dimensional $E_\lambda$-spaces and Weil-Deligne
representations over $E_\lambda$; those coming from $G_K$ are the ones for which the eigenvalues of $r(\Phi)$ are $\ell$-adic units. It is compatible with
$\oplus$, $\otimes$, duals, $\ad$ and twists by characters; $\Q_\ell(1)$ goes to $(\norm{\cdot},0)$.
\item (Stated.) For the Tate module $V_\ell E$ of an elliptic curve over $K$, the Frobenius semisimplification of $\WD(V_\ell E)$, viewed over
an algebraic closure of $\Q_\ell$, is defined over $\Q$ and does not depend on $\ell\ne p$ (\Cref{ch:galrep}, \Cref{ch:modgalrep}).
\end{enumerate}
\end{theorem}

\begin{proof}
(i) $\exp(-tN)$ is defined for all $t\in\Z_\ell$ ($N$ nilpotent). Let $w=\Phi^n\sigma$, $w'=\Phi^{n'}\sigma'$; then $ww'=\Phi^{n+n'}\sigma''\sigma'$ with
$\sigma''=\Phi^{-n'}\sigma\Phi^{n'}$, and $t_\zeta(\sigma'')=\chi_\ell(\Phi^{-n'})t_\zeta(\sigma)=q^{n'}t_\zeta(\sigma)$ by \Cref{thm:local:tame-relation}(i). The relation
$\rho(w')^{-1}N\rho(w')=\norm{w'}^{-1}N=q^{n'}N$ gives $\exp(-tN)\rho(w')=\rho(w')\exp(-tq^{n'}N)$, so
\[
r(w)r(w')=\rho(w)\rho(w')\exp\bigl(-q^{n'}t_\zeta(\sigma)N\bigr)\exp\bigl(-t_\zeta(\sigma')N\bigr)=\rho(ww')\exp\bigl(-t_\zeta(\sigma''\sigma')N\bigr)=r(ww').
\]
On $I_1$, $\rho(\sigma)=\exp(t_\zeta(\sigma)N)$, so $r(\sigma)=1$; the kernel of $r|_{I_K}$ is open. $r(w)Nr(w)^{-1}=\rho(w)N\rho(w)^{-1}=\norm wN$, as $\exp(tN)$ commutes with $N$.
(ii) Changing $\zeta$ to $\zeta^a$ leaves $r$ unchanged and replaces $N$ by $aN$. Let $k\ge1$ be such that conjugation by $r(\Phi^k)$ is trivial on the finite group
$r(I_K)$ (it is an automorphism of finite order). Then $r(\Phi^k)$ commutes with $r(W_K)$, so its generalised eigenspaces $V_\lambda$ are $r(W_K)$-stable, and
$N(V_\lambda)\subseteq V_{q^{-k}\lambda}$ because $r(\Phi^k)N=q^{-k}Nr(\Phi^k)$. Choose a representative $\lambda_0$ in each orbit of $q^{k\Z}$ on the eigenvalues and let $A$ act on
$V_{\lambda_0q^{-kj}}$ by $a^j$. Then $A$ commutes with $r(W_K)$ and $ANA^{-1}=aN$, so $(r,N)\simeq(r,aN)$. Here the eigenvalues lie in a finite Galois
extension $F$ of $E_\lambda$, and $A$ is first defined on $V\otimes F$; it is defined over $E_\lambda$ if the representatives are chosen compatibly with
$\Gal(F/E_\lambda)$. This is possible: if $\gamma\in\Gal(F/E_\lambda)$ maps an orbit to itself, then $\gamma(\lambda_0)=\lambda_0q^{-kj}$ for some $j$, and since
$\gamma$ has finite order $m$, $\lambda_0=\gamma^m(\lambda_0)=\lambda_0q^{-kjm}$, so $j=0$. Choosing one $\lambda_0$ in each $\Gal(F/E_\lambda)$-orbit of orbits and
$\gamma(\lambda_0)$ for its transforms is therefore well defined, and then $\gamma$ maps $V_{\lambda_0q^{-kj}}\otimes F$ semilinearly onto
$V_{\gamma(\lambda_0)q^{-kj}}\otimes F$, so $A$ commutes with $\Gal(F/E_\lambda)$ and descends to $E_\lambda$ (as $a\in\Z_\ell^\times\subseteq E_\lambda$). If $\Phi$ is replaced by $\Phi'=\Phi\sigma_0$ ($\sigma_0\in I_K$), write
$\Phi'^n=\Phi^n\tau_n$ with $\tau_n\in I_K$ and put $c_n=t_\zeta(\tau_n)$. From $\Phi'^{n+1}=\Phi^n\tau_n\Phi\sigma_0=\Phi^{n+1}(\Phi^{-1}\tau_n\Phi)\sigma_0$ and
$t_\zeta(\Phi^{-1}\tau_n\Phi)=\chi_\ell(\Phi^{-1})c_n=qc_n$ (\Cref{thm:local:tame-relation}(i)) we get $c_{n+1}=qc_n+t_\zeta(\sigma_0)$, with $c_0=0$; this recursion
determines $c_n$ for all $n\in\Z$ (it can be run backwards), and $c_n=t_\zeta(\sigma_0)\frac{q^n-1}{q-1}$ satisfies it. For $w=\Phi'^n\sigma=\Phi^n\tau_n\sigma$ the two
definitions give $r'(w)=\rho(w)\exp(-t_\zeta(\sigma)N)$ and $r(w)=\rho(w)\exp(-(c_n+t_\zeta(\sigma))N)$, so
$r'(w)=r(w)\exp\bigl(t_\zeta(\sigma_0)\frac{q^n-1}{q-1}N\bigr)$. Put $c=t_\zeta(\sigma_0)/(q-1)\in\Q_\ell$ and $B=\exp(-cN)$, which commutes with $N$. Since
$\norm w=q^{-n}$, $r(w)^{-1}Nr(w)=q^nN$, hence $r(w)^{-1}\exp(cN)r(w)=\exp(cq^nN)$ and
\[
B^{-1}r(w)B=r(w)\cdot r(w)^{-1}\exp(cN)r(w)\cdot\exp(-cN)=r(w)\exp\bigl(c(q^n-1)N\bigr)=r'(w);
\]
so $B$ is an isomorphism $(r',N)\simeq(r,N)$. (iii), (iv) We do not prove these; see Deligne \cite[Section 8]{Deligne73}, Tate \cite[Section 4.2]{Tate79}
and Rohrlich \cite{Rohrlich94}; (iv) for elliptic curves follows from the description of $V_\ell E$ in \Cref{ch:reduction,ch:galrep}.
\end{proof}

\begin{definition}[the special representations]\label{def:local:sp}
For $n\ge1$, $\Sp(n)$ is the Weil-Deligne representation on $V=\bigoplus_{i=0}^{n-1}\Omega e_i$ with $r(w)e_i=\norm w^ie_i$ and $Ne_i=e_{i+1}$ ($e_n:=0$). For a
character $\mu$ of $W_K$, $\Sp(n)\otimes\mu$ is the \emph{special} (or \emph{Steinberg}) representation twisted by $\mu$; for $n=2$ and $\mu$ unramified it is an
\emph{unramified twist of Steinberg}, and for $\mu$ ramified a \emph{ramified twist}.
\end{definition}

\begin{proposition}\label{prop:local:sp}
\begin{enumerate}[label=(\roman*)]
\item $\Sp(n)\otimes\mu$ is an indecomposable Weil-Deligne representation. On $\Sp(2)$, $\Phi$ acts on $e_0,e_1$ by $1,q^{-1}$ and $\Frob$ by $1,q$, and $\ker N=\Omega e_1$.
\item If $\dim V=2$ and $N\ne0$, then $(r,N)$ is automatically Frobenius-semisimple, and $(r,N)\simeq\Sp(2)\otimes\mu$ for a unique character $\mu$ of $W_K$,
namely the character by which $W_K$ acts on $V/\ker N$.
\item (Stated.) A Frobenius-semisimple $(r,0)$ of dimension $2$ is a sum of two characters or irreducible; an irreducible one is induced from a character of
$W_{K'}$ for a quadratic $K'/K$, except possibly when $p=2$. See \cite[Section 4.1]{Tate79} and \cite{Rohrlich94}.
\end{enumerate}
\end{proposition}

\begin{proof}
(i) $r(w)Nr(w)^{-1}e_i=\norm w^{-i}r(w)e_{i+1}=\norm we_{i+1}$, so the relation holds. $\ker N=\Omega e_{n-1}$ is a line. If $V=V'\oplus V''$ as Weil-Deligne
representations, $\ker N=\ker N|_{V'}\oplus\ker N|_{V''}$, and a nilpotent operator on a nonzero space has a nonzero kernel; so $V'$ or $V''$ is $0$.
(ii) $N^2=0$ and $\ker N=\im N$ is a line, $r(W_K)$-stable (\Cref{prop:local:wd-constructions}). Let $\mu'$ be the character of $W_K$ on $\ker N$ and $\mu$ that on
$V/\ker N$. From $Nr(w)=\norm w^{-1}r(w)N$ applied to $v\notin\ker N$ we get $\mu(w)Nv=\norm w^{-1}\mu'(w)Nv$, so $\mu'=\norm{\cdot}\mu$. Since $\mu'(\Phi)=q^{-1}\mu(\Phi)\ne\mu(\Phi)$, the
operator $r(\Phi)$ has two distinct eigenvalues, so it is semisimple, and there is an eigenvector $e_0\notin\ker N$ with $r(\Phi)e_0=\mu(\Phi)e_0$. For $\sigma\in I_K$, $r(\sigma)$ has finite order, hence is semisimple, with the
eigenvalue $\mu(\sigma)=\mu'(\sigma)$ on $\ker N$ and on $V/\ker N$; so $r(\sigma)=\mu(\sigma)$ is a scalar. Hence $r(w)e_0=\mu(w)e_0$ for all $w=\Phi^n\sigma$. Put $e_1=Ne_0$; then
$r(w)e_1=\norm wNr(w)e_0=\norm w\mu(w)e_1$. In the basis $e_0,e_1$, $(r,N)=\Sp(2)\otimes\mu$; and $\mu$ is determined as the character on $V/\ker N$.
\end{proof}

\begin{example}[the Tate shape]\label{ex:local:tate-shape-wd}
Let $\ell\ne p$, $x\in K^\times$ with $a=v_K(x)$, $\varepsilon\colon G_K\to\{\pm1\}$ unramified, and let $\rho$ act on $\Q_\ell e_1\oplus\Q_\ell e_2$ by
\[
\rho(g)=\begin{pmatrix}\chi_\ell(g)\varepsilon(g)&\varepsilon(g)\kappa_x(g)\\0&\varepsilon(g)\end{pmatrix},
\]
with $\kappa_x$ read in $\Z_\ell$ through $\zeta$. This is a homomorphism: the upper right entry of $\rho(g)\rho(h)$ is $\varepsilon(gh)(\chi_\ell(g)\kappa_x(h)+\kappa_x(g))=\varepsilon(gh)\kappa_x(gh)$. (It is
the twist by $\varepsilon$ of the Kummer representation $\sm{\chi_\ell}{\kappa_x}{0}{1}$ of \Cref{ch:galois}; the corner of the twist is $\varepsilon\kappa_x$, and the matrix with
corner $\kappa_x$ and diagonal $\chi_\ell\varepsilon,\varepsilon$ is not a homomorphism unless $\varepsilon=1$.)
\begin{enumerate}[label=(\roman*)]
\item On $I_K$, $\rho(\sigma)=\sm{1}{a\,t_\zeta(\sigma)}{0}{1}=\exp(t_\zeta(\sigma)N)$ with $N=aE_{12}$ ($Ne_2=ae_1$, $Ne_1=0$), by \Cref{thm:local:kummer-inertia}(ii); here $I_1=I_K$, and
$N\ne0$ iff $a\ne0$. With $\rho(\Frob)=\sm{\varepsilon q}{c'}{0}{\varepsilon}$ one checks $\rho(\Frob)N\rho(\Frob)^{-1}=qN$: for an upper triangular
$g=\sm{x}{y}{0}{z}$ one has $gE_{12}=xE_{12}$ and $E_{12}g=zE_{12}$, so $gE_{12}g^{-1}=(x/z)E_{12}$, and here $x/z=\varepsilon q/\varepsilon=q$.
\item If $a\ne0$: $r$ is unramified, $r(\Phi)=\rho(\Frob)^{-1}$ has the distinct eigenvalues $\varepsilon q^{-1}$ (on $e_1$) and $\varepsilon$ (on the quotient), with $\varepsilon=\varepsilon(\Frob)$, so
$\WD(\rho)\simeq\Sp(2)\otimes\mu_\varepsilon$ with $\mu_\varepsilon$ unramified, $\mu_\varepsilon(\Phi)=\varepsilon(\Frob)$ (\Cref{prop:local:sp}(ii)). Moreover $V^{I_K}=\Q_\ell e_1$, on which $\Phi$ acts by
$\varepsilon q^{-1}$ and $\Frob$ by $\varepsilon q$.
\item If $a=0$: $\rho$ is unramified, $N=0$, and $\Frob$ has the distinct eigenvalues $\varepsilon q$, $\varepsilon$.
\item On the lattice $\Z_\ell e_1\oplus\Z_\ell e_2$, $I_K$ acts trivially modulo $\ell^k$ iff $\ell^k\mid a$ (\Cref{thm:local:kummer-inertia}(v)).
\end{enumerate}
\Cref{ch:reduction,ch:galrep} show that $T_\ell E$, at a prime of multiplicative reduction, has exactly this shape, with $x=q_E$ the Tate parameter and
$\varepsilon$ trivial (split) or the unramified quadratic character (nonsplit).
\end{example}

\begin{paperlink}\label{pl:local:steinberg}
In Lemma lem:euler (type $\mathrm{I}_n$) and in the proof of Lemma lem:twist(iii), the paper uses that the local component is an unramified twist of the Steinberg
representation. On the Galois side this means $\WD(V_\ell E|_{G_{\Q_q}})\simeq\Sp(2)\otimes\mu$ with $\mu$ unramified, which is \Cref{ex:local:tate-shape-wd}(ii) for the Tate
shape; the passage between the automorphic and the Galois side is \Cref{ch:modgalrep}. The proof of Lemma lem:tamfactor uses exactly the matrices of
\Cref{ex:local:tate-shape-wd}, with corner $\varepsilon\kappa$ (the corner of the revised paper; the original version wrote $\kappa$, which does not satisfy the cocycle relation for
the twisted matrices), $te_2=e_2+n'e_1$ with
$n'=n\,t_\zeta(t)$, and $v_\ell(n')=v_\ell(n)$ for a generator $t$ of tame inertia.
\end{paperlink}

\section{Inertia invariants, local \texorpdfstring{$L$}{L}-factors and conductors}\label{sec:local:L}

The paper's local terms are built from three invariants of $V=V_\ell E$ and of $A_f=\ad V$ at a bad prime: the inertia invariants (the paper's
$\delta_p=\dim V^{I_p}$), the local $L$-factor, read with the geometric Frobenius, and the conductor exponent. This section defines them, proves that
they depend only on the Weil-Deligne representation, and fixes the Frobenius convention for $L$-factors. The paper uses them in Section 1.1 (the
definitions of $\delta_p$, $L_p(A_f,s)$ and $L^{\nv}_p$), in Definition def:terms, Lemma lem:euler and the table of Section 1, in Lemma lem:fpr
($L_p(A_f,1)=L_p(A_f(1),0)$), in the paragraph after Lemma lem:dfg21, and in Lemmas lem:twist and lem:IV and Remark rem:types (conductor exponents).

\begin{proposition}[inertia invariants]\label{prop:local:invariants}
Let $\rho\colon G_K\to\GL(V)$ be continuous, $\ell\ne p$, and $\WD(\rho)=(r,N)$ as in \Cref{thm:local:wd-equivalence}. Then
\[
V^{\rho(I_K)}=(\ker N)^{r(I_K)},
\]
and $\rho(\Phi)=r(\Phi)$ on this space. So $\dim V^{I_K}$ and the characteristic polynomial of $\Phi$ on $V^{I_K}$ depend only on the isomorphism class of $\WD(\rho)$.
\end{proposition}

\begin{proof}
For $\sigma\in I_K$, $\rho(\sigma)=r(\sigma)\exp(t_\zeta(\sigma)N)$. If $Nv=0$ and $r(\sigma)v=v$ for all $\sigma$, then $\rho(\sigma)v=v$. Conversely let $v$ be fixed by $\rho(I_K)$. Take
$\sigma_1\in I_1$ with $t_\zeta(\sigma_1)=\ell^m\ne0$; then $v=\exp(\ell^mN)v$, and $\exp(\ell^mN)-1=\ell^mNU$ with $U=\sum_{k\ge0}(\ell^mN)^k/(k+1)!$ unipotent and commuting with $N$, so
$NUv=UNv=0$ and $Nv=0$. Then $r(\sigma)v=\rho(\sigma)\exp(-t_\zeta(\sigma)N)v=\rho(\sigma)v=v$. Finally $r(\Phi)=\rho(\Phi)$ by definition. An isomorphism of Weil-Deligne representations
carries $(\ker N)^{r(I_K)}$ to the corresponding space and intertwines $r(\Phi)$.
\end{proof}

For example, $(\ker N)^{r(I_K)}$ is $\Omega e_1$ for $\Sp(2)\otimes\mu$ with $\mu$ unramified, and $0$ for $\Sp(2)\otimes\mu$ with $\mu$ ramified ($r(I_K)$ acts on $\ker N=\Omega e_1$ by the
nontrivial $\mu|_{I_K}$); for $(r,0)$ it is $V^{r(I_K)}$. The paper's $\delta_p$ is $\dim V^{I_p}$ for $V=V_\ell E$ (Section 1.1).

\begin{definition}[local $L$-factors]\label{def:local:L-factor}
For $\rho$ as in \Cref{prop:local:invariants}, or a Weil-Deligne representation $(r,N)$, the \emph{local $L$-factor} is
\[
L(V,s)=\det\bigl(1-\Phi\,q^{-s}\bigm|V^{I_K}\bigr)^{-1}=\det\bigl(1-r(\Phi)\,q^{-s}\bigm|(\ker N)^{r(I_K)}\bigr)^{-1},
\]
with $\Phi=\Frob^{-1}$ the geometric Frobenius.
For $K=\Q_p$ and a representation $M$ of $G_\Q$ we write $L_p(M,s)$; this is the paper's $L_p(\cdot,s)$ of Section 1.1 (Appendix B, convention B-C2(d)).
\end{definition}

\begin{proposition}\label{prop:local:L-factor}
\begin{enumerate}[label=(\roman*)]
\item $L(V,s)$ is well defined (any two geometric Frobenius lifts agree on $V^{I_K}$) and depends only on the isomorphism class of $\WD(V)$.
\item $L(V\oplus V',s)=L(V,s)L(V',s)$.
\item For $\mu$ unramified with $\mu(\Phi)=\lambda$, $L(\mu,s)=(1-\lambda q^{-s})^{-1}$. In particular $L(\Q_\ell(1),s)=(1-q^{-1-s})^{-1}$, $L(V(1),s)=L(V,s+1)$ and
$L(V\otimes\norm{\cdot}^k,s)=L(V,s+k)$.
\item With the arithmetic $\Frob$ in place of $\Phi$ one would get $(1-q^{1-s})^{-1}$ for $\Q_\ell(1)$ and $L(V(1),s)=L(V,s-1)$.
\item $L(\Sp(2)\otimes\mu,s)=(1-\mu(\Phi)q^{-1-s})^{-1}$ for $\mu$ unramified, and $=1$ for $\mu$ ramified. If $E/\Q$ has good reduction at $p\ne\ell$ and
$\alpha,\beta$ are the roots of $X^2-a_pX+p$, the arithmetic Frobenius on $V_\ell E$ has eigenvalues $\alpha,\beta$ (\Cref{ch:ellcurves}), the geometric one
$\alpha^{-1},\beta^{-1}$, and
\begin{align*}
L_p(V_\ell E,s)&=(1-a_pp^{-1-s}+p^{-1-2s})^{-1}=L_p(E,s+1),\\
L_p(V_\ell E^\vee,s)&=(1-a_pp^{-s}+p^{1-2s})^{-1}=L_p(E,s).
\end{align*}
\end{enumerate}
\end{proposition}

\begin{proof}
(i) Two lifts differ by an element of $I_K$, which acts trivially on $V^{I_K}$; $V^{I_K}$ is $\Phi$-stable because $I_K$ is normal. The second expression and
\Cref{prop:local:invariants} give the independence. (ii) Invariants and determinants are additive. (iii) $I_K$ acts trivially on $\Q_\ell(1)$ and $\Phi$ by
$\chi_\ell(\Phi)=q^{-1}$; $V(1)^{I_K}=V^{I_K}(1)$, on which $\Phi$ acts by $q^{-1}$ times its action on $V^{I_K}$, so $\det(1-\Phi q^{-s}\mid V(1)^{I_K})=\det(1-\Phi q^{-s-1}\mid V^{I_K})$.
The same holds for $\norm{\cdot}^k$. (iv) $\Frob$ acts on $\Q_\ell(1)$ by $q$. (v) The first statement is \Cref{prop:local:invariants} and the example after it,
$\Phi$ acting on $e_1$ by $\mu(\Phi)\norm\Phi=\mu(\Phi)q^{-1}$. For $E$: $\det(1-\Phi X\mid V_\ell E)=(1-X/\alpha)(1-X/\beta)=1-a_pX/p+X^2/p$; with $X=p^{-s}$ this is
$1-a_pp^{-1-s}+p^{-1-2s}$, which is $1-a_pp^{-(s+1)}+p^{1-2(s+1)}$. On $V^\vee$ the geometric Frobenius has the inverse eigenvalues $\alpha,\beta$, giving
$1-a_pX+pX^2$.
\end{proof}

\begin{warning}[the paper's $L$-factors are geometric]\label{warn:local:geometric}
The paper's $L_p(A_f,s)$, the table of Section 1, Lemma lem:euler, Lemma lem:fpr ($L_p(A_f,1)^{-1}=L_p(A_f(1),0)^{-1}$), the values of PARI's
\texttt{lfunsympow} and the conventions of DFG04 \cite[Section 1.7.2]{DFG04} all use the geometric Frobenius: by \Cref{prop:local:L-factor}(iii), (iv) the identity
$L_p(M(1),s)=L_p(M,s+1)$ of Lemma lem:fpr holds for $\Phi$ and fails for $\Frob$, and at a multiplicative prime the factor $1-q^{-1-s}$ of the table
requires $\Phi$ (\Cref{prop:local:steinberg-shape}(iii)). On $\ad V$ at a good prime the eigenvalues $\alpha/\beta,1,\beta/\alpha$ are stable under inversion, so there
the choice is invisible.
\end{warning}

\begin{paperlink}\label{pl:local:Lfactor}
Section 1.1 defines $L_p(A_f,s)=\det(1-\Frob_p^{-1}p^{-s}\mid A_f^{I_p})^{-1}$: this is \Cref{def:local:L-factor} for $M=A_f$, the geometric Frobenius
at $p$ being $\Frob_p^{-1}$. The naive factor
$L_p^{\nv}(A_f,s)$ is $L_p(A_f,s)$ when $\delta_p=\dim V^{I_p}>0$ and $1$ when $\delta_p=0$. The paragraph after Lemma lem:dfg21 writes
``$\det(1-\Frob_p^{-1}\mid(V^D)^{I_p})=L_p(A_f,1)^{-1}\ne0$'' with $V^D=A_f(1)$: the left side is $L_p(A_f(1),0)^{-1}$, which equals $L_p(A_f,1)^{-1}$ by
\Cref{prop:local:L-factor}(iii). The original version of the paper wrote $\Frob_p$ in this sentence; it has to be read as the geometric Frobenius.
\end{paperlink}

\subsection{Swan and Artin conductors}

\begin{definition}[Swan and Artin conductors]\label{def:local:swan}
Let $V$ be a finite-dimensional representation over a field $\Omega$ of $G_K$, $W_K$ or $I_K$, on which $I_K$ acts through a finite quotient. Choose a finite Galois
$L/K$ such that $I_K\cap G_L$ acts trivially (possible because the kernel is open in $I_K$), so that $I_K$ acts through its image $G_0=I(L/K)$ in $\Gal(L/K)$
(\Cref{thm:local:inertia-structure}(i)); let $G_i=G_i(L/K)$ and $g_i=\#G_i$. The \emph{Swan conductor} and the \emph{Artin conductor} (exponent) are
\[
\operatorname{sw}(V)=\sum_{i\ge1}\frac{g_i}{g_0}\dim\bigl(V/V^{G_i}\bigr),\qquad a(V)=\sum_{i\ge0}\frac{g_i}{g_0}\dim\bigl(V/V^{G_i}\bigr)=\dim V-\dim V^{I_K}+\operatorname{sw}(V).
\]
For a continuous $\ell$-adic $\rho$ ($\ell\ne p$), $P_K$ acts through a finite quotient (\Cref{prop:local:wild-image}), and $\operatorname{sw}(\rho)$ is defined by the same sum,
which involves only the $G_i$ with $i\ge1$, the images of $P_K$.
\end{definition}

\begin{theorem}[properties of the conductors]\label{thm:local:swan}
\begin{enumerate}[label=(\roman*)]
\item (Stated.) $\operatorname{sw}(V)$ and $a(V)$ do not depend on the choice of $L$ (by Herbrand's theorem), and in upper numbering
\[
\operatorname{sw}(V)=\int_0^\infty\dim\bigl(V/V^{G_K^u}\bigr)\,du .
\]
\item (Stated; Artin.) $\operatorname{sw}(V)$ and $a(V)$ are nonnegative integers.
\item $\operatorname{sw}(V)=0$ iff $P_K$ acts trivially on $V$; $\operatorname{sw}(V)$ depends only on $V|_{P_K}$; both are additive in direct sums, unchanged by twisting by an
unramified character and by restriction to $G_{K'}$ for $K'/K$ finite unramified; and if $V$ is tamely ramified, $a(V)=\dim V-\dim V^{I_K}$.
\item For a character $\chi$: $a(\chi)=0$ if $\chi$ is unramified, $a(\chi)=1$ if it is tamely ramified, and $a(\chi)=1+(\text{the largest }i\text{ with }\chi(G_i)\ne1)$ for
the quadratic characters of \Cref{ex:local:ramgroups}. Hence $a(\chi_{-3}|_{G_{\Q_3}})=1$, $a(\chi_{q^*}|_{G_{\Q_q}})=1$ for $q$ odd, $a(\chi_{-4}|_{G_{\Q_2}})=2$ and
$a(\chi_{\pm8}|_{G_{\Q_2}})=3$, in agreement with \Cref{thm:local:lcft}(iii) and with the conductors $3$, $q$, $4$, $8$. (Stated: for a general wildly ramified
character, $a(\chi)=1+c$ with $c$ the largest upper jump.)
\item The same definitions apply verbatim to representations over $\F_\ell$, where the image of inertia is finite; this gives the local factor at $q$ of
the Serre conductor of $\bar\rho_{E,\ell}$ (\Cref{ch:galrep}).
\end{enumerate}
\end{theorem}

\begin{proof}[Proof of (iii), (iv); references for (i), (ii)]
(iii) If $P_K$ acts trivially, so does each $G_i$, $i\ge1$, and all terms vanish. If $\operatorname{sw}(V)=0$, then either $G_1=1$, or the term $i=1$ (with $g_1/g_0>0$) vanishes
and $G_1$ acts trivially; in both cases $P_K$, which maps onto $G_1$, acts trivially. The sum for $i\ge1$ involves only $V|_{G_1}$, the image of $V|_{P_K}$. Additivity
holds because $(V\oplus V')^{G_i}=V^{G_i}\oplus V'^{G_i}$. An unramified character is trivial on $I_K$, so it does not change any $V^{G_i}$. If $K'/K$ is unramified and
$L'=LK'$, the inertia group of $L'/K$ fixes $K'$, so the ramification groups of $L'/K'$, which are $G_i(L'/K)\cap\Gal(L'/K')$ (\Cref{prop:local:ramgroups}(i)), equal
$G_i(L'/K)$ for $i\ge0$; by (i) one may compute both conductors with $L'$. If $V$ is tame, $P_K$ acts trivially and $\operatorname{sw}=0$.
(iv) For a character, $\dim V/V^{G_i}$ is $1$ if $\chi(G_i)\ne1$ and $0$ otherwise. If $\chi$ is tamely ramified, $\chi(G_0)\ne1=\chi(G_1)$ and $a=1$. For the
quadratic characters of \Cref{ex:local:ramgroups} the field $L$ is the quadratic extension itself, $g_i/g_0=1$ while $G_i=G$; so $a=1+1=2$ for the jump
$1$ ($\Q_2(i)$, $\Q_2(\sqrt3)$) and $a=1+2=3$ for the jump $2$ ($\Q_2(\sqrt{\pm2})$, $\Q_2(\sqrt{\pm10})$). The values for the global characters follow from
\Cref{prop:local:global-quadratic-local}(ii), (iii). (i) and (ii) are Serre \cite[Chapter IV, Section 3, and Chapter VI, Section 2]{SerreLocalFields};
the integrality is Artin's theorem, proved there through the Artin character (see also \cite[Chapter 19]{SerreLinearReps}). (v) is a definition.
\end{proof}

\begin{definition}[conductor of a Weil-Deligne representation]\label{def:local:conductor}
The conductor exponent of a Weil-Deligne representation $(r,N,V)$ is
\[
a(r,N)=a(r)+\dim V^{r(I_K)}-\dim(\ker N)^{r(I_K)}=\dim V-\dim(\ker N)^{r(I_K)}+\operatorname{sw}(r),
\]
and for a continuous $\ell$-adic $\rho$, $a(\rho)=a(\WD(\rho))$.
\end{definition}

\begin{proposition}\label{prop:local:conductor}
\begin{enumerate}[label=(\roman*)]
\item $a(\rho)=\dim V-\dim V^{\rho(I_K)}+\operatorname{sw}(\rho)$, where $\operatorname{sw}(\rho)=\operatorname{sw}(\rho|_{P_K})$.
\item $a$ is additive in direct sums, unchanged by unramified twists and by unramified base change, $a(\rho)=0$ iff $\rho$ is unramified, and
$a(\rho)=\dim V-\dim V^{I_K}$ if $\rho$ is tamely ramified.
\item If $V$ is unramified and $\chi$ is a character, $a(V\otimes\chi)=\dim V\cdot a(\chi)$. If $\chi$ is ramified and $\mu$ unramified, $a(\Sp(2)\otimes\mu)=1$ and
$a(\Sp(2)\otimes\mu\chi)=2a(\chi)$.
\item For an elliptic curve $E/\Q$ and $\ell\ne p$, \Cref{ch:reduction} defines the exponent of $p$ in the conductor of $E$ to be $a(V_\ell E|_{G_{\Q_p}})$;
it does not depend on $\ell$. (Stated.) Ogg's formula (Saito in general) gives it as $v_p(\Delta_{\min})+1-m$, where $m$ is the number of irreducible
components over $\Fbar_p$ of the special fibre of the minimal regular model of $E$ over $\Z_p$ (\Cref{ch:reduction}).
\end{enumerate}
\end{proposition}

\begin{proof}
(i) By \Cref{prop:local:invariants}, $(\ker N)^{r(I_K)}=V^{\rho(I_K)}$; and $r=\rho$ on $P_K$, because $t_\zeta$ vanishes on $P_K$; so $\operatorname{sw}(r)=\operatorname{sw}(\rho)$.
(ii) follows from (i), \Cref{thm:local:swan}(iii) and \Cref{prop:local:unram-basechange}; $a(\rho)=0$ forces $V^{I_K}=V$. (iii) $I_K$ acts on $V\otimes\chi$ by the
scalars $\chi(\sigma)$, so $(V\otimes\chi)^{G_i}$ is $V\otimes\chi$ or $0$ according as $\chi(G_i)=1$ or not, and each term of the sum is $\dim V$ times that for $\chi$. For
$\Sp(2)\otimes\mu$, $r$ is unramified and $\dim(\ker N)^{r(I_K)}=1$, so $a=2-1+0=1$. For $\Sp(2)\otimes\mu\chi$, $r=(\mu\chi\oplus\norm{\cdot}\mu\chi)$ has $a(r)=2a(\chi)$
by the previous case, and $(\ker N)^{r(I_K)}=V^{r(I_K)}=0$, so $a=2a(\chi)$. (iv) The first sentence is a definition. We do not prove the independence of $\ell$ or Ogg's formula; see \cite[Section 2]{SerreTate68}, Ogg \cite{Ogg}
and \Cref{ch:reduction}, where Saito's general case is cited.
\end{proof}

\begin{theorem}[epsilon factors]\label{thm:local:epsilon}
(Langlands, Deligne.) Fix a nontrivial additive character $\psi$ of $K$ and a Haar measure $dx$ on $K$. There is a unique way to attach to each representation $r$
of $W_K$ over $\C$ with $r(I_K)$ finite a constant $\epsilon(r,\psi,dx)\in\C^\times$ that is multiplicative in short exact sequences, inductive in degree $0$, and
equal to Tate's local constant for characters. For a Weil-Deligne representation one puts
\[
\epsilon(r,N,\psi,dx)=\epsilon(r,\psi,dx)\det\bigl(-r(\Phi)\bigm|V^{r(I_K)}/(\ker N)^{r(I_K)}\bigr).
\]
We do not prove this; see Deligne \cite[Sections 4, 5 and 8]{Deligne73} and Tate \cite[Sections 3 and 4]{Tate79} (Deligne's normalisation of the Artin
map is the geometric one, \Cref{warn:local:artin}). The paper does not use epsilon factors; \Cref{ch:modgalrep,ch:lfunctions} recall this statement for
root numbers and functional equations.
\end{theorem}

\section{The local shapes and their adjoints}\label{sec:local:shapes}

This section computes, for each shape of local representation that the paper meets, the inertia invariants, the conductor exponent and the
$L$-factor of $V$ and of its trace-zero adjoint $A=\ad V$. It is the Galois-theoretic content of Lemma lem:euler, of the table of Section 1, of
Proposition prop:shapes and of the conductor statements in Lemmas lem:IV and lem:twist. The linear algebra of the inertia actions (unipotent
elements, commutants, the table of cyclic inertia shapes) is \Cref{ch:reps}; the dictionary between shapes and Kodaira types is proved in
\Cref{ch:reduction,ch:galrep} and only quoted here.

Standing hypotheses: $K=\Q_q$, $\ell\ne q$, $V$ a two-dimensional representation of $G_{\Q_q}$ over a finite extension $\Omega$ of $\Q_\ell$ with $\det V|_{I_q}=1$
(true for $V_\ell E$, whose determinant $\chi_\ell$ is unramified at $q$), $A=\ad V$ with $g\cdot x=\rho(g)x\rho(g)^{-1}$, and $\Phi$ the geometric Frobenius. In the
basis $e_1,e_2$ of \Cref{ex:local:tate-shape-wd}, $e=e_1\otimes e_2^\vee$, $h=e_1\otimes e_1^\vee-e_2\otimes e_2^\vee$, $f=e_2\otimes e_1^\vee$ is the paper's basis of $\ad T$ (here $e$
is a vector, not a ramification index).

\begin{proposition}[unramified $V$]\label{prop:local:unramified-shape}
If $V$ is unramified with arithmetic Frobenius eigenvalues $\alpha,\beta$, then $\WD(V)=(r,0)$, $a(V)=0$, $L(V,s)^{-1}=(1-\alpha^{-1}q^{-s})(1-\beta^{-1}q^{-s})$, and $A$ is
unramified with eigenvalues $\alpha/\beta,1,\beta/\alpha$ for both $\Frob$ and $\Phi$. Hence
\[
L(A,s)^{-1}=\Bigl(1-\frac\alpha\beta q^{-s}\Bigr)(1-q^{-s})\Bigl(1-\frac\beta\alpha q^{-s}\Bigr),\qquad L(A,1)^{-1}=\frac{q-1}q\cdot\frac{(q+1)^2-a^2}{q^2},
\]
the last formula when $\alpha+\beta=a$ and $\alpha\beta=q$.
\end{proposition}

\begin{proof}
$N=0$ since $\rho(I_q)=1$ is finite, and $a=0$ by \Cref{prop:local:conductor}(ii). The eigenvalues of $\Phi$ on $V$ are $\alpha^{-1},\beta^{-1}$. If $\Frob$ has eigenvectors
$v_1,v_2$ (over an extension of $\Omega$, or after triangularising), it acts on $v_i\otimes v_j^\vee$ by $\alpha_i/\alpha_j$; the trace-zero part has eigenvalues $\alpha/\beta$,
$1$, $\beta/\alpha$, a set stable under inversion. For $s=1$, $(1-\frac\alpha{\beta q})(1-\frac\beta{\alpha q})=1-\frac{\alpha^2+\beta^2}{\alpha\beta q}+\frac1{q^2}=1-\frac{a^2-2q}{q^2}+\frac1{q^2}=\frac{(q+1)^2-a^2}{q^2}$
(\Cref{ch:rings}), using $\alpha^2+\beta^2=(\alpha+\beta)^2-2\alpha\beta=a^2-2q$; the middle factor contributes $1-q^{-1}=\frac{q-1}q$.
\end{proof}

For example, \curve{11a1} has good reduction at $2$ with $a_2=-2$ (PARI, \texttt{ellap}), so $L_2(A,1)^{-1}=\frac12\cdot\frac{9-4}4=\frac58$. Note that
$(q+1)^2-a^2=(q+1-a)(q+1+a)$, and $q+1-a$ is the number of points of the reduced curve over $\F_q$ (\Cref{ch:ellcurves}); here
$\#\tilde E(\F_2)=5$, and the factor $5$ in $\frac58$ is that count.

\begin{proposition}[finite cyclic tame image of inertia]\label{prop:local:tame-cyclic-shape}
Suppose $\rho(I_q)$ is cyclic of order $m\in\{2,3,4,6\}$ with $q\nmid m$, generated by $g$. Then $N=0$, the ramification is tame, and:
\begin{enumerate}[label=(\roman*)]
\item If $m=2$, then $g=-1$, $V^{I_q}=0$, $a(V)=2$ and $L(V,s)=1$. For $q$ odd and any ramified quadratic character $\chi$ of $G_{\Q_q}$, $V\otimes\chi$ is unramified, and
$A=\ad(V\otimes\chi)$ is unramified, with $L(A,s)$ as in \Cref{prop:local:unramified-shape} for $V\otimes\chi$.
\item If $m\in\{3,4,6\}$, $g$ has the eigenvalues $\zeta,\zeta^{-1}$ with $\zeta$ a primitive $m$-th root of unity; $V^{I_q}=0$, $a(V)=2$ and $L(V,s)=1$. Put $m'=3$, $\chi=\chi_{-3}$ if
$m\in\{3,6\}$, and $m'=4$, $\chi=\chi_{-4}$ if $m=4$. Then $\rho(\Frob)g\rho(\Frob)^{-1}=g^{\chi(q)}$. If $\chi(q)=1$, $V\otimes\Omega(\zeta)$ is a sum of two characters of $G_{\Q_q}$
restricting to $\zeta^{\pm1}$ on $g$; if $\chi(q)=-1$, $\rho(\Frob)$ exchanges the two eigenlines of $g$, and $V\otimes\Omega(\zeta)$ is induced from a character of
$G_{K_2}$, $K_2$ the unramified quadratic extension of $\Q_q$.
\item In case (ii), $A^{I_q}$ is the line spanned by $2g-\tr g$ (which is $\pm(2\sigma+1)$, of square $-3$, for $m\in\{3,6\}$ with $\sigma$ the element of order $3$ in
$\langle g\rangle$, and $g$ itself for $m=4$), and both $\Frob$ and $\Phi$ act on it by $\chi(q)\in\{\pm1\}$. Hence
\[
L(A,s)^{-1}=1-\chi(q)\,q^{-s},\qquad a(A)=2 .
\]
\end{enumerate}
\end{proposition}

\begin{proof}
$\rho(I_q)$ is finite, so $N=0$ (\Cref{thm:local:grothendieck}), and of order prime to $q$, so $\rho(P_q)=1$ (\Cref{cor:local:finite-tame-image}). Since $\det g=1$, the
eigenvalues of $g$ (semisimple, of finite order in characteristic $0$) are $\lambda,\lambda^{-1}$ with $\lambda$ of order $m$. (i) For $m=2$, $\lambda=\lambda^{-1}$ forces $\lambda=-1$
and $g=-1$. In all cases $1$ is not an eigenvalue, so $V^{I_q}=0$, $L(V,s)=1$ and $a(V)=2-0=2$ (tame). If $q$ is odd and $\chi$ ramified quadratic, then
$\chi|_{I_q}=\theta_2$ (\Cref{prop:local:quadratic-ext}(ii)), and $\rho|_{I_q}$ is a character of $I_q$ of order $2$ times the identity, so it is $\theta_2\cdot1$; hence
$(V\otimes\chi)|_{I_q}=\theta_2^2=1$. Twisting by a character does not change $\ad$. (ii) By \Cref{cor:local:finite-tame-image},
$\rho(\Frob)g\rho(\Frob)^{-1}=g^q$. For $m\in\{3,4\}$, $q\equiv\chi(q)\pmod{m'}$ and $g^q=g^{\chi(q)}$, as $g$ has order $m=m'$. For $m=6$, $q$ is prime to $6$, so $q\equiv\pm1\pmod6$ and
again $g^q=g^{\pm1}$ with the sign $\chi_{-3}(q)$. If $\rho(\Frob)$ commutes with $g$, it preserves the two eigenlines of $g$ in $V\otimes\Omega(\zeta)$, and so does $\rho(I_q)=\langle g\rangle$;
by \Cref{cor:local:split} these lines are stable under $G_{\Q_q}$. If $\rho(\Frob)g\rho(\Frob)^{-1}=g^{-1}$, $\rho(\Frob)$ maps the $\zeta$-eigenline $D$ to the $\zeta^{-1}$-eigenline;
$D$ is stable under the index-two subgroup $G_{K_2}=I_q\rtimes\overline{\langle\Frob^2\rangle}$, and $V\otimes\Omega(\zeta)=D\oplus\rho(\Frob)D$ is induced from the character of $G_{K_2}$ on
$D$. (iii) The commutant of $g$ in $\End V$ is $\Omega+\Omega g$, since $g$ has distinct eigenvalues (\Cref{ch:reps}), so its trace-zero part is spanned by
$2g-\tr(g)$; and $A^{I_q}$ is that trace-zero commutant. By Cayley-Hamilton, $g^{-1}=\tr(g)-g$, so $2g^{-1}-\tr g=-(2g-\tr g)$; therefore $\Frob$, which conjugates
$g$ to $g^{\chi(q)}$, acts on the line by $\chi(q)$, and so does $\Phi=\Frob^{-1}$. For $m=3$, $\tr\sigma=-1$ and $(2\sigma+1)^2=4(\sigma^2+\sigma)+1=-3$; for $m=6$,
$g=-\sigma$; for $m=4$, $\tr g=0$. Finally $A$ is tame (its inertia image is a quotient of $\langle g\rangle$), so $a(A)=3-1=2$.
\end{proof}

\begin{proposition}[Steinberg and its twists]\label{prop:local:steinberg-shape}
Let $\mu$ be an unramified and $\chi$ a ramified quadratic character of $G_{\Q_q}$.
\begin{enumerate}[label=(\roman*)]
\item $V=\Sp(2)\otimes\mu$: $V^{I_q}=\ker N$ is a line on which $\Phi$ acts by $\mu(\Phi)q^{-1}$ (and $\Frob$ by $\mu(\Frob)q$); $a(V)=1$; $L(V,s)=(1-\mu(\Phi)q^{-1-s})^{-1}$.
\item $V=\Sp(2)\otimes\mu\chi$: $V^{I_q}=0$, $L(V,s)=1$ and $a(V)=2a(\chi)$, which is $2$ for $q$ odd and $4$ or $6$ for $q=2$.
\item For $V=\Sp(2)\otimes\mu$ and $V=\Sp(2)\otimes\mu\chi$, $\ad V=\ad\Sp(2)\simeq\Sp(3)\otimes\norm{\cdot}^{-1}$. In the basis $E_{01}$, $h=E_{00}-E_{11}$, $E_{10}=N$ (with
$E_{ij}=e_i\otimes e_j^\vee$ in the basis of \Cref{def:local:sp}), $\Phi$ acts by $q,1,q^{-1}$, and $[N,E_{01}]=-h$, $[N,h]=2N$, $[N,N]=0$. So $A^{I_q}=\ker[N,\cdot\,]=\Omega N$, $\Phi$ acts on it
by $q^{-1}$ ($\Frob$ by $q$), $a(A)=2$, and
\[
L(\ad V,s)=(1-q^{-1-s})^{-1}.
\]
With the arithmetic Frobenius one would get $(1-q^{1-s})^{-1}$.
\item The line $\Omega N$ is the line of $e=e_1\otimes e_2^\vee$ in the basis of \Cref{ex:local:tate-shape-wd}, the line $\Z_\ell e$ of the proof of Lemma lem:tamfactor.
\end{enumerate}
\end{proposition}

\begin{proof}
(i) and (ii) are \Cref{prop:local:invariants}, \Cref{prop:local:L-factor}(v) and \Cref{prop:local:conductor}(iii), with $a(\chi)=1$ for $q$ odd and $a(\chi)\in\{2,3\}$ for $q=2$
(\Cref{thm:local:swan}(iv)). (iii) Twisting by a character does not change $\ad$. $r(\Phi)$ acts on $e_0,e_1$ by $1,q^{-1}$, hence on $E_{ij}$ by the ratio of the eigenvalues
of $e_i$ and $e_j$: $q$ on $E_{01}$, $1$ on $h$, $q^{-1}$ on $E_{10}$. The brackets: $E_{10}E_{01}-E_{01}E_{10}=E_{11}-E_{00}=-h$; $E_{10}h-hE_{10}=E_{10}+E_{10}=2E_{10}$. So
$f_0=E_{01}$, $f_1=-h$, $f_2=-2E_{10}$ satisfy $N_{\ad}f_i=f_{i+1}$, with $\Phi$-eigenvalues $q^{1-i}=\norm\Phi^{i-1}$: this is $\Sp(3)\otimes\norm{\cdot}^{-1}$. In characteristic $0$,
$\ker N_{\ad}=\Omega E_{10}=\Omega N$, and $r(I_q)$ acts trivially on $\ad$ (it acts on $V$ by scalars), so $A^{I_q}=\Omega N$ (\Cref{prop:local:invariants}), with $\Phi$-eigenvalue $q^{-1}$.
$a(A)=3-1+0=2$. (iv) In \Cref{ex:local:tate-shape-wd}, $N=aE_{12}$ maps $e_2$ to $e_1$, so $N$ is a multiple of $e_1\otimes e_2^\vee$.
\end{proof}

\begin{theorem}[the shapes at a tame prime]\label{thm:local:shapes-table}
Under the standing hypotheses, suppose that $q\ge5$. (Columns $1$ to $7$, which do not involve Kodaira types, hold more generally, with the same proofs:
rows $1$ and $3$ for every $q$, rows $2$ and $4$ for every odd $q$, and rows $5$ and $6$ whenever $q\nmid\#\rho(I_q)$. The Kodaira column is asserted
only for $q\ge5$.) The following table
lists, for each shape of $\WD(V)$, the image of inertia, $\dim V^{I_q}$, $a(V)$, $L(V,s)^{-1}$, $\dim A^{I_q}$ and $L(A,s)^{-1}$ ($\Phi$ geometric); the last column,
the Kodaira types of an elliptic curve with $V=V_\ell E$, is quoted from \Cref{ch:reduction,ch:galrep}.
\begin{center}\footnotesize\setlength{\tabcolsep}{3pt}
\begin{tabular}{llcclcll}
\toprule
$\WD(V)$ & $\rho(I_q)$ & $\delta_q$ & $a(V)$ & $L(V,s)^{-1}$ & $\dim A^{I_q}$ & $L(A,s)^{-1}$ & types\\
\midrule
unramified & $1$ & $2$ & $0$ & $(1-\frac X\alpha)(1-\frac X\beta)$ & $3$ & $(1-\frac\alpha\beta X)(1-X)(1-\frac\beta\alpha X)$ & $\mathrm{I}_0$\\
unram.${}\otimes\chi$ & $\{\pm1\}$ & $0$ & $2$ & $1$ & $3$ & as in row $1$ & $\mathrm{I}_0^*$\\
$\Sp(2)\otimes\mu$ & unipotent & $1$ & $1$ & $1-\mu(\Phi)q^{-1}X$ & $1$ & $1-q^{-1}X$ & $\mathrm{I}_n$\\
$\Sp(2)\otimes\mu\chi$ & $\pm$unipotent & $0$ & $2$ & $1$ & $1$ & $1-q^{-1}X$ & $\mathrm{I}_n^*$\\
order $3$ or $6$ & $\langle\sigma\rangle$, $\langle-\sigma\rangle$ & $0$ & $2$ & $1$ & $1$ & $1-\chi_{-3}(q)X$ & $\mathrm{IV},\mathrm{IV}^*$; $\mathrm{II},\mathrm{II}^*$\\
order $4$ & $\langle\tau\rangle$ & $0$ & $2$ & $1$ & $1$ & $1-\chi_{-4}(q)X$ & $\mathrm{III},\mathrm{III}^*$\\
\bottomrule
\end{tabular}
\end{center}
Here $X=q^{-s}$, $\delta_q=\dim V^{I_q}$, $\alpha,\beta$ are the arithmetic Frobenius eigenvalues of the unramified representation (in the second row, of $V\otimes\chi$),
$\mu$ is unramified, $\chi$ is a ramified quadratic character, $N=0$ in the last two rows, and $n\ge1$. So $\delta_q$ is $2,0,1,0,0,0$ down the rows; the naive factor is $L(A,s)$ when $\delta_q>0$ and $1$ when $\delta_q=0$;
and the factor that the naive $L$-function omits is $L(A,s)^{-1}$ in rows $2$, $4$, $5$, $6$ and nothing in rows $1$, $3$.
\end{theorem}

\begin{proof}
Every column but the last is \Cref{prop:local:unramified-shape,prop:local:tame-cyclic-shape,prop:local:steinberg-shape}. That these are all the shapes with $q\nmid\#\rho(I_q)$
when $\rho(I_q)$ is finite of order in $\{1,2,3,4,6\}$ follows from \Cref{prop:local:sp}(ii) when $N\ne0$ (and $\det V|_{I_q}=1$ forces $\mu^2$ to be trivial on
$I_q$, so $\mu$ is unramified or an unramified twist of a ramified quadratic character, \Cref{cor:local:split}), and from \Cref{ch:reps} when $N=0$. The orders $1,2,3,4,6$ and the Kodaira column are
\Cref{ch:reduction,ch:galrep} (Serre and Tate \cite{SerreTate68}, Kraus \cite{Kraus90}); we do not prove them here.
\end{proof}

\begin{paperlink}\label{pl:local:euler}
Lemma lem:euler and the table of Section 1 are the last column of \Cref{thm:local:shapes-table} read through the Kodaira types: $L_q(A_f,s)^{-1}$ is $1-q^{-1-s}$ for
$\mathrm{I}_n$ and $\mathrm{I}_n^*$, $1-\chi_{-3}(q)q^{-s}$ for $\mathrm{II},\mathrm{II}^*,\mathrm{IV},\mathrm{IV}^*$, $1-\chi_{-4}(q)q^{-s}$ for $\mathrm{III},\mathrm{III}^*$ and the good
factor of the twist for $\mathrm{I}_0^*$; $e_\ell(E,q)=v_\ell(L_q(A_f,1)^{-1})$ in the rows with $\delta_q=0$ gives $v_\ell(q^2-1)$, $v_\ell(q-\chi_{-3}(q))$, $v_\ell(q-\chi_{-4}(q))$ and
$v_\ell(q-1)+v_\ell((q+1)^2-a_q(E')^2)$, and $e_\ell=0$ in the rows with $\delta_q>0$. ``Frobenius centralises $\sigma$ if $q\equiv1\pmod3$ and inverts it if $q\equiv2$, for
either choice of Frobenius'' is \Cref{prop:local:tame-cyclic-shape}(ii), (iii). ``$(\ad V)^{I_q}$ is a line on which the geometric Frobenius has eigenvalue $q^{-1}$''
(type $\mathrm{I}_n$) is \Cref{prop:local:steinberg-shape}(iii). Proposition prop:shapes(i) to (iii) is the same computation organised by $P_q(X)=L_q(A_f,s-1)^{-1}$,
$X=q^{-s}$: $1-X$ for the Steinberg rows, $(1-qX)(1-(a^2-2q)X+q^2X^2)$ for row 2, $1-cqX$ with $c=\chi(q)$ for rows 5 and 6. In Lemma lem:twist(iii), the fact that the twist of an
unramified twist of the Steinberg representation by a ramified quadratic character has conductor exponent $2$ ($q$ odd) is
\Cref{prop:local:steinberg-shape}(ii). In Lemma lem:twist(ii), ``the conductor exponents are the tame ones, $\dim V-\dim V^{I_q}$'' holds for the exchanges listed there
because twisting a tame representation by the tame $\chi_{q^*}$ keeps it tame (\Cref{prop:local:global-quadratic-local}(v)); at $q=3$ a general additive curve can be
wild (\Cref{rem:local:wild-ec}). In Lemma lem:IV, ``the conductor exponent of $E$ is $\dim V-\dim V^{I_q}=2$, there being no wild part'' is row 5 with
\Cref{prop:local:conductor}(ii), and ``the Serre conductor exponent of $\bar\rho_{E,3}$ at $q$ is $\dim E[3]-\dim E[3]^{I_q}=1$'' is \Cref{thm:local:swan}(v) for the tame
mod $3$ representation.
\end{paperlink}

\begin{proposition}[the primes $2$ and $3$]\label{prop:local:wild-shapes}
Keep the standing hypotheses, with $q\in\{2,3\}$.
\begin{enumerate}[label=(\roman*)]
\item If $q=3$ and $\rho(I_3)$ is finite of order divisible by $3$, then $\rho$ is wildly ramified and $a(V)\ge3$. If moreover $\rho(I_3)$ is cyclic of order $3$ or
$6$, $A^{I_3}$ is still the line spanned by $2\sigma+1$, and $\Frob$ still acts on it by $+1$ or $-1$; but the sign is no longer determined by $q$,
since the tame relation is not available on wild inertia.
\item If $q=2$ and $\rho(I_2)$ is cyclic of order $3$, then \Cref{prop:local:tame-cyclic-shape} applies: $a(V)=2$, $L(A,s)^{-1}=1-\chi_{-3}(2)2^{-s}=1+2^{-s}$, and
$L(A,1)^{-1}=\frac32$.
\item If $\rho(I_q)$ is finite and not abelian, then $V$ is an absolutely irreducible representation of $\rho(I_q)$, $A^{I_q}=0$ and $L(A,s)=1$.
\item For $q=2$, $a(\Sp(2)\otimes\mu\chi)$ is $4$ or $6$ according as $\chi$ has ramification jump $1$ or $2$.
\end{enumerate}
\end{proposition}

\begin{proof}
(i) $I_3/P_3$ has no quotient of order $3$ (\Cref{thm:local:tame-character}(ii)), so an element of order $3$ in $\rho(I_3)$ comes from $P_3$ (\Cref{prop:local:wild-image}(ii)),
and $\operatorname{sw}(V)\ge\frac{g_1}{g_0}\dim(V/V^{G_1})>0$. An element of order $3$ and determinant $1$ has the eigenvalues $\zeta_3^{\pm1}$ and fixes no vector, so
$V^{I_3}=0$; as $\operatorname{sw}$ is an integer (\Cref{thm:local:swan}(ii)), $a(V)=2+\operatorname{sw}(V)\ge3$. The computation of the commutant in \Cref{prop:local:tame-cyclic-shape}(iii) uses only the cyclicity. $\Frob$ normalises $\rho(I_3)$, hence its unique subgroup $\langle\sigma\rangle$ of
order $3$, so $\rho(\Frob)\sigma\rho(\Frob)^{-1}=\sigma^{\pm1}$; as $2\sigma^{-1}+1=2\sigma^2+1=-(2\sigma+1)$ (from $\sigma^2+\sigma+1=0$), $\Frob$ acts on the line by
$\pm1$. In the tame case the sign was read off from $g\mapsto g^q$, which used the tame relation; that is not available on wild inertia. (ii) $2\nmid3$, so the image is tame; $\chi_{-3}(2)=-1$. (iii) A two-dimensional
faithful representation of a finite non-abelian group in characteristic $0$ that fixed a line over $\overline\Omega$ would be triangular, hence (being semisimple)
diagonal and abelian; so it is absolutely irreducible, its commutant is the scalars (Schur, \Cref{ch:reps}), and the trace-zero commutant is $0$. (iv) is
\Cref{prop:local:conductor}(iii) with \Cref{thm:local:swan}(iv).
\end{proof}

\begin{remark}[what the paper takes from this at $2$ and $3$]\label{rem:local:wild-paper}
\begin{enumerate}[label=(\alph*)]
\item For a curve of type $\mathrm{IV}$ or $\mathrm{IV}^*$ at $2$, the image of inertia on $T_\ell E$ is cyclic of order exactly $3$; this is a theorem proved in
\Cref{ch:galrep} from $E[3]^{I_2}\simeq\Phig_2[3]$ and in \Cref{ch:reduction} from the conductor exponent. With
\Cref{prop:local:wild-shapes}(ii) and $L(\Sym f,s)=L(A_f,s-1)$ (\Cref{ch:lfunctions}), $P_2(X)=1+2X$ and $e_3(E,2)=v_3(3/2)=1$, as Lemma lem:IV and
Remark rem:wild state.
\item For type $\mathrm{II}$ or $\mathrm{II}^*$ at $2$ the Kodaira type does not determine the image of inertia. When the image is non-abelian ($Q_8$ or $\SL_2(\F_3)$,
acting absolutely irreducibly), $P_2=1$ by \Cref{prop:local:wild-shapes}(iii); but \curve{336f1} has type $\mathrm{II}$ at $2$ and $P_2=1+2X$, and \curve{1232l1} has type
$\mathrm{II}^*$ at $2$ and a cubic $P_2$. The revised Remark rem:wild says this; the original version of the paper asserted $P_2=1$ for every type $\mathrm{II}$ at $2$.
\item At $q=2$ a curve of type $\mathrm{I}_n^*$ may have potentially good reduction, so the Kodaira type does not determine the Weil-Deligne shape (Definition
def:terms: $t_\ell=0$ then); and in the potentially multiplicative case the conductor exponent is $4$ or $6$ by \Cref{prop:local:wild-shapes}(iv).
\item At $q=3$ the image of inertia for types $\mathrm{IV}$, $\mathrm{IV}^*$ has order divisible by $3$ (\Cref{ch:reduction,ch:galrep}), so by (i) these types
are wild; the paper says ``At $q=3$ the types $\mathrm{IV}$, $\mathrm{IV}^*$ are wildly ramified, and the lemma does not apply.'' This is why Lemma lem:IV excludes $q=3$ and Lemma lem:euler assumes $q\ge5$ for the potentially good
rows.
\end{enumerate}
\end{remark}

\section{Exercises}\label{sec:local:exercises}

\begin{exercise}\label{ex:local:ex-expansions}
Find the first six $5$-adic digits of $1/3$ and of $-1/3$, and check, carry by carry, that their sum is $0$. Which rational number has the $7$-adic
expansion $1+2\cdot7+1\cdot7^2+2\cdot7^3+\cdots$ (digits $1,2$ repeated)? (Answers: $2,3,1,3,1,3$ and $3,1,3,1,3,1$; $-5/16$.)
\end{exercise}

\begin{exercise}\label{ex:local:ex-not-iso}
(a) Show that $\Q_p$ is not algebraically closed. (b) Show that $u\in\Q_p^\times$ lies in $\Z_p^\times$ iff $u$ is an $m$-th power for infinitely many integers $m\ge1$
(use \Cref{thm:local:powers} and \Cref{thm:local:units-structure}). (c) Deduce that every field isomorphism $\Q_p\to\Q_{p'}$ maps $\Z_p^\times$ onto $\Z_{p'}^\times$, and
that $\Q_p\not\simeq\Q_{p'}$ as fields for $p\ne p'$. (Hints: (a) $X^2-p$. (b) A unit is $[\bar u]u_1$, and $u_1$ is an $m$-th power for every $m$ prime
to $p$; an $m$-th power $x$ has $m\mid v_p(x)$. (c) The characterisation in (b) uses only the field structure. A field isomorphism is the identity on $\Q$, and
$p\notin\Z_p^\times$ while $p\in\Z_{p'}^\times$.)
\end{exercise}

\begin{exercise}\label{ex:local:ex-hensel}
(a) Find the two square roots of $-1$ in $\Z_{13}$ modulo $13^3$. (b) Show that $2$ has exactly one cube root in $\Z_5$ and find it modulo $5^3$. (c) Show that
$X^3-3$ has no root in $\Z_7$. (d) Compute the Teichm\"uller lift $[2]\in\Z_7$ modulo $7^3$ and check that it is a root of $X^2+X+1$. (Answers: $239$, $1958$; $53$;
$324$.)
\end{exercise}

\begin{exercise}\label{ex:local:ex-q2}
For $d\in\{-7,6,14,-15,17,-3,13\}$ decide whether $\Q_2(\sqrt d)$ is $\Q_2$, the unramified quadratic extension, or ramified, and in the last case give the
conductor exponent of its character. Compare with the $2$-adic valuation of the discriminant of $\Q(\sqrt d)$. (Answer: $\Q_2$ for $-7$, $-15$, $17$
($d\equiv1\pmod8$); unramified for $-3$, $13$ ($d\equiv5\pmod8$); ramified with conductor exponent $3$ for $6\in-10(\Q_2^\times)^2$ and
$14\in-2(\Q_2^\times)^2$. The discriminants $-7,24,56,-15,17,-3,13$ have $2$-adic valuations $0,3,3,0,0,0,0$: in every case the conductor exponent.)
\end{exercise}

\begin{exercise}\label{ex:local:ex-pth-powers}
Let $p$ be odd. (a) Show that $(1+p\Z_p)^p=1+p^2\Z_p$ and that $\Q_p^\times/(\Q_p^\times)^p\simeq(\Z/p)^2$. (b) Compare with the formula
$\#K^\times/(K^\times)^{\ell^k}=\ell^k\cdot\#\mu_{\ell^k}(K)$ of \Cref{thm:local:powers}(iv), valid for $\ell\ne p$; which factor is missing when $\ell=p$? (c) Show that $\log$ is not
injective on $1+2\Z_2$, and explain why this does not contradict \Cref{thm:local:units-structure}(ii). (Hints: (a) $\log$ carries $1+p\Z_p$ onto $p\Z_p$ and
$p$-th powers to $p\Z_p\cdot p=p^2\Z_p$. (b) The factor coming from the $\Z_p$ in $\Z_p^\times\simeq\Z/(p-1)\times\Z_p$; for $\ell\ne p$ the corresponding $U^{(1)}$ is
$\ell$-divisible. (c) $\log(-1)=\frac12\log((-1)^2)=0$; the theorem concerns $1+4\Z_2$ only.)
\end{exercise}

\begin{exercise}\label{ex:local:ex-cubes13}
Determine $\Q_{13}^\times/(\Q_{13}^\times)^3$ and $\Q_{13}^\times/(\Q_{13}^\times)^9$, and decide which of $2$, $3$, $5$, $8$, $65$ are cubes in $\Q_{13}$.
(Answer: $\Z/3\times\Z/3$ and $\Z/9\times\Z/3$, as $v_3(12)=1$. The cubes of $\F_{13}^\times$ are $\{1,5,8,12\}$, so $5$ and $8$ are cubes; $2$ and $3$ are not;
$65=13\cdot5$ is not, as $3\nmid v_{13}(65)$.)
\end{exercise}

\begin{exercise}\label{ex:local:ex-11a1}
Let $E=\curve{11a1}$, which has split multiplicative reduction at $11$ with $n_{11}=v_{11}(\Delta_E)=5$. Using \Cref{thm:local:kummer-inertia}(v) and
\Cref{ex:local:tate-shape-wd}, show that $I_{11}$ acts trivially on $E[5]$ but not on $E[25]$ (the paper's Lemma lem:tamagawa-exp). Using \Cref{ex:local:powers-examples}(c),
show more: $E[5]\simeq\Z/5\oplus\mu_5$ as a $G_{\Q_{11}}$-module.
\end{exercise}

\begin{exercise}\label{ex:local:ex-tame-level}
Redo \Cref{ex:local:tame-explicit} for $(q,\ell,k)=(5,3,2)$ and $(11,3,2)$: give the order of $\Gal(\Q_q(\mu_9,q^{1/9})/\Q_q)$, and compute $\Frob\,t\,\Frob^{-1}$ and
$\Frob^{-1}t\,\Frob$ as powers of $t$. Then show directly, in the affine model, that $\Phi\,t\,\Phi^{-1}=t^{q'}$ for the geometric Frobenius $\Phi=\Frob^{-1}$.
\end{exercise}

\begin{exercise}\label{ex:local:ex-14a1}
The curve $E=\curve{14a1}$, $y^2+xy+y=x^3+4x-6$, has split multiplicative reduction of type $\mathrm{I}_3$ at $7$, and PARI gives the Tate parameter
$q_E=6\cdot7^3+4\cdot7^4+6\cdot7^5+\cdots$. Compute $n$, the reduction $\bar u$ of the unit part, the class $c$ of Lemma lem:tamfactor at $\ell=3$ and $v_3(c)$, and deduce
$v_3(\Tam_7(\ad T_3E))$ from the lemma. Do the same for $\curve{11a1}$ at $q=11$, $\ell=5$. (Answers: $2$ in both cases.)
\end{exercise}

\begin{exercise}\label{ex:local:ex-cprime}
Let $q\ne\ell$ be primes, $x\in\Q_q^\times$ with $v_q(x)=n\ge1$, and $\kappa$ the Kummer cocycle of $x$ for a chosen compatible system of $\ell$-power roots,
with $\Z_\ell(1)\simeq\Z_\ell$. Working directly from the definition of $\kappa$ (without quoting \Cref{thm:local:kummer-inertia}(iv)): (a) show that
replacing the system of roots by another one changes $\kappa$ by the coboundary $g\mapsto(\chi_\ell(g)-1)z$ for some $z\in\Z_\ell$, so that
$c'=\kappa(\Frob)$ changes by $(q-1)z$; (b) show that replacing $\Frob$ by $\Frob\sigma$ ($\sigma\in I_q$) changes $c'$ by $q\,n\,t_\ell(\sigma)$; (c) deduce that
$\min(v_\ell(n),v_\ell(q-1),v_\ell(c'))$ does not depend on any choice. This is the quantity in the formula of the paper's Lemma lem:tamfactor. (d) For
$q=7$, $\ell=3$ and $x=7^3u$ with $u\equiv3\pmod7$, compute this minimum. (Answer: $0$, since $3$ is not a cube in $\F_7^\times$; contrast
\Cref{ex:local:ex-14a1}, where $\bar u=6=(-1)^3$ is a cube.)
\end{exercise}

\begin{exercise}\label{ex:local:ex-ramified-basechange}
Let $q$ be odd, $K=\Q_q(\sqrt{q^*})$ and $\ell\ne q$. (a) Show that $I_K$ has index $2$ in $I_q$ and that $t_{\ell,\Q_q}|_{I_K}=2t_{\ell,K}$. (b) For $\rho$ the Tate shape of
\Cref{ex:local:tate-shape-wd} with $x\in\Q_q^\times$, $a=v_q(x)$, show that $\rho\bmod\ell^k$ is unramified on $G_K$ iff $\ell^k\mid2a$. (c) Conclude that the criterion of
Lemma lem:tamagawa-exp is insensitive to this ramified base change for $\ell$ odd but not for $\ell=2$, in contrast with the unramified base change of its proof.
\end{exercise}

\begin{exercise}\label{ex:local:ex-sym2}
Let $\rho$ be the Tate shape of \Cref{ex:local:tate-shape-wd} with $a\ne0$ and $\varepsilon=1$. Show that $\WD(\Sym\rho)\simeq\Sp(3)$, and compute $L(\Sym V,s)$, $\dim(\Sym V)^{I_K}$
and $a(\Sym V)$. Check your answer against $\Sym V\simeq\ad V\otimes\det V$ and \Cref{prop:local:steinberg-shape}(iii).
\end{exercise}

\begin{exercise}\label{ex:local:ex-twist-conductor}
The curve $\curve{14a1}$ has multiplicative reduction at $2$ and $7$. Using \Cref{prop:local:conductor} and \Cref{prop:local:global-quadratic-local}, predict the
conductors of its quadratic twists by $\chi_{-4}$, $\chi_8$, $\chi_{-3}$ and $\chi_{-7}$. (PARI gives $112$, $448$, $126$ and $98$.)
\end{exercise}

\begin{exercise}\label{ex:local:ex-121d2}
$\curve{11a1}$ has type $\mathrm{I}_5$ at $11$, and its twist by $\chi_{-11}$ is $\curve{121d2}$, of type $\mathrm{I}_5^*$ at $11$. Using
\Cref{thm:local:shapes-table}, compute $L_{11}(A_f,1)^{-1}$ and the naive factor for both curves, and $e_3(E,11)$, $e_5(E,11)$. Explain why the true factor is
invariant under the twist and the naive one is not (Definition def:terms).
\end{exercise}

\begin{exercise}\label{ex:local:ex-e3}
For a prime $q\ne3$ at which $\rho(I_q)$ is cyclic of order $3$, $e_3(E,q)=v_3(q-\chi_{-3}(q))$ (\Cref{prop:local:tame-cyclic-shape}). Find all such $q<60$ with
$e_3(E,q)\ge2$, and the largest value of $e_3$ among them.
\end{exercise}

\begin{exercise}\label{ex:local:ex-zeta9}
Let $L=\Q_3(\zeta_9)$, $G=\Gal(L/\Q_3)\simeq(\Z/9)^\times$. (a) Using $v_L(\zeta^{a-1}-1)=3^{v_3(a-1)}$ for $a\not\equiv1\pmod9$, show that $G_0=G$, $G_1=G_2$ is the subgroup of
order $3$, and $G_3=1$. (b) Compute the Swan and Artin conductors of the characters of $G$ of orders $2$, $3$ and $6$, and compare with the conductors of the
Dirichlet characters modulo $9$. (c) Which of them is $\chi_{-3}$?
\end{exercise}

\begin{exercise}\label{ex:local:ex-order4}
Suppose $\rho(I_q)=\langle\tau\rangle$ is cyclic of order $4$ with $\det\rho|_{I_q}=1$ and $q\equiv3\pmod4$. Show that $\tr\rho(\Frob)=0$ for every Frobenius lift, that $\rho(\Frob)^2$
commutes with $\tau$, and that $V\otimes\Omega(i)$ is irreducible as a representation of $G_{\Q_q}$.
\end{exercise}

\begin{exercise}\label{ex:local:ex-ell-equals-p}
(a) Show that for $\ell\ne p$ every continuous character $G_K\to\Q_\ell^\times$ has finite image on $I_K$. (b) Show that the cyclotomic character $\chi_p$ of
$G_{\Q_p}$ has infinite image on $I_p$, and explain which step of the proof of \Cref{thm:local:grothendieck} fails for $\ell=p$.
\end{exercise}

\begin{exercise}\label{ex:local:ex-krasner}
(a) For $p$ odd, show that $\Q_p(\sqrt{p+p^3})=\Q_p(\sqrt p)$, once by Hensel's lemma and once by \Cref{lem:local:krasner}(i). (b) Show that $\Q_2(\sqrt{10})\ne\Q_2(\sqrt2)$
although $\abs{\sqrt{10}-\sqrt2}_2=\abs{\sqrt2-(-\sqrt2)}_2$: the strict inequality in Krasner's lemma cannot be relaxed.
\end{exercise}

\section{Notes and references}\label{sec:local:notes}

Absolute values, completions, $\Q_p$, Hensel's lemma and the structure of $\Q_p^\times$ follow Serre's \emph{A Course in Arithmetic} \cite[Chapter II]{SerreCourse} and
Neukirch \cite[Chapter II, Sections 1 to 5]{Neukirch}; the factorisation form of Hensel's lemma and its corollary on coefficients of irreducible polynomials
are in Neukirch, Chapter II, Section 4, and many variants are collected in Keith Conrad's notes \cite{KConradHensel}. Cassels' chapter ``Global fields'' in
Cassels and Fr\"ohlich \cite{CasselsFrohlich} is a compact alternative. Extensions of the absolute value, $e$ and $f$, Eisenstein polynomials, unramified and
tamely ramified extensions, and Krasner's lemma follow Serre's \emph{Local Fields} \cite[Chapters I to IV]{SerreLocalFields} and Neukirch, Chapter II,
Sections 6 to 10; the comparison of completions of number fields with the global theory (\Cref{thm:local:completion-numberfield},
\Cref{prop:local:decomposition}) is Neukirch, Chapter II, Sections 8 and 9.

The ramification groups, their quotients and the tame character are Serre \cite[Chapter IV, Sections 1 and 2]{SerreLocalFields}; Herbrand's theorem, upper
numbering and Hasse-Arf are Chapter IV, Section 3, and Chapter V, Section 7 there; the Artin and Swan conductors and Artin's integrality theorem are Chapter VI,
Section 2, and the representation-theoretic side is in \cite[Chapter 19]{SerreLinearReps}. The structure of $\Gal(K^{\tame}/K)$ as a profinite group with two
generators is in Neukirch, Schmidt and Wingberg \cite[Chapter VII, Section 5]{NSW}. The use of the tame quotient and of Kummer classes for $\ell$-adic
representations, and the Tate curve, are in Serre \cite[Chapters I and IV]{SerreAbelian} and Silverman \cite[Chapter V]{SilvermanATAEC}; Serre and Tate
\cite[Sections 1 and 2]{SerreTate68} prove Grothendieck's theorem in the form used here (their appendix), the criterion of N\'eron, Ogg and Shafarevich, and the
independence of $\ell$ of the image of inertia for potentially good reduction. Local class field theory is stated after Serre's chapter in
\cite{CasselsFrohlich} and Milne \cite[Chapter I]{MilneCFT}.

Weil-Deligne representations, local $L$-factors, conductors and epsilon factors follow Deligne \cite[Section 8]{Deligne73} and Tate \cite[Section 4]{Tate79};
Rohrlich \cite{Rohrlich94} is a gentle account written for elliptic curves, with the classification of two-dimensional representations. The convention that local
$L$-factors use the geometric Frobenius is the one of \cite[Section 1.7.2]{DFG04}, which the paper follows.

All chapter and section numbers of books cited in this chapter (Serre, Neukirch, Neukirch-Schmidt-Wingberg, Milne, Cassels-Fr\"ohlich, Silverman, Deligne, Tate,
Rohrlich) are recalled and were not checked against printed copies; they are pointers, and no proof in the chapter depends on them. No theorem numbers from
these sources are cited. Every numerical example and the answers to the exercises were checked with PARI/GP \cite{PARI} through \texttt{cypari2}, with
commands such as the following.
\begin{lstlisting}[language=GP]
polrootspadic(x^2+1, 5, 4)            \\ 182 and 443 modulo 5^4
lift(teichmuller(2 + O(5^4)))         \\ 182
polrootspadic(x^2+x+1, 7, 2)          \\ 30 and 18 modulo 49
E = ellinit([0,-1,1,-10,-20], O(11^8)); E.tate[3]   \\ q_E of 11a1: 10*11^5 + ...
E = ellinit([1,0,1,4,-6], O(7^6)); E.tate[3]        \\ q_E of 14a1 at 7: 6*7^3 + ...
ellglobalred(ellinit(elltwist(ellinit([1,0,1,4,-6]), -4)))[1]   \\ 112
valuation(coredisc(-2), 2)            \\ 3: conductor exponent of chi_{-8} at 2
\end{lstlisting}
